% Options for packages loaded elsewhere
% Options for packages loaded elsewhere
\PassOptionsToPackage{unicode}{hyperref}
\PassOptionsToPackage{hyphens}{url}
\PassOptionsToPackage{dvipsnames,svgnames,x11names}{xcolor}
%
\documentclass[
  czech,
  11pt,
  letterpaper,
  DIV=11,
  numbers=noendperiod]{scrreprt}
\usepackage{xcolor}
\usepackage[lmargin=1in,rmargin=1in,tmargin=1in,bmargin=1in]{geometry}
\usepackage{amsmath,amssymb}
\setcounter{secnumdepth}{5}
\usepackage{iftex}
\ifPDFTeX
  \usepackage[T1]{fontenc}
  \usepackage[utf8]{inputenc}
  \usepackage{textcomp} % provide euro and other symbols
\else % if luatex or xetex
  \usepackage{unicode-math} % this also loads fontspec
  \defaultfontfeatures{Scale=MatchLowercase}
  \defaultfontfeatures[\rmfamily]{Ligatures=TeX,Scale=1}
\fi
\usepackage{lmodern}
\ifPDFTeX\else
  % xetex/luatex font selection
\fi
% Use upquote if available, for straight quotes in verbatim environments
\IfFileExists{upquote.sty}{\usepackage{upquote}}{}
\IfFileExists{microtype.sty}{% use microtype if available
  \usepackage[]{microtype}
  \UseMicrotypeSet[protrusion]{basicmath} % disable protrusion for tt fonts
}{}
\usepackage{setspace}
\makeatletter
\@ifundefined{KOMAClassName}{% if non-KOMA class
  \IfFileExists{parskip.sty}{%
    \usepackage{parskip}
  }{% else
    \setlength{\parindent}{0pt}
    \setlength{\parskip}{6pt plus 2pt minus 1pt}}
}{% if KOMA class
  \KOMAoptions{parskip=half}}
\makeatother
% Make \paragraph and \subparagraph free-standing
\makeatletter
\ifx\paragraph\undefined\else
  \let\oldparagraph\paragraph
  \renewcommand{\paragraph}{
    \@ifstar
      \xxxParagraphStar
      \xxxParagraphNoStar
  }
  \newcommand{\xxxParagraphStar}[1]{\oldparagraph*{#1}\mbox{}}
  \newcommand{\xxxParagraphNoStar}[1]{\oldparagraph{#1}\mbox{}}
\fi
\ifx\subparagraph\undefined\else
  \let\oldsubparagraph\subparagraph
  \renewcommand{\subparagraph}{
    \@ifstar
      \xxxSubParagraphStar
      \xxxSubParagraphNoStar
  }
  \newcommand{\xxxSubParagraphStar}[1]{\oldsubparagraph*{#1}\mbox{}}
  \newcommand{\xxxSubParagraphNoStar}[1]{\oldsubparagraph{#1}\mbox{}}
\fi
\makeatother

\usepackage{color}
\usepackage{fancyvrb}
\newcommand{\VerbBar}{|}
\newcommand{\VERB}{\Verb[commandchars=\\\{\}]}
\DefineVerbatimEnvironment{Highlighting}{Verbatim}{commandchars=\\\{\}}
% Add ',fontsize=\small' for more characters per line
\usepackage{framed}
\definecolor{shadecolor}{RGB}{241,243,245}
\newenvironment{Shaded}{\begin{snugshade}}{\end{snugshade}}
\newcommand{\AlertTok}[1]{\textcolor[rgb]{0.68,0.00,0.00}{#1}}
\newcommand{\AnnotationTok}[1]{\textcolor[rgb]{0.37,0.37,0.37}{#1}}
\newcommand{\AttributeTok}[1]{\textcolor[rgb]{0.40,0.45,0.13}{#1}}
\newcommand{\BaseNTok}[1]{\textcolor[rgb]{0.68,0.00,0.00}{#1}}
\newcommand{\BuiltInTok}[1]{\textcolor[rgb]{0.00,0.23,0.31}{#1}}
\newcommand{\CharTok}[1]{\textcolor[rgb]{0.13,0.47,0.30}{#1}}
\newcommand{\CommentTok}[1]{\textcolor[rgb]{0.37,0.37,0.37}{#1}}
\newcommand{\CommentVarTok}[1]{\textcolor[rgb]{0.37,0.37,0.37}{\textit{#1}}}
\newcommand{\ConstantTok}[1]{\textcolor[rgb]{0.56,0.35,0.01}{#1}}
\newcommand{\ControlFlowTok}[1]{\textcolor[rgb]{0.00,0.23,0.31}{\textbf{#1}}}
\newcommand{\DataTypeTok}[1]{\textcolor[rgb]{0.68,0.00,0.00}{#1}}
\newcommand{\DecValTok}[1]{\textcolor[rgb]{0.68,0.00,0.00}{#1}}
\newcommand{\DocumentationTok}[1]{\textcolor[rgb]{0.37,0.37,0.37}{\textit{#1}}}
\newcommand{\ErrorTok}[1]{\textcolor[rgb]{0.68,0.00,0.00}{#1}}
\newcommand{\ExtensionTok}[1]{\textcolor[rgb]{0.00,0.23,0.31}{#1}}
\newcommand{\FloatTok}[1]{\textcolor[rgb]{0.68,0.00,0.00}{#1}}
\newcommand{\FunctionTok}[1]{\textcolor[rgb]{0.28,0.35,0.67}{#1}}
\newcommand{\ImportTok}[1]{\textcolor[rgb]{0.00,0.46,0.62}{#1}}
\newcommand{\InformationTok}[1]{\textcolor[rgb]{0.37,0.37,0.37}{#1}}
\newcommand{\KeywordTok}[1]{\textcolor[rgb]{0.00,0.23,0.31}{\textbf{#1}}}
\newcommand{\NormalTok}[1]{\textcolor[rgb]{0.00,0.23,0.31}{#1}}
\newcommand{\OperatorTok}[1]{\textcolor[rgb]{0.37,0.37,0.37}{#1}}
\newcommand{\OtherTok}[1]{\textcolor[rgb]{0.00,0.23,0.31}{#1}}
\newcommand{\PreprocessorTok}[1]{\textcolor[rgb]{0.68,0.00,0.00}{#1}}
\newcommand{\RegionMarkerTok}[1]{\textcolor[rgb]{0.00,0.23,0.31}{#1}}
\newcommand{\SpecialCharTok}[1]{\textcolor[rgb]{0.37,0.37,0.37}{#1}}
\newcommand{\SpecialStringTok}[1]{\textcolor[rgb]{0.13,0.47,0.30}{#1}}
\newcommand{\StringTok}[1]{\textcolor[rgb]{0.13,0.47,0.30}{#1}}
\newcommand{\VariableTok}[1]{\textcolor[rgb]{0.07,0.07,0.07}{#1}}
\newcommand{\VerbatimStringTok}[1]{\textcolor[rgb]{0.13,0.47,0.30}{#1}}
\newcommand{\WarningTok}[1]{\textcolor[rgb]{0.37,0.37,0.37}{\textit{#1}}}

\usepackage{longtable,booktabs,array}
\newcounter{none} % for unnumbered tables
\usepackage{calc} % for calculating minipage widths
% Correct order of tables after \paragraph or \subparagraph
\usepackage{etoolbox}
\makeatletter
\patchcmd\longtable{\par}{\if@noskipsec\mbox{}\fi\par}{}{}
\makeatother
% Allow footnotes in longtable head/foot
\IfFileExists{footnotehyper.sty}{\usepackage{footnotehyper}}{\usepackage{footnote}}
\makesavenoteenv{longtable}
\usepackage{graphicx}
\makeatletter
\newsavebox\pandoc@box
\newcommand*\pandocbounded[1]{% scales image to fit in text height/width
  \sbox\pandoc@box{#1}%
  \Gscale@div\@tempa{\textheight}{\dimexpr\ht\pandoc@box+\dp\pandoc@box\relax}%
  \Gscale@div\@tempb{\linewidth}{\wd\pandoc@box}%
  \ifdim\@tempb\p@<\@tempa\p@\let\@tempa\@tempb\fi% select the smaller of both
  \ifdim\@tempa\p@<\p@\scalebox{\@tempa}{\usebox\pandoc@box}%
  \else\usebox{\pandoc@box}%
  \fi%
}
% Set default figure placement to htbp
\def\fps@figure{htbp}
\makeatother


% definitions for citeproc citations
\NewDocumentCommand\citeproctext{}{}
\NewDocumentCommand\citeproc{mm}{%
  \begingroup\def\citeproctext{#2}\cite{#1}\endgroup}
\makeatletter
 % allow citations to break across lines
 \let\@cite@ofmt\@firstofone
 % avoid brackets around text for \cite:
 \def\@biblabel#1{}
 \def\@cite#1#2{{#1\if@tempswa , #2\fi}}
\makeatother
\newlength{\cslhangindent}
\setlength{\cslhangindent}{1.5em}
\newlength{\csllabelwidth}
\setlength{\csllabelwidth}{3em}
\newenvironment{CSLReferences}[2] % #1 hanging-indent, #2 entry-spacing
 {\begin{list}{}{%
  \setlength{\itemindent}{0pt}
  \setlength{\leftmargin}{0pt}
  \setlength{\parsep}{0pt}
  % turn on hanging indent if param 1 is 1
  \ifodd #1
   \setlength{\leftmargin}{\cslhangindent}
   \setlength{\itemindent}{-1\cslhangindent}
  \fi
  % set entry spacing
  \setlength{\itemsep}{#2\baselineskip}}}
 {\end{list}}
\usepackage{calc}
\newcommand{\CSLBlock}[1]{\hfill\break\parbox[t]{\linewidth}{\strut\ignorespaces#1\strut}}
\newcommand{\CSLLeftMargin}[1]{\parbox[t]{\csllabelwidth}{\strut#1\strut}}
\newcommand{\CSLRightInline}[1]{\parbox[t]{\linewidth - \csllabelwidth}{\strut#1\strut}}
\newcommand{\CSLIndent}[1]{\hspace{\cslhangindent}#1}

\ifLuaTeX
\usepackage[bidi=basic,shorthands=off]{babel}
\else
\usepackage[bidi=default,shorthands=off]{babel}
\fi
\ifLuaTeX
  \usepackage{selnolig} % disable illegal ligatures
\fi


\setlength{\emergencystretch}{3em} % prevent overfull lines

\providecommand{\tightlist}{%
  \setlength{\itemsep}{0pt}\setlength{\parskip}{0pt}}



 


\usepackage{amsmath}
\usepackage{amsfonts}
\usepackage{amssymb}    
\KOMAoption{captions}{tableheading}
\makeatletter
\@ifpackageloaded{tcolorbox}{}{\usepackage[skins,breakable]{tcolorbox}}
\@ifpackageloaded{fontawesome5}{}{\usepackage{fontawesome5}}
\definecolor{quarto-callout-color}{HTML}{909090}
\definecolor{quarto-callout-note-color}{HTML}{0758E5}
\definecolor{quarto-callout-important-color}{HTML}{CC1914}
\definecolor{quarto-callout-warning-color}{HTML}{EB9113}
\definecolor{quarto-callout-tip-color}{HTML}{00A047}
\definecolor{quarto-callout-caution-color}{HTML}{FC5300}
\definecolor{quarto-callout-color-frame}{HTML}{acacac}
\definecolor{quarto-callout-note-color-frame}{HTML}{4582ec}
\definecolor{quarto-callout-important-color-frame}{HTML}{d9534f}
\definecolor{quarto-callout-warning-color-frame}{HTML}{f0ad4e}
\definecolor{quarto-callout-tip-color-frame}{HTML}{02b875}
\definecolor{quarto-callout-caution-color-frame}{HTML}{fd7e14}
\makeatother
\makeatletter
\@ifpackageloaded{bookmark}{}{\usepackage{bookmark}}
\makeatother
\makeatletter
\@ifpackageloaded{caption}{}{\usepackage{caption}}
\AtBeginDocument{%
\ifdefined\contentsname
  \renewcommand*\contentsname{Obsah}
\else
  \newcommand\contentsname{Obsah}
\fi
\ifdefined\listfigurename
  \renewcommand*\listfigurename{Seznam obrázků}
\else
  \newcommand\listfigurename{Seznam obrázků}
\fi
\ifdefined\listtablename
  \renewcommand*\listtablename{Seznam tabulek}
\else
  \newcommand\listtablename{Seznam tabulek}
\fi
\ifdefined\figurename
  \renewcommand*\figurename{Obrázek}
\else
  \newcommand\figurename{Obrázek}
\fi
\ifdefined\tablename
  \renewcommand*\tablename{Tabulka}
\else
  \newcommand\tablename{Tabulka}
\fi
}
\@ifpackageloaded{float}{}{\usepackage{float}}
\floatstyle{ruled}
\@ifundefined{c@chapter}{\newfloat{codelisting}{h}{lop}}{\newfloat{codelisting}{h}{lop}[chapter]}
\floatname{codelisting}{Výpis}
\newcommand*\listoflistings{\listof{codelisting}{Seznam výpisů}}
\makeatother
\makeatletter
\makeatother
\makeatletter
\@ifpackageloaded{caption}{}{\usepackage{caption}}
\@ifpackageloaded{subcaption}{}{\usepackage{subcaption}}
\makeatother
\usepackage{bookmark}
\IfFileExists{xurl.sty}{\usepackage{xurl}}{} % add URL line breaks if available
\urlstyle{same}
\hypersetup{
  pdftitle={Havi Notes   Poznámky z pravděpodobnosti a   statistiky},
  pdfauthor={Jirka Haviger},
  pdflang={cs},
  colorlinks=true,
  linkcolor={blue},
  filecolor={Maroon},
  citecolor={Blue},
  urlcolor={Blue},
  pdfcreator={LaTeX via pandoc}}


\title{\textbf{Havi Notes} Poznámky z pravděpodobnosti a statistiky}
\author{Jirka Haviger}
\date{2025-12-15}
\begin{document}
\maketitle

\renewcommand*\contentsname{Obsah}
{
\setcounter{tocdepth}{1}
\tableofcontents
}

\setstretch{1.2}
\bookmarksetup{startatroot}

\chapter*{Úvodní poznámky}\label{uxfavodnuxed-poznuxe1mky}
\addcontentsline{toc}{chapter}{Úvodní poznámky}

\markboth{Úvodní poznámky}{Úvodní poznámky}

Na této stránce se budou objevovat zdroje a materiály vhodné k pochopení
základních i pokročilejších myšlenek z různých obastí. Stránky tedy
nejsou rigorózní učebnicí, jsou zaměřené na \textbf{vizualizaci
základních myšlenek a ideí}, od které jsou přesné definice, věty a
algorimty odvozené.

\begin{itemize}
\tightlist
\item
  \textbf{Pravděpodobnost}
\item
  \textbf{Statistika}
\end{itemize}

\begin{tcolorbox}[enhanced jigsaw, arc=.35mm, bottomrule=.15mm, bottomtitle=1mm, breakable, colback=white, colbacktitle=quarto-callout-tip-color!10!white, colframe=quarto-callout-tip-color-frame, coltitle=black, left=2mm, leftrule=.75mm, opacityback=0, opacitybacktitle=0.6, rightrule=.15mm, title=\textcolor{quarto-callout-tip-color}{\faLightbulb}\hspace{0.5em}{Tip}, titlerule=0mm, toprule=.15mm, toptitle=1mm]

styl používaný pro tipy

\end{tcolorbox}

\begin{tcolorbox}[enhanced jigsaw, arc=.35mm, bottomrule=.15mm, breakable, colback=white, colframe=quarto-callout-tip-color-frame, left=2mm, leftrule=.75mm, opacityback=0, rightrule=.15mm, toprule=.15mm]
\begin{minipage}[t]{5.5mm}
\textcolor{quarto-callout-tip-color}{\faLightbulb}
\end{minipage}%
\begin{minipage}[t]{\textwidth - 5.5mm}

styl používaný pro tipy, verze ``simple''

\end{minipage}%
\end{tcolorbox}

\begin{tcolorbox}[enhanced jigsaw, arc=.35mm, bottomrule=.15mm, breakable, colback=white, colframe=quarto-callout-tip-color-frame, left=2mm, leftrule=.75mm, opacityback=0, rightrule=.15mm, toprule=.15mm]

styl používaný pro tipy, verze ``minimal''

\end{tcolorbox}

\begin{tcolorbox}[enhanced jigsaw, arc=.35mm, bottomrule=.15mm, bottomtitle=1mm, breakable, colback=white, colbacktitle=quarto-callout-note-color!10!white, colframe=quarto-callout-note-color-frame, coltitle=black, left=2mm, leftrule=.75mm, opacityback=0, opacitybacktitle=0.6, rightrule=.15mm, title=\textcolor{quarto-callout-note-color}{\faInfo}\hspace{0.5em}{Poznámka}, titlerule=0mm, toprule=.15mm, toptitle=1mm]

styl používaný pro poznámky

\end{tcolorbox}

\begin{tcolorbox}[enhanced jigsaw, arc=.35mm, bottomrule=.15mm, bottomtitle=1mm, breakable, colback=white, colbacktitle=quarto-callout-warning-color!10!white, colframe=quarto-callout-warning-color-frame, coltitle=black, left=2mm, leftrule=.75mm, opacityback=0, opacitybacktitle=0.6, rightrule=.15mm, title=\textcolor{quarto-callout-warning-color}{\faExclamationTriangle}\hspace{0.5em}{Varování}, titlerule=0mm, toprule=.15mm, toptitle=1mm]

styl používaný pro varování

\end{tcolorbox}

\begin{tcolorbox}[enhanced jigsaw, arc=.35mm, bottomrule=.15mm, bottomtitle=1mm, breakable, colback=white, colbacktitle=quarto-callout-important-color!10!white, colframe=quarto-callout-important-color-frame, coltitle=black, left=2mm, leftrule=.75mm, opacityback=0, opacitybacktitle=0.6, rightrule=.15mm, title=\textcolor{quarto-callout-important-color}{\faExclamation}\hspace{0.5em}{Důležité}, titlerule=0mm, toprule=.15mm, toptitle=1mm]

Styl pro důležité informace

\end{tcolorbox}

\begin{tcolorbox}[enhanced jigsaw, arc=.35mm, bottomrule=.15mm, bottomtitle=1mm, breakable, colback=white, colbacktitle=quarto-callout-caution-color!10!white, colframe=quarto-callout-caution-color-frame, coltitle=black, left=2mm, leftrule=.75mm, opacityback=0, opacitybacktitle=0.6, rightrule=.15mm, title=\textcolor{quarto-callout-caution-color}{\faFire}\hspace{0.5em}{Upozornění}, titlerule=0mm, toprule=.15mm, toptitle=1mm]

Styl pro důležité informace

\end{tcolorbox}

\begin{tcolorbox}[enhanced jigsaw, arc=.35mm, bottomrule=.15mm, bottomtitle=1mm, breakable, colback=white, colbacktitle=quarto-callout-tip-color!10!white, colframe=quarto-callout-tip-color-frame, coltitle=black, left=2mm, leftrule=.75mm, opacityback=0, opacitybacktitle=0.6, rightrule=.15mm, title=\textcolor{quarto-callout-tip-color}{\faLightbulb}\hspace{0.5em}{Řešení}, titlerule=0mm, toprule=.15mm, toptitle=1mm]

Výsledná pravděpodobnost je \(P(A) = 0.5\), protože\ldots{}

\end{tcolorbox}

\begin{tcolorbox}[enhanced jigsaw, arc=.35mm, bottomrule=.15mm, breakable, colback=white, colframe=quarto-callout-tip-color-frame, left=2mm, leftrule=.75mm, opacityback=0, rightrule=.15mm, toprule=.15mm]

\vspace{-3mm}\textbf{Řešení v minimal verzi}\vspace{3mm}

Výsledná pravděpodobnost je \(P(A) = 0.5\), protože\ldots{}

\end{tcolorbox}

Autor čerpá inspiraci z různých zdrojů, z některých tu dobrou a z
nkterých i tu, kterou zde používat nebudeme. Pro pochopení základních
pojmů a ideí doporučuji také následující zdroje:

Na stránce \textbf{Betterexpleined} ukazuje (Azad, 2024) vizuálně pěkně
jednotlivé matematické koncepty od běžných operací s čísly přes maticový
počet po komplexní čísla a používá velmi srozumitelnou angličtinu.

Na stránce \textbf{stats.libretext} je volně dostupné větší množství
materiálů o statistice a pravděpodobnosti. Doporučuji zejména (Grinstead
\& Snell, 2012), který má velmi pěkně vysvětlenou teorii
pravděpodobnosti, zvláště markovovy řetězce.

další přibudou později \ldots{}

J.Haviger, 2025

\emph{Tato stránka je založená na quarto book stylu s využitím plotly.}

\bookmarksetup{startatroot}

\chapter{}\label{section}

A druhý pokus

a ctřetí

\part{\textbf{Deskriptivní statistika}}

\chapter{Deskriptivní statistiky}\label{deskriptivnuxed-statistiky}

Představete si, že máte seznam studentů a u každého z nich výsledek test
ze statistiky (údaje můžeme mít uložené například jako slopupec hodnot v
excelu). \textbf{Deskriptivní statistiky} jsou čísla, která souhrně
charakterizují \textbf{velikost} hodnot a jejich \textbf{rozložení}
nějaké \textbf{proměnné}, v našem příkladu výsledků testu. Nejprve si
pojďme ujasnit co se rozumí pod pojmem proměnná a jaké typy proměnných
máme.

\section{Proměnná}\label{promux11bnnuxe1}

Proměnná je jakákoli \textbf{měřená} (zaznamenávaná) \textbf{veličina},
která \textbf{může nabývat různých hodnot}. Například čas průjezdu kola
v motocyklové soutěži, počet nehod v kraji za měsíc, výdaje domácností
na stravu, spokojenost s výukou, barva očí, pohlaví pacienta, \ldots{}
Hodnoty kterých tyto veličiny mohou nabývat jsou různé a podle toho s
nimi pak můžeme pracovat. Proto se proměnné rozdělují na následující
typu, kterým se také říká \textbf{úroveň měření}

\textbf{Nominální proměnná}: objekty jen pojmenováváme. Např.
\(\Omega = \{\text{muž, žena}\}\) nebo
\(\Omega = \{Královéhradecký, Pardubický, ...\}\) jako kraj rodiště
studentů.Matematicky můžeme kategorie popsat pouze četností (např kolik
je žen) a relativní četností (jaké je procentuální zastoupení žen).
Patří mezi \textbf{kvalitativní} proměnné, speciálním příadem je
\textbf{dichotomická}, která má jen dva možné výsledky.

\textbf{Ordinální proměnná}: Kategorie už mají jasné pořadí. Např.
známky ve škole, pocity spokojenosti (nízká, střední, vysoká) nebo
pořadí v závodu. Zde už má smysl mluvit o tom, že hodnoty rostou a na
ose x je uspořádat od nejmenší po největší. Patří také mezi
\textbf{kvalitativní} proměnné.

\textbf{Kardinální diskrétní proměnná}: Na ose \(x\) jsou skutečná čísla
(počet dětí, počet padlých šestek). Rozestupy mezi nimi jsou matematicky
definované. Stále jsou to ale izolované body -- mezi 1 a 2 dětmi
neexistuje žádná hodnota.

\textbf{Kardinální spojitá proměnná}: Body na ose \(x\) se k sobě
přiblíží natolik, že mezi jakýmikoliv dvěma hodnotami vždy najdeme
nekonečně mnoho dalších. Tímto se budeme zabývat v jiné kapitole.

Obě kardinální porměnné patří mezi \textbf{kvantitativní}, alternativně
se pro ně používají i označení \textbf{metrická} nebo \textbf{škála}.
Někdy se dělí na \textbf{intervalové}, které nemají pevnou nulu (teplota
ve stuních celsia, datum) a \textbf{poměrové} které mají nulu
definovanou (věk, plat, hmotnost, rychlost, \ldots).

\section{Deskriptivní statistiky
obecně}\label{deskriptivnuxed-statistiky-obecnux11b}

Když máme nějakou proměnnou, můžeme provádět její měření (nebo
zaznamenávat její hodnoty), kterých však může mýt mnoho. Proto je vhodné
těchto ``mnoho'' hodnot nějak popsat a k tomu slouží popisná neboli
\textbf{deskriptivní statitika}. Jaké nástroje nám tedy nabízí? V zásadě
dvoje:

\begin{itemize}
\tightlist
\item
  \textbf{míry polohy} určují \textbf{kde} na ose jsou data umístěna.
\item
  \textbf{míry variability} ukazují \textbf{jak} jsou data nahuštěna
  nebo naopak rozprostřena.
\end{itemize}

V grafu jsou původní data, data posunutá (+2) a data rozptýlená (.3). U
posunu se jen změnila poloha těch dat ale vzájemné rozestupy zůstaly
stejné. U vynásobení se změnily rozestupy i celkový rozsah.

\begin{Shaded}
\begin{Highlighting}[]
\NormalTok{//| echo: false}
\NormalTok{// 1. Příprava dat}
\NormalTok{dataZdroj = [0, 1, 1.8, 2.3, 3, 3.5, 4, 5, 8, 10]}

\NormalTok{transformData = (data, groupName, transformFn = d =\textgreater{} d) =\textgreater{} \{}
\NormalTok{  const transformed = data.map(d =\textgreater{} transformFn(d));}
\NormalTok{  const minVal = d3.min(transformed);}
\NormalTok{  const maxVal = d3.max(transformed);}
\NormalTok{  return transformed.map(d =\textgreater{} (\{ }
\NormalTok{    hodnota: d, }
\NormalTok{    skupina: groupName, }
\NormalTok{    isMinMax: (d === minVal || d === maxVal) // Nová proměnná pro rozlišení min/max}
\NormalTok{  \}));}
\NormalTok{\};}

\NormalTok{minMaxLines = dataProGraf.reduce((acc, curr) =\textgreater{} \{}
\NormalTok{  if (!acc[curr.skupina]) \{}
\NormalTok{    acc[curr.skupina] = \{ min: Infinity, max: {-}Infinity \};}
\NormalTok{  \}}
\NormalTok{  acc[curr.skupina].min = Math.min(acc[curr.skupina].min, curr.hodnota);}
\NormalTok{  acc[curr.skupina].max = Math.max(acc[curr.skupina].max, curr.hodnota);}
\NormalTok{  return acc;}
\NormalTok{\}, \{\});}

\NormalTok{minMaxLinesArray = Object.keys(minMaxLines).map(key =\textgreater{} (\{}
\NormalTok{  skupina: key,}
\NormalTok{  x1: minMaxLines[key].min,}
\NormalTok{  x2: minMaxLines[key].max}
\NormalTok{\}));}
\NormalTok{dataProGraf = [}
\NormalTok{  ...transformData(dataZdroj, "Data"),}
\NormalTok{  ...transformData(dataZdroj, "Změna polohy (+5)", d =\textgreater{} d + 5),}
\NormalTok{  ...transformData(dataZdroj, "Změna variability (.3)", d =\textgreater{} d * 3)}
\NormalTok{];}

\NormalTok{Plot.plot(\{}
\NormalTok{  width: 800,}
\NormalTok{  height: 250, // Trochu vyšší pro lepší přehled}
\NormalTok{  grid: true,}
\NormalTok{  marginLeft: 150, }
\NormalTok{  x: \{ domain: [0, 34], label: "Hodnota →" \},}
\NormalTok{  y: \{ }
\NormalTok{    domain: ["Data", "Změna polohy (+5)", "Změna variability (.3)"],}
\NormalTok{    label: null,}
\NormalTok{    inset: 20 }
\NormalTok{  \},  }
\NormalTok{  marks: [}
\NormalTok{    // Vodorovné linky pro každou skupinu}
\NormalTok{    Plot.ruleY(["Data", "Změna polohy", "Změna variability"], \{stroke: "lightgrey"\}),}
\NormalTok{    // Úsečky mezi min a max hodnotou}
\NormalTok{Plot.link(minMaxLinesArray, \{}
\NormalTok{      x1: "x1", }
\NormalTok{      x2: "x2", }
\NormalTok{      y1: "skupina", }
\NormalTok{      y2: "skupina", }
\NormalTok{      stroke: "steelblue", }
\NormalTok{      strokeWidth: 4,}
\NormalTok{      strokeOpacity: 0.6 // trošku průhlednosti, aby nezanikly tečky}
\NormalTok{    \}),}
\NormalTok{    // Samotné body {-} nyní s podmíněnou barvou}
\NormalTok{    Plot.dot(dataProGraf, \{}
\NormalTok{      x: "hodnota",}
\NormalTok{      y: "skupina",}
\NormalTok{      fill: d =\textgreater{} d.isMinMax ? "steelblue" : "gray", // Podmíněná barva}
\NormalTok{      stroke: "white",}
\NormalTok{      r: 6}
\NormalTok{    \})}
\NormalTok{  ]}
\NormalTok{\})}
\end{Highlighting}
\end{Shaded}

\subsection*{Příklad}\label{pux159uxedklad}

Abychom nemluvli jen teoreticky, pojďme si ukázat několik základních
deskriptivních statistik na konkrétním a malém příkladu. Máme deset
domácností a sledujeme, jaký mají počet motorových vozidel (automobilů,
motocyklů, elektrokoloběžek). Získaná data po uspořádání od nejmenšího
po největší jsou \(\{1, 2, 2, 2, 3, 3, 4, 5, 8, 11\}\). Data jsou
\textbf{kardinální poměrová}, nula je korektně definovaná.

První otázkou je, kolik aut má ``normální'' domácnost, kolik je minimum,
maximum, případně jaký počet aut má dolních 25\% domácností. Jinými
slovy kde v našich datech leží střed, kde leží minimum, kde čtvrtina.
Protože se ptáme otázkou \textbf{kde}, na jakém místě osy, nazývají se
takovéto hodnoty \textbf{míry polohy}.

Nejznámější mírou polohy je aritmetický \textbf{průměr \(\bar{x}\)}.
Spočítáme ho tak, že sečteme všechny hodnoty a vydělíme je jejich
počtem. Protože však souvisí s dalšími momentovými charakteristikami,
budeme se mu věnovat v jiné kapitole. \textbf{Průměr (aritmetický
průměr)}: \[
   m = \bar{x}= \frac{1 + 2 + 2 + 2 + 3 + 3 + 4 + 5 + 8 + 11}{10} = \frac{41}{10} = 4.1
   \]

Můžeme tedy říci, že průměrná domácnost má \(4.1\) motorových vozidel

\section{Míra založená na
četnostech}\label{muxedra-zaloux17eenuxe1-na-ux10detnostech}

Mírou polohy založenou na četnostech je \textbf{Modus (mod)}, což je
hodnota, která se v datovém souboru vyskytuje nejčastěji. V našem
případě je to hodnota 2, která se vyskytuje třikrát.

\(\{1, \mathbf{ 2, 2, 2}, 3, 3, 4, 5, 8, 11\}\), tedy \(mds=2\).

Můžeme tedy říci, že nejčastěji má domácnost dvě motorová vozidla.

\begin{tcolorbox}[enhanced jigsaw, arc=.35mm, bottomrule=.15mm, bottomtitle=1mm, breakable, colback=white, colbacktitle=quarto-callout-note-color!10!white, colframe=quarto-callout-note-color-frame, coltitle=black, left=2mm, leftrule=.75mm, opacityback=0, opacitybacktitle=0.6, rightrule=.15mm, title=\textcolor{quarto-callout-note-color}{\faInfo}\hspace{0.5em}{Poznámka}, titlerule=0mm, toprule=.15mm, toptitle=1mm]

Modus ůžeme určit pro věechny typy proměnných, nominální, ordinální i
kardinální.

\end{tcolorbox}

\section{Sloupcový graf četností a
histogram}\label{sloupcovuxfd-graf-ux10detnostuxed-a-histogram}

Předchozí data můžeme velmi snadno vizualizovat pomocí sloupcovho grafu

\begin{Shaded}
\begin{Highlighting}[]
\NormalTok{  //| screenshot: true}
\NormalTok{  //| echo: false}
\NormalTok{  //| }
\NormalTok{  // 1. Definice dat}
\NormalTok{  data01 = [1, 2, 2, 2, 3, 3, 4, 5, 8, 11]}

\NormalTok{  // 2. Vykreslení grafu pomocí Plot}
\NormalTok{  Plot.plot(\{}
\NormalTok{    marks: [}
\NormalTok{      // rectY automaticky spočítá četnosti (binning není nutný pro diskrétní čísla)}
\NormalTok{      Plot.rectY(data01, Plot.binX(\{y: "count"\}, \{thresholds: d3.range(0.5, 12.5, 1), fill: "steelblue"\})),}
\NormalTok{      // Přidání pravítka na nule pro lepší čitelnost}
\NormalTok{      Plot.ruleY([0])}
\NormalTok{    ],}
\NormalTok{    x: \{}
\NormalTok{      label: "Hodnota",}
\NormalTok{      ticks: [1, 2, 3, 4, 5, 6, 7, 8, 9, 10, 11]}
\NormalTok{    \},}
\NormalTok{    y: \{}
\NormalTok{      label: "Četnost",}
\NormalTok{      grid: true,}
\NormalTok{      ticks: [0, 1, 2, 3, 4, 5, 6, 7, 8, 9, 10]}
\NormalTok{    \},}
\NormalTok{    style: \{}
\NormalTok{      fontSize: "14px",}
\NormalTok{      background: "transparent"}
\NormalTok{    \}}
\NormalTok{  \})}
\end{Highlighting}
\end{Shaded}

Podobný graf můžeme vytvořit i pro proměnné, které jsou nominální.
Ukázka je z oboru kvantitativní lingvistiky: v následujících textech
necháme spočítat četnosti jednotlivých písmen a vložíme je do grafu
uspořádané podle velikosti od největší rel.četnosti

\begin{tcolorbox}[enhanced jigsaw, arc=.35mm, bottomrule=.15mm, bottomtitle=1mm, breakable, colback=white, colbacktitle=quarto-callout-note-color!10!white, colframe=quarto-callout-note-color-frame, coltitle=black, left=2mm, leftrule=.75mm, opacityback=0, opacitybacktitle=0.6, rightrule=.15mm, title=\textcolor{quarto-callout-note-color}{\faInfo}\hspace{0.5em}{Poznámka}, titlerule=0mm, toprule=.15mm, toptitle=1mm]

\# použité texty \textbf{Švejk (Hašek)} \emph{Tak nám zabili Ferdinanda,
řekla posluhovačka panu Švejkovi, který opustiv před léty vojenskou
službu, když byl definitivně prohlášen vojenskou lékařskou komisí za
blba, živil se prodejem psů, ošklivých nečistokrevných oblud, kterým
falšoval rodokmeny. Kromě tohoto zaměstnání byl stižen revmatismem a
mazal si právě kolena opodeldokem. Kterýho Ferdinanda, paní Müllerová?}

\textbf{R.U.R. (Čapek)} \emph{DOMIN: kromě toho dodáme za stejnou cenu
patnáct tisíc kusů Robotů pro vykládání lodí. Značka: Roboti pro těžkou
práci. Nakládací výkon jednoho Robota: dvě stě šedesát kilogramů za
hodinu. Roboti pracují ve dne v noci. Vyžadují jen nepatrné náklady na
stravu. Sulla, pošlete to!}

\textbf{Škola hrou (Komenský)} \emph{Vznešení a urození, ctihodní a
slovutní pánové, i všechno vzácné diváctvo! Mluvit ve škole o škole a
pro školu je věc zvlášť případná. Poněvadž jsme se tedy tak slavnostně
shromáždili ve škole, mluvme o škole a čiňme pro školu, cokoli můžeme.
Ale co je to, co můžeme? Za prvé vylévati prosby\ldots{}}

\end{tcolorbox}

\begin{Shaded}
\begin{Highlighting}[]
\NormalTok{// 1. Definice dat (bez složených závorek kolem celého bloku)}
\NormalTok{texts = (\{}
\NormalTok{  "Švejk (Hašek)": "„Tak nám zabili Ferdinanda,“ řekla posluhovačka panu Švejkovi, který opustiv před léty vojenskou službu, když byl definitivně prohlášen vojenskou lékařskou komisí za blba, živil se prodejem psů, ošklivých nečistokrevných oblud, kterým falšoval rodokmeny. Kromě tohoto zaměstnání byl stižen revmatismem a mazal si právě kolena opodeldokem. „Kterýho Ferdinanda, paní Müllerová?“",}
  
\NormalTok{  "R.U.R. (Čapek)": "DOMIN: „kromě toho dodáme za stejnou cenu patnáct tisíc kusů Robotů pro vykládání lodí. Značka: Roboti pro těžkou práci. Nakládací výkon jednoho Robota: dvě stě šedesát kilogramů za hodinu. Roboti pracují ve dne v noci. Vyžadují jen nepatrné náklady na stravu.“ Sulla, pošlete to!",}
  
\NormalTok{  "Škola hrou (Komenský)": "„Vznešení a urození, ctihodní a slovutní pánové, i všechno vzácné diváctvo! Mluvit ve škole o škole a pro školu je věc zvlášť případná. Poněvadž jsme se tedy tak slavnostně shromáždili ve škole, mluvme o škole a čiňme pro školu, cokoli můžeme. Ale co je to, co můžeme? Za prvé vylévati prosby...“"}
\NormalTok{\})}

\NormalTok{// 2. Funkce pro výpočet ABSOLUTNÍ četnosti}
\NormalTok{function getLetterCounts(text) \{}
\NormalTok{  const cleaned = text.toLowerCase().normalize("NFD").replace(/[\textbackslash{}u0300{-}\textbackslash{}u036f]/g, "").replace(/[\^{}a{-}z]/g, "");}
\NormalTok{  const counts = new Map();}
\NormalTok{  for (const char of cleaned) counts.set(char, (counts.get(char) || 0) + 1);}
  
\NormalTok{  return Array.from(counts, ([letter, count]) =\textgreater{} (\{ }
\NormalTok{    letter: letter.toUpperCase(), }
\NormalTok{    count: count }
\NormalTok{  \})).sort((a, b) =\textgreater{} b.count {-} a.count);}
\NormalTok{\}}

\NormalTok{// 3. Samostatné viewof pro výběr}
\NormalTok{viewof selection = Inputs.select(Object.keys(texts), \{label: "Vyberte text"\})}

\NormalTok{// 4. Samotný graf}
\NormalTok{Plot.plot(\{}
\NormalTok{  title: \textasciigrave{}Četnost písmen: $\{selection\}\textasciigrave{},}
\NormalTok{  x: \{ }
\NormalTok{    label: "Písmeno", }
\NormalTok{    domain: getLetterCounts(texts[selection]).map(d =\textgreater{} d.letter) }
\NormalTok{  \},}
\NormalTok{  y: \{ }
\NormalTok{    label: "Počet výskytů", }
\NormalTok{    grid: true,}
\NormalTok{    tickFormat: "d" // Odsud zmizí desetinná čísla, zůstanou jen celá (1, 2, 3...)}
\NormalTok{  \},}
\NormalTok{  marks: [}
\NormalTok{    Plot.barY(getLetterCounts(texts[selection]), \{}
\NormalTok{      x: "letter", }
\NormalTok{      y: "count", }
\NormalTok{      fill: "steelblue"}
\NormalTok{    \}),}
\NormalTok{    Plot.ruleY([0])}
\NormalTok{  ]}
\NormalTok{\})}
\end{Highlighting}
\end{Shaded}

Podobný způsob vizualizace můžeme provést i pro větší množství dat z
proměnné kardinální. Tam může nastat problém, že každá hodnota je
unikátní, takže vykreslení každé hodnoty by nám mnoho nepomohlo. Pro
takovéto případy se hodnoty \textbf{agregují}, tedy slučují do několika
\textbf{kategorií} (\emph{bins}) a následně se vykresluje četnost v dané
kategorii. V následujícím grafu si můžete vyzkoušt různá data (body u
osy x) a různý počet kategorií a histogram vytvořený podle těchto
kdategorií.

\begin{Shaded}
\begin{Highlighting}[]
\NormalTok{// 1. Generování dat na základě výběru}
\NormalTok{data02 = \{}
\NormalTok{  let v;}
\NormalTok{  const n = 500;}
\NormalTok{  if (distType02 === "Normální") \{}
\NormalTok{    v = Array.from(\{length: n\}, d3.randomNormal(0, 1));}
\NormalTok{  \} else if (distType02 === "Zešikmené") \{}
\NormalTok{    // Log{-}normální rozdělení je přirozeně zešikmené}
\NormalTok{    v = Array.from(\{length: n\}, d3.randomLogNormal(0, 0.5));}
\NormalTok{  \} else if (distType02 === "Uniformní") \{}
\NormalTok{    v = Array.from(\{length: n\}, d3.randomUniform({-}2, 2));}
\NormalTok{  \}}
\NormalTok{  return v.map(d =\textgreater{} (\{val: d\}));}
\NormalTok{\}}

\NormalTok{// 2. Ovládací prvky}
\NormalTok{viewof distType02 = Inputs.select(}
\NormalTok{  ["Normální", "Zešikmené", "Uniformní"], }
\NormalTok{  \{label: "Typ rozdělení"\}}
\NormalTok{)}

\NormalTok{viewof binCount = Inputs.range([2, 200], \{value: 20, step: 1, label: "Počet bins"\})}

\NormalTok{// 3. Graf}
\NormalTok{Plot.plot(\{}
\NormalTok{  height: 400,}
\NormalTok{  grid: true,}
\NormalTok{  marks: [}
\NormalTok{    // Histogram (Density)}
\NormalTok{    Plot.rectY(data02, Plot.binX(\{y: "count"\}, \{x: "val", thresholds: binCount, fill: "steelblue", fillOpacity: 0.7\})),}
    
\NormalTok{  // dodgeY vytvoří "mrak" bodů, anchor: "bottom" je drží u spodní hrany}
\NormalTok{    Plot.dot(data02, Plot.dodgeY(\{}
\NormalTok{      x: "val", }
\NormalTok{      anchor: "bottom", }
\NormalTok{      padding: 0.5, // Tímto zvětšuješ/zmenšuješ rozestupy mezi body}
\NormalTok{      r: 2, }
\NormalTok{      fill: "black",}
\NormalTok{      fillOpacity: 0.4}
\NormalTok{    \})),}
    
\NormalTok{    Plot.ruleY([0])}
\NormalTok{  ],}
\NormalTok{  y: \{}
\NormalTok{    label: "Četnost / Hustota",}
\NormalTok{    tickFormat: "d"}
\NormalTok{  \},}
\NormalTok{  x: \{}
\NormalTok{    label: "Hodnota x"}
\NormalTok{  \}}
\NormalTok{\})}
\end{Highlighting}
\end{Shaded}

Pro málo kategorií nám graf příliš neřekne, pro příliš mnoho je pak
nezachycuje trend v datech ale spíše odchylky od trendu. Proto se
stanovilo několik pravidel, jak ``optimálně'' nastavit počet kategorií.
Nejčastěji využívané je Sturgessovo pravidlo
\(k=\lfloor 1+\log_2(n)\rfloor\), kde uvedené závorky značí zaokrouhlení
na celá čísla dolů. Toto pravilo je však někdy nahrazováno jinými, viz
(R Core Team, 2024)

\section{Střední hodnota založená na
uspořádání}\label{stux159ednuxed-hodnota-zaloux17eenuxe1-na-uspoux159uxe1duxe1nuxed}

Mírou polohy založenou na uspořádání je \textbf{Medián (mdn)},
prostřední hodnota v uspořádaném datovém souboru, případně polovina z
dvou prostředních hodnot, pokud máme sudý počet dat. V našem případě je
medián mezi 3 a 3, \(\{1, 2, 2, 2, \mathbf{ 3, 3}, 4, 5, 8, 11\}\) tedy
\(\text{mdn} = \frac{3 + 3}{2} = 3\).

Můžeme tedy říci, že 50 \% domácností má tři nebo méně aut.

\begin{tcolorbox}[enhanced jigsaw, arc=.35mm, bottomrule=.15mm, bottomtitle=1mm, breakable, colback=white, colbacktitle=quarto-callout-note-color!10!white, colframe=quarto-callout-note-color-frame, coltitle=black, left=2mm, leftrule=.75mm, opacityback=0, opacitybacktitle=0.6, rightrule=.15mm, title=\textcolor{quarto-callout-note-color}{\faInfo}\hspace{0.5em}{Poznámka}, titlerule=0mm, toprule=.15mm, toptitle=1mm]

Medián i další kvantily můžeme určit pro typy proměnných, které mají
alespoň uspořádání, tedy proměnné ordinální a kardinální.

\end{tcolorbox}

Máme tedy tři možnosti jak určit střed, porovnání těchto tří možností
jak určit střed našich dat (středí hodnoty) je v následujícím grafu

\begin{Shaded}
\begin{Highlighting}[]
\NormalTok{//| echo: false}
\NormalTok{data = [1, 2, 2, 2, 3, 3, 4, 5, 8, 11].map(d =\textgreater{} (\{ hodnota: d \}))}

\NormalTok{// Statistiky pro čáry a legendu}
\NormalTok{stats = [}
\NormalTok{  \{ label: "Modus", val: 2, color: "green", dash: "2,2", yOffset: 15 \},}
\NormalTok{  \{ label: "Medián", val: 3, color: "blue", dash: "2.2", yOffset: 35 \},}
\NormalTok{  \{ label: "Průměr", val: 4.1, color: "red", dash: "2.2", yOffset: 55 \}}
\NormalTok{]}

\NormalTok{Plot.plot(\{}
\NormalTok{  width: 600,}
\NormalTok{  height: 200,}
\NormalTok{  grid: true,}
\NormalTok{  x: \{  domain: [0, 12], }
\NormalTok{        ticks: d3.range(1, 12), }
\NormalTok{        label: "Hodnota →" \},}
\NormalTok{  y: \{  axis: null, }
\NormalTok{        domain: [0, 3] \}, }
\NormalTok{  marks: [}
\NormalTok{    // Tečky nad sebou}
\NormalTok{    Plot.dot(data, Plot.stackY(\{ x: "hodnota", fill: "gray", stroke: "white", r: 6 \})),}

\NormalTok{    // Svislé čáry na pozicích statistik}
\NormalTok{    stats.map(s =\textgreater{} Plot.ruleX([s.val], \{ stroke: s.color, strokeWidth: 1, strokeDasharray: s.dash \})),}

\NormalTok{    // Popisky v pravém horním rohu pod sebou}
\NormalTok{    stats.map(s =\textgreater{} }
\NormalTok{      Plot.text([0], \{ // [0] je zde jen zástupný, pozici určí frameAnchor}
\NormalTok{        text: [\textasciigrave{}$\{s.label\}: $\{s.val\}\textasciigrave{}],}
\NormalTok{        frameAnchor: "top{-}right",}
\NormalTok{        dx: {-}20,        // odsazení od pravého okraje}
\NormalTok{        dy: s.yOffset,  // vertikální rozestup mezi popisky}
\NormalTok{        fill: s.color,}
\NormalTok{        // fontWeight: "bold",}
\NormalTok{        fontSize: 14,}
\NormalTok{        textAnchor: "end"}
\NormalTok{      \})}
\NormalTok{    )}
\NormalTok{  ]}
\NormalTok{\})}
\end{Highlighting}
\end{Shaded}

\section*{Kvantily - další míry založené na
uspořádání}\label{kvantily---dalux161uxed-muxedry-zaloux17eenuxe9-na-uspoux159uxe1duxe1nuxed}

\markright{Kvantily - další míry založené na uspořádání}

Míra polohy však nemusí kazovat na střed jako \textbf{medián (mdn)}, ale
také na krajní hodnoty: \textbf{minimum (min)}, \textbf{maximum (max)}.
Pokud chceme podrobnější dělení našich dat, pak můžeme využít i další
kvantily. \textbf{Kvantil} je hodnota, ``pod kterou'' leží určité
procento dat (přesněji jsou menší nebo stejné). Důležité kvantily mají
svá spěcifická označení: \textbf{medián} pro polovinu, \textbf{kvartily}
pro čtvrtiny, \textbf{decily} pro desetiny a \textbf{percentily} pro
setiny. Proto platí že např medián = druhý kvartil = pátý decil =
padesátý percentil. Postup výpočtu se může lišit, my uvedeme ten nejvíce
používaný, který odpovídá funkci PERCENTIL.INC v MS EXCEL. Kvantil \(p\)
hledáme \textbf{v setříděném souboru} na pozici
\textbf{\(k=1+(n-1)\cdot p\)} a to buď přímo, pokud \(k\) je celé, nebo
interpolací, pokud \(k\) celé není. V druhém příapdě číslo \(k\)
rozdělíme na celou část \(c\) (index souseda vlevo) a desetinnou část
\(d\) (poměrná vzdálenost k sousedovi vpravo). Percentil je pak hodnota
na spojnici těchto dvou sousedů:
\textbf{\(x_{c} + (x_{c+1}-x_c) \cdot d\)}. Ukažme to na příkladu našich
dat.

\begin{tcolorbox}[enhanced jigsaw, arc=.35mm, bottomrule=.15mm, bottomtitle=1mm, breakable, colback=white, colbacktitle=quarto-callout-note-color!10!white, colframe=quarto-callout-note-color-frame, coltitle=black, left=2mm, leftrule=.75mm, opacityback=0, opacitybacktitle=0.6, rightrule=.15mm, title=\textcolor{quarto-callout-note-color}{\faInfo}\hspace{0.5em}{Určete první kvartil a osmý decil}, titlerule=0mm, toprule=.15mm, toptitle=1mm]

Pro první kvartil je \(p=0.25\). Vypočeteme pořadí
\(k=1+(10-1)\cdot 0.25=3.25\), protože to není číslo celé, rozdělíme ho
\(3.25 = 3 + 0.25\). Požadovaný kvartil tedy hledáme mezi třetí a
čtvrtou, \(\{1,  2, \mathbf{ 2, 2}, 3, 3, 4, 5, 8, 11\}\). Můžeme použít
vzorec \(x_3+(x_4-x_3)\cdot 0.25\), ale protože \(x_3=x_4=2\), tak
hodnota prvního kvartilu je \(Q_1=2\).

Pro osmý decil popíšeme postup stručněji v bodech

\begin{itemize}
\tightlist
\item
  osmý decil \(\Longrightarrow p=0.8, n=10\)
\item
  \(k=1+(10-1)\cdot 0.8=8.2 \Longrightarrow 8.2 = 8 + 0.2\).
\item
  \(\{1, 2, 2, 2, 3, 3, 4, \mathbf{ 5, 8}, 11\}\)
\item
  \(x_8+(x_9-x_8)\cdot 0.2 \Longrightarrow 5+(8-5)\cdot0.2\)
\item
  \(D_8=5,6\).
\end{itemize}

v grafu na vodorovné ose hodnoty, na svislé ose pořadí, ukázka lineární
interpolace když na svislou osu vyneseme 8.2 a pomocí úsečky získáme
hodnotu 5.6 v datech.

\begin{Shaded}
\begin{Highlighting}[]
\NormalTok{//| echo: false}

\NormalTok{// 1. Definice bodů}
\NormalTok{body\_data = [}
\NormalTok{  \{ x: 4, y: 7, label: "x₇ = 4" \},}
\NormalTok{  \{ x: 5, y: 8, label: "x₈ = 5" \},}
\NormalTok{  \{ x: 8, y: 9, label: "x₉ = 8" \}}
\NormalTok{]}
\NormalTok{// Cílové hodnoty}
\NormalTok{target\_y = 8.2}
\NormalTok{target\_x = 5.6}
\NormalTok{Plot.plot(\{}
\NormalTok{  width: 600,}
\NormalTok{  height: 200,}
\NormalTok{  grid: true,}
\NormalTok{  // Rozsah osy X s trochou místa pro popisek x7}
\NormalTok{  x: \{ domain: [3, 9], label: "Hodnota (osa x) →" \},}
\NormalTok{  y: \{ domain: [6.8, 9.2], label: "Pořadí (index) ⭡", ticks: [7, 8, 8.2, 9] \},}
\NormalTok{  marks: [}
\NormalTok{    // Lomená čára propojující body x7 {-}\textgreater{} x8 {-}\textgreater{} x9}
\NormalTok{    Plot.line(body\_data, \{x: "x", y: "y", stroke: "\#ccc", strokeWidth: 1\}),}

\NormalTok{    // Svislá pomocná čára (průmět na osu X)}
\NormalTok{    Plot.ruleX([target\_x], \{y1: 6.8, y2: target\_y, stroke: "darkred", strokeDasharray: "4,4" \}),}
    
\NormalTok{    // Vodorovná pomocná čára (zadání indexu 8.2)}
\NormalTok{    Plot.ruleY([target\_y], \{x1: 3.5, x2: target\_x, stroke: "grey", strokeDasharray: "4,4"\}),}
    
\NormalTok{    // Body x7, x8 a x9 v grafu (v prostoru)}
\NormalTok{    Plot.dot(body\_data, \{x: "x", y: "y", r: 5, fill: "gray", stroke: "white"\}),}
\NormalTok{    Plot.text(body\_data, \{x: "x", y: "y", text: "label", dy: {-}15, fontWeight: "bold"\}),}

\NormalTok{    // Cílový bod (8. decil) na spojnici}
\NormalTok{    Plot.dot([\{x: target\_x, y: target\_y\}], \{x: "x", y: "y", r: 5, fill: "darkred"\}),}
    
\NormalTok{    // Původní data na ose x (šedá/modrá)}
\NormalTok{    Plot.dot(body\_data, \{x: "x", y: 6.8, r: 8, fill: "steelblue", stroke: "white" \}),}
    
\NormalTok{    // Zvýraznění výsledné hodnoty přímo na ose X (červená tečka a číslo)}
\NormalTok{    Plot.dot( [\{x: target\_x, y: 6.8\}], \{x: "x", y: "y", r: 8, fill: "darkred", stroke: "white", strokeWidth: 2 \}),}
\NormalTok{    // Popisek u cílového bodu na spojnici }
\NormalTok{    Plot.text([\{x: target\_x, y: 6.8\}], \{x: "x", y: "y", text: ["osmý decil"], dy: 15, fontWeight: "bold", fill: "darkred"\}),}
\NormalTok{    // Plot.text([\{x: target\_x, y: 6.8\}], \{ text: [target\_x.toFixed(1)], fill: "darkred", fontWeight: "bold", dy: 20, fontSize: 14 \})}
\NormalTok{  ]}
\NormalTok{\})}
\end{Highlighting}
\end{Shaded}

\end{tcolorbox}

Pokud využijeme vzorce pro celý rozsah dat, získáme \textbf{kvantilovou
funkci} \(q = Q(p)\), která popisuje empirircky (z dat) popisuje vztah
hodnotou \textbf{\(q\)} nazývanou \textbf{kvantil, percentil, \ldots{}}
a hodnotou \textbf{\(p\in(0,1)\)} nazývaný \textbf{hladina,
(kumulativní) pravděpodobnost, percentuální podíl, \ldots{}}. Z pohledu
pravděpodobnosti je tedy \(q\) taková hodnota, že pravděpodobnost že při
náhodném výběru je vybrané hodnota x menší nebo rovna \(q\) je rovna
\(p\).

\[ Q(p)=q \Longleftrightarrow P(x\leq q)=p\]

\begin{tcolorbox}[enhanced jigsaw, arc=.35mm, bottomrule=.15mm, bottomtitle=1mm, breakable, colback=white, colbacktitle=quarto-callout-note-color!10!white, colframe=quarto-callout-note-color-frame, coltitle=black, left=2mm, leftrule=.75mm, opacityback=0, opacitybacktitle=0.6, rightrule=.15mm, title=\textcolor{quarto-callout-note-color}{\faInfo}\hspace{0.5em}{Poznámka}, titlerule=0mm, toprule=.15mm, toptitle=1mm]

Interpretačně nám odpovídá na otázky například \emph{jaký plat má
spodních 10\% ve firmě?} (\(p=0.10, q=...\)) Nebo pro inverzní funkci
\emph{kolik procent osob ve firmě má plat menší než 40 000?}
(\(q=40000, p=...\))\$

\end{tcolorbox}

\begin{Shaded}
\begin{Highlighting}[]
\NormalTok{//| echo: false}
\NormalTok{\{}
\NormalTok{  const data = [1, 2, 2, 2, 3, 3, 4, 5, 8, 11];}
\NormalTok{  const n = data.length;}

\NormalTok{  // Vytvoření dvojic [p, hodnota] pro každý prvek}
\NormalTok{  // index = (n{-}1)*p + 1  =\textgreater{}  p = (index {-} 1) / (n {-} 1)}
\NormalTok{  const points = data.map((d, i) =\textgreater{} [i / (n {-} 1), d]);}

\NormalTok{  return Plot.plot(\{}
\NormalTok{    width: width,}
\NormalTok{    height: width * 0.8, // Snaha o čtvercový poměr}
\NormalTok{    grid: true,}
\NormalTok{    x: \{}
\NormalTok{      label: "p (kvantil)",}
\NormalTok{      domain: [0, 1],}
\NormalTok{      ticks: 10}
\NormalTok{    \},}
\NormalTok{    y: \{}
\NormalTok{      label: "Hodnota",}
\NormalTok{      domain: [0, 11],}
\NormalTok{      ticks: 11}
\NormalTok{    \},}
\NormalTok{    marks: [}
\NormalTok{      // Spojnice (lineární interpolace)}
\NormalTok{      Plot.line(points, \{}
\NormalTok{        x: d =\textgreater{} d[0], }
\NormalTok{        y: d =\textgreater{} d[1], }
\NormalTok{        stroke: "steelblue", }
\NormalTok{        strokeWidth: 3}
\NormalTok{      \}),}
\NormalTok{      // Samotné body}
\NormalTok{      Plot.dot(points, \{}
\NormalTok{        x: d =\textgreater{} d[0], }
\NormalTok{        y: d =\textgreater{} d[1], }
\NormalTok{        fill: "steelblue", }
\NormalTok{        r: 8}
\NormalTok{      \}),}
\NormalTok{      // Pravítko na nule}
\NormalTok{      Plot.ruleY([0])}
\NormalTok{    ]}
\NormalTok{  \});}
\NormalTok{\}  }
\end{Highlighting}
\end{Shaded}

\begin{tcolorbox}[enhanced jigsaw, arc=.35mm, bottomrule=.15mm, bottomtitle=1mm, breakable, colback=white, colbacktitle=quarto-callout-note-color!10!white, colframe=quarto-callout-note-color-frame, coltitle=black, left=2mm, leftrule=.75mm, opacityback=0, opacitybacktitle=0.6, rightrule=.15mm, title=\textcolor{quarto-callout-note-color}{\faInfo}\hspace{0.5em}{Poznámka}, titlerule=0mm, toprule=.15mm, toptitle=1mm]

Jiné přístupy se liší v tom, jak vypočtou \(k\) (např.
\(k=(n+1)\cdot p\)) případně jak pracují s desetinnou částí (například
ji ignorují tedy percentil na pozici \(k=6.3\) je vyjádřen hodnotou
\(x_6\)). Podrobněji viz manuál R, (R Core Team, 2023).

\end{tcolorbox}

Doposud jsme se ptali, \textbf{kde} leží střed, minimum, čtvrtina,
\ldots. Druhou otázkou je, jak jsou data ``nahuštěna'' kolem střední
hodnoty nebo naopak ``rozprostřena''. To se snaží popsat \textbf{míry
variability}.

\section{Míry variabiliy odvozená z
kvantilů}\label{muxedry-variabiliy-odvozenuxe1-z-kvantilux16f}

Základní idea vychází z uspořádaného souboru a pořadí jednotivých
hodnot.

\textbf{Rozpětí} (\emph{range}), někdy také \textbf{variační rozpětí}:
rozdíl mezi největší a nejmenš hodnotou, \textbf{\(range = max-min\)}.

\textbf{Mezikvartilové rozpětí} (\emph{inter-quartile-range, IQR}):
rozdíl mezi třetím kvartilem a prvním kvartilem,
\textbf{\(IQR=Q_3-Q_1\)}.

\begin{tcolorbox}[enhanced jigsaw, arc=.35mm, bottomrule=.15mm, bottomtitle=1mm, breakable, colback=white, colbacktitle=quarto-callout-note-color!10!white, colframe=quarto-callout-note-color-frame, coltitle=black, left=2mm, leftrule=.75mm, opacityback=0, opacitybacktitle=0.6, rightrule=.15mm, title=\textcolor{quarto-callout-note-color}{\faInfo}\hspace{0.5em}{Poznámka}, titlerule=0mm, toprule=.15mm, toptitle=1mm]

\textbf{Range} nám říká, kolik ``prostoru'' data zabírají. Pozor na to,
že rozsah, min a max se nemusí shodovat z rozsahem škály. Například
měřené hodnoty mohou mít maximální rozsah 0-20, ale data mohou obsahovat
hodnoty 1-11 a tedy min=1 a max=11, range = 10.

\textbf{IQR} nám říká, kolik prostoru si zabírá ``střední polovina''
hodnot, tedy když ignorujeme čtvrtinu nejmenších a čtvrtinu největších
hodnot.

\end{tcolorbox}

V našem příkladu je maximální rozdíl mezi domácnostmi \(10\) vozidel, v
prostřední polovině je maximální rozdíl \(3\) vozidla.

\section{Odlehlé hodnoty a
box-plot}\label{odlehluxe9-hodnoty-a-box-plot}

Často je třeba stanovit, zda v souboru nejsou data, která výzraně
vybočují od středu. Ve statistice se nejčastěji používá rozsah
\(4\cdot \operatorname{IQR}\) pro hodnoty akceptovatelné - prostřední
IQR máme dané kvartliy, na každou stranu přidáme ještě
\(1.5\cdot \operatorname{IQR}\) a získáme hranice, hradby. \textbf{Dolní
hradbu} (\emph{Lower fence, LF}) určíme jako
\(Q1-1.5\cdot \operatorname{IQR}\), \textbf{horní hradbu} (\emph{Upper
fence, UF}) jako \(Q3+1.5\cdot
IQR\). Z našich dat

\begin{align}
    k_{0.25}      &= 1+(10-1)\cdot 0.25= 3.25\\
    Q_1           &= x_3+(x_4-x_3)\cdot 0.25= 2\\
    k_{0.75}      &= 1+(10-1)\cdot 0.25= 7.75\\
    Q_3           &= x_7+(x_8-x_7)\cdot 0.75= 4+(5-4)\cdot 0.75= 4.75\\
    IQR           &= 4.75-2= 2.75\\
    \mathbf{LF}   &= Q_1 - 1,5 \cdot \operatorname{IQR} = 2 - 1.5 \cdot 2.75 = -2.125\\
    \mathbf{UF}   &=Q_3 + 1.5 \cdot \operatorname{IQR} = 4.75 + 1.5 \cdot 2.75 = 8.875
\end{align}

V našich datech tedy nemáme žádné odlehlé hodnoty pod dolní hradbou ale
jednu odlehlou hodnoty nad horní hradbou.

Důležité kvantily a odlehlé hodnoty se často znázorňují do grafu, který
se nazývá \textbf{box-plot}: obdélník ukazuje na 1Q, med a 3Q, levý
``vous'' ukazuje buď minimum, nebo dolní hradbu, pravý ``vous'' maximum
nebo horní hradbu, červenou tečkou je znázorněna odlehlá hodnota. Můžete
se setkat s různými variantami, ale tato je nejčastější.

\begin{Shaded}
\begin{Highlighting}[]
\NormalTok{//| echo: false}

\NormalTok{const data\_base=[1,2,2,2,3,3,4,5,8,11]}

\NormalTok{// Funkce pro výpočet percentilu metodou R{-}Type 7 (1 + (n{-}1)p)}
\NormalTok{function quantileR7(data, p) \{}
\NormalTok{  const sorted = data.slice().sort(d3.ascending);}
\NormalTok{  const n = sorted.length;}
\NormalTok{  if (n === 0) return undefined;}
\NormalTok{  if (p \textless{}= 0) return sorted[0];}
\NormalTok{  if (p \textgreater{}= 1) return sorted[n {-} 1];}

\NormalTok{  // Teoretický index (rank {-} 1)}
\NormalTok{  const h = (n {-} 1) * p;}
\NormalTok{  const i = Math.floor(h);}
\NormalTok{  const a = sorted[i];}
\NormalTok{  const b = sorted[i + 1];}

\NormalTok{  // Lineární interpolace}
\NormalTok{  return a + (h {-} i) * (b {-} a);}
\NormalTok{\}}

\NormalTok{// Výpočty statistik}
\NormalTok{q1 = quantileR7(data\_base, 0.25)}
\NormalTok{q2 = quantileR7(data\_base, 0.5)}
\NormalTok{q3 = quantileR7(data\_base, 0.75)}

\NormalTok{iqr = q3 {-} q1}
\NormalTok{lowerWhisker = Math.max(d3.min(data\_base), q1 {-} 1.5 * iqr)}
\NormalTok{upperWhisker = Math.min(d3.max(data\_base), q3 + 1.5 * iqr)}
\NormalTok{outliers = data\_base.filter(d =\textgreater{} d \textless{} lowerWhisker || d \textgreater{} upperWhisker)}

\NormalTok{// Graf}
\NormalTok{Plot.plot(\{}
\NormalTok{  width: 600,}
\NormalTok{  height: 200,}
\NormalTok{  x: \{domain: [0, 12], grid: true\},}
\NormalTok{  y: \{domain: [{-}1, 3], axis: null\},}
\NormalTok{  marks: [}
\NormalTok{    // Stackované body}
\NormalTok{    Plot.dot(data\_base.map(d =\textgreater{} (\{x: d\})), Plot.stackY(\{}
\NormalTok{      x: "x",}
\NormalTok{      fill: "steelblue",}
\NormalTok{      stroke: "white",}
\NormalTok{      r: 6}
\NormalTok{    \})),}
    
\NormalTok{    // Boxplot (box)}
\NormalTok{    Plot.rect([\{\}], \{}
\NormalTok{      x1: q1, x2: q3, y1: {-}0.3, y2: 0.3,}
\NormalTok{      fill: "lightblue", stroke: "black", fillOpacity: 0.5}
\NormalTok{    \}),}
    
\NormalTok{    // Medián}
\NormalTok{    Plot.ruleX([q2], \{y1: {-}0.3, y2: 0.3, stroke: "black", strokeWidth: 2\}),}
    
\NormalTok{    // Vousy (Linky)}
\NormalTok{    Plot.link([\{x1: q1, x2: lowerWhisker\}], \{x1: "x1", x2: "x2", y: 0, stroke: "black"\}),}
\NormalTok{    Plot.link([\{x1: q3, x2: upperWhisker\}], \{x1: "x1", x2: "x2", y: 0, stroke: "black"\}),}
    
\NormalTok{    // Konce vousů}
\NormalTok{    Plot.ruleX([lowerWhisker, upperWhisker], \{y1: {-}0.2, y2: 0.2, stroke: "black"\}),}
    
\NormalTok{    // Odlehlé hodnoty (zobrazeny červeně)}
\NormalTok{    Plot.dot(outliers.map(d =\textgreater{} (\{x: d\})), \{x: "x", y: 0, r: 4, fill: "red"\})}
\NormalTok{  ]}
\NormalTok{\})}
\end{Highlighting}
\end{Shaded}

\section*{Klíčové poznatky}\label{kluxedux10dovuxe9-poznatky}
\addcontentsline{toc}{section}{Klíčové poznatky}

\markright{Klíčové poznatky}

{\def\LTcaptype{none} % do not increment counter
\begin{longtable}[]{@{}
  >{\raggedright\arraybackslash}p{(\linewidth - 2\tabcolsep) * \real{0.5000}}
  >{\raggedright\arraybackslash}p{(\linewidth - 2\tabcolsep) * \real{0.5000}}@{}}
\toprule\noalign{}
\endhead
\bottomrule\noalign{}
\endlastfoot
Typy proměnných & nominální, ordinální, kardinální (intervalová,
poměrová) \\
co popisují míry & poloha (kde data jsou), variabilita (jak jsou
rozložená) \\
průměr & \(\bar{x}=\frac{\sum x_i}{n}\) \\
modus & nejčetněješí hodnota \\
medián & střední hodnota uspořádaného souboru \\
kvantily & poloha určitého procenta dat, kvartily, percentily \\
index kvantilu & \(k=1+(n-1)p\) \\
kvantil & \(q_p=x_c+(x_{c+1}-x_c)\cdot d\), kde \(k=c+d\), celá část +
desetinná část \\
variabilita kvantilová & rozsah \(range = max-min\), rozsah střední
poloviny \(IQR=Q_3-Q_1\) \\
odlehlé hodnoty & box plot, odlehlé hodnoty, hranice \\
\end{longtable}
}

\chapter{Momentové statistiky}\label{momentovuxe9-statistiky}

Dosud jsem po proměné požadovali aby měla uspořádání, tedy byla alespoň
ordinální. To nám umožnilo definovat medián, kvantily a viualizaci
rozloření dat pomocí box-plot. Abychom mohli přejít k aritmetickému
průměru a statistikám z něj odvozeným, musíme přejít do vyšší úrovně
měření, tedy k proměnným \textbf{kardinálním poměrovým}.

\section{Momentová míra polohy}\label{momentovuxe1-muxedra-polohy}

Jak již bylo zmíněno, nejznámější statistikou je \textbf{(aritmetický)
průměr} označovaný jako \textbf{\(\bar{x}\)} nebo \textbf{m (mean)}.
Jedná se o statistiku momentovou, ale co to vlastně znamená?

Slovo \textbf{moment} pochází původně z fyziky, kde popisuje situaci na
páce - \textbf{moment síly je součin hmotnosti \(x_i\) a vzdálenosti
\(d_i\) od bodu otáčení}. Artimetický průměr v této situaci ukazuje na
\textbf{těžiště}, tedy když páku podepřeme právě v tomto místě, bude v
rovnováze. Hmotnost je ve statistice nahrazena relativní četností
\(p_i\), vzdálenost je nahrazena pojmem \textbf{odchylka od průměru}
\(d_i\).

\begin{Shaded}
\begin{Highlighting}[]
\NormalTok{//| echo: false}
\NormalTok{//| label: fig{-}teziste{-}vizualizace}
\NormalTok{data\_base = [1, 2, 2, 2, 3, 3, 4, 5, 8, 11]}
\NormalTok{mean\_vypocet = d3.sum(data\_base) / data\_base.length}

\NormalTok{Plot.plot(\{}
\NormalTok{  width: 600,}
\NormalTok{  height: 180,}
\NormalTok{  grid: true,}
\NormalTok{  x: \{ domain: [0, 12], label: "Hodnota (poloha na páce) →" \},}
\NormalTok{  y: \{ axis: null, domain: [0, 4] \},}
\NormalTok{  marks: [}
\NormalTok{    // Tyč páky}
\NormalTok{    Plot.ruleY([0.5], \{ strokeWidth: 2, stroke: "\#333" \}),}

\NormalTok{    // Závaží (datové body) naskládané na sebe}
\NormalTok{    Plot.dot(data\_base.map(d =\textgreater{} (\{x: d\})), Plot.stackY(\{}
\NormalTok{      x: "x",}
\NormalTok{      y: 0.3,}
\NormalTok{      dy: {-}10,}
\NormalTok{      fill: "steelblue",}
\NormalTok{      stroke: "white",}
\NormalTok{      r: 6}
\NormalTok{    \})),}

\NormalTok{    // Těžiště (Průměr) jako podpěra páky}
\NormalTok{    Plot.dot([\{x: mean\_vypocet, y: 0.5\}], \{}
\NormalTok{      x: "x",}
\NormalTok{      y: "y",}
\NormalTok{      symbol: "triangle",}
\NormalTok{      r: 10,}
\NormalTok{      rotate: 0, // otočení trojúhelníku, aby vypadal jako podstavec}
\NormalTok{      fill: "darkred",}
\NormalTok{      dy: 15}
\NormalTok{    \}),}

\NormalTok{    // Popisek těžiště}
\NormalTok{    Plot.text([\{x: mean\_vypocet, y: 0\}], \{}
\NormalTok{      x: "x",}
\NormalTok{      y: "y",}
\NormalTok{      text: [\textasciigrave{}Těžiště (průměr) = $\{mean\_vypocet.toFixed(1)\}\textasciigrave{}],}
\NormalTok{      dy: 25,}
\NormalTok{      fill: "red",}
\NormalTok{      fontWeight: "bold"}
\NormalTok{    \})}
\NormalTok{  ]}
\NormalTok{\})}
\end{Highlighting}
\end{Shaded}

\section{Momentové míry
variability}\label{momentovuxe9-muxedry-variability}

Pojďme se podívat na to, jak na lze charakterizovat variabilitu dat
vzhledem k průměru. Můžeme zjistit, jak daleko od průměru jsou
jednotlivé body a získáme \textbf{odchylky jednotlivých hodnot od
průměru}, \(e_i=x_i-\bar{x}\). Naším zájmem je získat \textbf{průměrnou
odchylku}, což vypadá snadno: sečteme všechny odchylky a vydělíme
počtem. Jenže (!) problém je ten, že když je sečteme, dostaneme vždy
nulu.

\begin{tcolorbox}[enhanced jigsaw, arc=.35mm, bottomrule=.15mm, bottomtitle=1mm, breakable, colback=white, colbacktitle=quarto-callout-note-color!10!white, colframe=quarto-callout-note-color-frame, coltitle=black, left=2mm, leftrule=.75mm, opacityback=0, opacitybacktitle=0.6, rightrule=.15mm, title=\textcolor{quarto-callout-note-color}{\faInfo}\hspace{0.5em}{Odvození, že \(\sum e_i=0\)}, titlerule=0mm, toprule=.15mm, toptitle=1mm]

\begin{align}
  \sum(e_i)         &= \sum(x_i-\bar{x})                      \\
                    &= (x_1-\bar{x}) + (x_2-\bar{x})+\dots    \\
                    &= x_1+x_2+\dots -(\bar{x}+\bar{x}+\dots) \\
                    &= \sum x_i-\sum \bar{x}                  \\
                    &= \sum x_i-n\cdot \bar{x}               \\
                    &=\sum x_i-n\cdot \frac{\sum x_i}{n}      \\
                    &=0
    \end{align}

\end{tcolorbox}

\begin{tcolorbox}[enhanced jigsaw, arc=.35mm, bottomrule=.15mm, bottomtitle=1mm, breakable, colback=white, colbacktitle=quarto-callout-note-color!10!white, colframe=quarto-callout-note-color-frame, coltitle=black, left=2mm, leftrule=.75mm, opacityback=0, opacitybacktitle=0.6, rightrule=.15mm, title=\textcolor{quarto-callout-note-color}{\faInfo}\hspace{0.5em}{Poznámka}, titlerule=0mm, toprule=.15mm, toptitle=1mm]

To že součet odchylek je nulový ukazuje na to, že se opravdu jedná o
těžiště: součet odchylek záporných = součet odchylek kladných, žádná
strana není zvýhodněna.

\end{tcolorbox}

Abychom získali průměrnou odchylku můžeme buď použít absolutní hodnoty a
získat tak hodnotu zvanou \textbf{\emph{mean absolute deviance, MAD}} :
\(MAD=\frac{\sum |x_i-\bar{x}|}{n}\), Nebo spočítat nejprve kvadratické
odchylky \((x_i-\bar{x})^2\), určit průměrnou kvadratickou odchylku
\(\frac{\sum (x_i-\bar{x})^2}{n}\) a poté ji odmocnit. Získáme tím
hodnotu které se říká \textbf{směrodatná odchylka} a značí se symbolem
\emph{sd (standard deviation)}

\[ sd=\sqrt{\frac{\sum (x_i-\bar{x})^2}{n}} \]

Koncept směrodatné odchylky lze získat také jinou, geometrickou cestou.
Pro zájemce ji zde uvedeme.

\begin{tcolorbox}[enhanced jigsaw, arc=.35mm, bottomrule=.15mm, bottomtitle=1mm, breakable, colback=white, colbacktitle=quarto-callout-note-color!10!white, colframe=quarto-callout-note-color-frame, coltitle=black, left=2mm, leftrule=.75mm, opacityback=0, opacitybacktitle=0.6, rightrule=.15mm, title=\textcolor{quarto-callout-note-color}{\faInfo}\hspace{0.5em}{Směrodatná odchylka lineárně algebraicky}, titlerule=0mm, toprule=.15mm, toptitle=1mm]

Abychom mohli vizualizovat potřebujeme jednodušší příklad: mějme dvě
domácnosti, kteří omají počet vozidel v rozsahu 0-10. Na ose x budeme
mít počet u první domácnosti, na ose y druhé. V defaultím nastavení má
první 7 vozidel, druhá 3 vozidla, což reprezentuje modrý bod, který má
souřadice {[}7,3{]}. Pro lepší vyjádření použijeme vektorovou formu
zápisu, která je bližší i sloupcům v tabulce dat
\(\mathbf{x}=\begin{pmatrix} x_1\\ x_2\end{pmatrix}=\begin{pmatrix} 7\\ 3\end{pmatrix}\)

Naším zájmem je nalézt jedno číslo, které charakterizuje naše data.
Fakticky tedy řešíme soustavy rovnic

\begin{align} x_1&=c\\ x_2&=c\end{align}

Tato soustava však v praktických situacích nemá řešení. Zkusme se však
na situaci podívat pohledem geometrie:

\begin{align} x_1&=c \\ x_2&=c \end{align}

\[
  \Longrightarrow \begin{pmatrix} x_1\\x_2 \end{pmatrix} =c \cdot \begin{pmatrix} 1 \\ 1 \end{pmatrix} \Longrightarrow \mathbf{ x}=c \cdot\mathbf{1}
\]

Na levé straně máme vektor dat, na pravé straně neznámý násobek vektoru
jedničkového.

\textbf{Na přímce dané jedničkovým vektorem tedy hledáme takové c, aby
výsledný vektor byl co nejblíže vektoru našich dat.}

\begin{Shaded}
\begin{Highlighting}[]
\NormalTok{//| echo: false}

\NormalTok{// Statické hodnoty}
\NormalTok{x1\_static = 7}
\NormalTok{x2\_static = 3}
\NormalTok{pru\_static = (x1\_static + x2\_static) / 2}

\NormalTok{// Vizualizace}
\NormalTok{Plot.plot(\{}
\NormalTok{  grid: true,}
\NormalTok{  width: 500,}
\NormalTok{  height: 500,}
\NormalTok{  x: \{domain: [0, 10], label: "Osa x₁"\},}
\NormalTok{  y: \{domain: [0, 10], label: "Osa x₂"\},}
\NormalTok{  marks: [}
\NormalTok{    // Přímka y = x}
\NormalTok{    Plot.line([[0,0], [10,10]], \{stroke: "darkgrey", strokeDasharray: "4,4", strokeWidth: 2\}),}
    
\NormalTok{    // VEKTOR jedniček}
\NormalTok{    Plot.link([\{x: pru\_static, y: pru\_static\}], \{}
\NormalTok{      x1: 0, y1: 0, x2: 1, y2: 1, }
\NormalTok{      stroke: "\#666", strokeWidth: 2,markerEnd: "arrow"}
\NormalTok{    \}),}
\NormalTok{    Plot.text([\{x: pru\_static/2, y: pru\_static/2\}], \{}
\NormalTok{      x: "x", y: "y", text: ["prostor generovaný 1"], }
\NormalTok{      rotate: {-}45, dy: 15, fill: "\#666"}
\NormalTok{    \}),}
    
\NormalTok{    // BOD DAT}
\NormalTok{    Plot.dot([\{x: x1\_static, y: x2\_static\}], \{}
\NormalTok{      x: "x", y: "y",}
\NormalTok{      fill: "steelblue",}
\NormalTok{      r: 5  // poloměr bodu}
\NormalTok{    \}),}
\NormalTok{    Plot.text([\{x: x1\_static, y: x2\_static\}], \{}
\NormalTok{      x: "x", }
\NormalTok{      y: "y", }
\NormalTok{      text: [\textasciigrave{}Data X = [$\{x1\_static.toFixed(1)\}, $\{x2\_static.toFixed(1)\}]\textasciigrave{}], }
\NormalTok{      dy: 20, dx: {-}5, }
\NormalTok{      fontWeight: "bold",}
\NormalTok{      fill: "steelblue"}
\NormalTok{    \})}
\NormalTok{  ]}
\NormalTok{\})}
\end{Highlighting}
\end{Shaded}

Provedeme lineárně algebraické úpravy rovnice\\
\begin{align} 
    c \cdot\mathbf{1} &=\mathbf{x} \\
    c \cdot\mathbf{1}^T \mathbf{1} &=\mathbf{1}^T\mathbf{x} \\
    c &= (\mathbf{1}^T \mathbf{1})^{-1}\mathbf{1}^T\mathbf{x}
  \end{align}

V našem případě však

\begin{align}
  \mathbf{1}^T \mathbf{1} &= (1,1) \cdot \begin{pmatrix} 1 \\ 1   \end{pmatrix} = 2 \\
  \mathbf{1}^T\mathbf{x}  &= (1,1) \cdot \begin{pmatrix} x_1\\x_2 \end{pmatrix} =\sum x_i\\
  c                       &= (\mathbf{1}^T \mathbf{1})^{-1} \cdot \mathbf{1}^T\mathbf{x}\\
                          &= 2^{-1} \cdot \sum x_i\\
                          &= \frac{\sum x_i}{2} \\
                          &= \bar{x}
  \end{align}

Aritmetický \textbf{průměr} tedy popisuje nejbliší bod na přímce
\textbf{1} od bodu našich dat \textbf{x}.

Můžeme tedy sestrojit vektor odchylek našich dat od vektoru průměru

\[
  \mathbf{e}=\mathbf{x}-\bar{x}\mathbf{1}=\begin{pmatrix} 7\\3\end{pmatrix}-5\begin{pmatrix}1\\ 1\end{pmatrix}=\begin{pmatrix} 2\\ -2\end{pmatrix}
\]

Vektor \(\mathbf{x}\) je tedy rozložen do vektoru do směru vektoru
\(\mathbf{1}\) a směru kolmého, takže platí vektorově

\[\mathbf{x}=\bar{x}\mathbf{1}+\mathbf{e}\]

Vektor dat jsme rozložili na dvě složky, složku průměru a složku
odchylek.

\begin{Shaded}
\begin{Highlighting}[]
\NormalTok{//| echo: false}
\NormalTok{// 1. Ovládací prvky (stále stejné)}
\NormalTok{viewof x1 = Inputs.range([0, 10], \{step: 1, value: 7, label: " Respondent 1"\})}
\NormalTok{viewof x2 = Inputs.range([0, 10], \{step: 1, value: 3, label: " Respondent 2"\})}
\NormalTok{// 2. Výpočty}
\NormalTok{pru = (x1 + x2) / 2}
\NormalTok{norma\_e = Math.sqrt(Math.pow(x1 {-} pru, 2) + Math.pow(x2 {-} pru, 2))}
\NormalTok{norma\_x = Math.sqrt(Math.pow(x1, 2) + Math.pow(x2, 2))}
\NormalTok{sigma\_val = norma\_e / Math.sqrt(2)}
\NormalTok{// 3. Vizualizace}
\NormalTok{Plot.plot(\{}
\NormalTok{  grid: true,}
\NormalTok{  width: 500,}
\NormalTok{  height: 500,}
\NormalTok{  x: \{domain: [0, 10], label: "Osa x₁"\},}
\NormalTok{  y: \{domain: [0, 10], label: "Osa x₂"\},}
\NormalTok{  marks: [}
\NormalTok{    // Přímka y = x}
\NormalTok{    Plot.line([[0,0], [10,10]], \{stroke: "\#ccc", strokeDasharray: "4,4"\}),}
    
\NormalTok{    // VEKTOR jedniček}
\NormalTok{    Plot.link([\{x: pru\_static, y: pru\_static\}], \{}
\NormalTok{      x1: 0, y1: 0, x2: 1, y2: 1, }
\NormalTok{      stroke: "\#666", strokeWidth: 3,markerEnd: "arrow"}
\NormalTok{    \}),}
\NormalTok{    // VEKTOR PRŮMĚRU (Projekce)}
\NormalTok{    Plot.link([\{x: pru, y: pru\}], \{x1: 0, y1: 0, x2: "x", y2: "y", stroke: "\#666", strokeWidth: 2\}),}
\NormalTok{    Plot.text([\{x: pru/2, y: pru/2\}], \{}
\NormalTok{      x: "x", y: "y", text: ["vektor průměru (x̄.1)"], }
\NormalTok{      rotate: {-}45, dy: 15, fill: "\#666"}
\NormalTok{    \}),}
    
\NormalTok{    // VEKTOR ODCHYLEK (Kolmice)}
\NormalTok{    Plot.link([\{x1: x1, y1: x2, x2: pru, y2: pru\}], \{}
\NormalTok{      x1: "x1", y1: "y1", x2: "x2", y2: "y2", }
\NormalTok{      stroke: "darkred", strokeWidth: 3, }
\NormalTok{    \}),}
\NormalTok{    Plot.text([\{x: (x1+pru)/2, y: (x2+pru)/2\}], \{}
\NormalTok{      x: "x", y: "y", text: ["vektor odchylek (e)"], }
\NormalTok{      rotate: 45, dx: 25, fill: "darkred"}
\NormalTok{    \}),    }
    
\NormalTok{    // Pravý úhel (symbol čtverce)}
\NormalTok{    Plot.dot([\{x: pru, y: pru\}], \{}
\NormalTok{        symbol: "square", r: 6, rotate: 45, fill: "none", stroke: "black", }
\NormalTok{        dx: (x1 \textgreater{} x2 ? {-}4 : 4), dy: (x1 \textgreater{} x2 ? 4 : {-}4)}
\NormalTok{    \}),}

\NormalTok{    // VEKTOR DAT (Přepona)}
\NormalTok{    Plot.link([\{x: x1, y: x2\}], \{x1: 0, y1: 0, x2: "x", y2: "y", stroke: "steelblue", strokeWidth: 1, markerEnd: "dot"\}),}
\NormalTok{    Plot.text([\{x: x1/2, y: x2/2\}], \{}
\NormalTok{      x: "x", y: "y", text: ["vektor dat (x)"], }
\NormalTok{      //rotate: Math.atan2(x2, x1) * (180/Math.PI), }
\NormalTok{      dy: {-}12, fill: "steelblue"}
\NormalTok{    \}),}
\NormalTok{    // BOD DAT}
\NormalTok{    Plot.dot([\{x: x1, y: x2\}], \{}
\NormalTok{      x: "x", y: "y",}
\NormalTok{      fill: "steelblue",}
\NormalTok{      r: 5  // poloměr bodu}
\NormalTok{    \}),    }
\NormalTok{    Plot.text([\{x: x1, y: x2\}], \{}
\NormalTok{      x: "x", }
\NormalTok{      y: "y", }
\NormalTok{      text: [\textasciigrave{}X = [$\{x1.toFixed(1)\}, $\{x2.toFixed(1)\}]\textasciigrave{}], }
\NormalTok{      dy: 20,       // Posun nad bod}
\NormalTok{      dx: {-}5, }
\NormalTok{      fontWeight: "bold",}
\NormalTok{      fill: "steelblue"}
\NormalTok{    \}),}


\NormalTok{    // Legenda s výsledky// Legenda s výpočty}
\NormalTok{    Plot.text([\{x: 0.5, y: 10.0\}],\{x: "x", y: "y", text: [\textasciigrave{}||x||² = $\{Math.pow(norma\_x,2).toFixed(2)\}\textasciigrave{}],  textAnchor: "start", fill: "black" \}),}
\NormalTok{    Plot.text([\{x: 0.5, y: 9.5\}], \{x: "x", y: "y", text: [\textasciigrave{}x̄= $\{Math.pow(pru,1).toFixed(1)\}\textasciigrave{}],            textAnchor: "start", fill: "black" \}),}
\NormalTok{    Plot.text([\{x: 0.5, y: 9.0\}], \{x: "x", y: "y", text: [\textasciigrave{}||x̄||²= $\{Math.pow(pru,2).toFixed(2)\}\textasciigrave{}],       textAnchor: "start", fill: "black" \}),}
\NormalTok{    Plot.text([\{x: 0.5, y: 8.5\}], \{x: "x", y: "y", text: [\textasciigrave{}||e||² = $\{Math.pow(norma\_e,2).toFixed(2)\}\textasciigrave{}],  textAnchor: "start", fill: "black" \}),}
\NormalTok{    Plot.text([\{x: 0.5, y: 8.0\}], \{x: "x", y: "y", text: [\textasciigrave{}σ² = ||e||² / 2 = $\{(Math.pow(norma\_e, 2) / 2).toFixed(2)\}\textasciigrave{}], textAnchor: "start", fill: "black" \}),}
\NormalTok{    Plot.text([\{x: 0.5, y: 7.5\}], \{x: "x", y: "y", text: [\textasciigrave{}σ = √σ² = $\{sigma\_val.toFixed(2)\}\textasciigrave{}],           textAnchor: "start", fill: "black"\})}
\NormalTok{  ]}
\NormalTok{\})}
\end{Highlighting}
\end{Shaded}

Zaměřme se na velikosti těchto vektorů:Pomocí Pythogorovy věty dostáváme
rovnici

\[\|\mathbf{x}\|^2 = \|\mathbf{\bar{x}}\|^2 + \|\mathbf{e}\|^2\]

Víme, že \(\|\mathbf{x}\|=\sqrt{x_1^2+x_2^2}\), proto pro nás bude
výhodnější zůstat u druhých mocnin. Můžeme vyjádřit \(\|\mathbf{e}\|^2\)
ve tvaru

\begin{align}
  \|\mathbf{e}\|^2 
    &=e_1^2+e_2^2\\
    &=(x_1-\bar{x})^2+(x_2-\bar{x})^2\\
    &=\sum (x_i-\bar{x})^2, 
  \end{align}

nebo ve tvaru

\begin{align}
  \|\mathbf{e}\|^2
    &=\|\mathbf{x}\|^2 - \|\mathbf{\bar{x}}\|^2
    &=(x_1^2+x_2^2)-(\bar{x}^2+\bar{x}^2)
    &=\sum x_i^2-n\cdot\bar{x}^2 
  \end{align}

\(\|\mathbf{e}\|^2\) vyjadřuje celkovou odchylku dat od vekotru průměrů,
přepočítáno na jedno pozorování dostaneme veličiny zvanou
\textbf{rozptyl} a označovanou jako \(var\) nebo \(\sigma^2\) a která se
tedy rovná \(= \frac{\|\mathbf{e}\|^2}{n}\)

\[
  var=\sigma^2 = \frac{\sum (x_i - \bar{x})^2}{n}
\]

Protože tohle je velikost na druhou, ``průměrnou odchylku'' získáme jako
odmocninu z rozptylu

\[
  sd=\sqrt{\frac{\sum (x_i - \bar{x})^2}{n}}
\]

\end{tcolorbox}

\section{Momentové míry tvaru
rozdělení}\label{momentovuxe9-muxedry-tvaru-rozdux11blenuxed}

Připomeňme, že základním nástrojem vizualizace kardinálních dat je
histogram. Tvarem rozdělení tedy rozumíme tvar histogramu - jestli je
symetrický, zešikmený , vodorovný, \(\dots\) .

\begin{Shaded}
\begin{Highlighting}[]
\NormalTok{// 1. Generování dat na základě výběru}
\NormalTok{data02 = \{}
\NormalTok{  let v;}
\NormalTok{  const n = 500;}
\NormalTok{  if (distType02 === "Normální") \{}
\NormalTok{    v = Array.from(\{length: n\}, d3.randomNormal(0, 1));}
\NormalTok{  \} else if (distType02 === "Zešikmené") \{}
\NormalTok{    // Log{-}normální rozdělení je přirozeně zešikmené}
\NormalTok{    v = Array.from(\{length: n\}, d3.randomLogNormal(0, 0.5));}
\NormalTok{  \} else if (distType02 === "Uniformní") \{}
\NormalTok{    v = Array.from(\{length: n\}, d3.randomUniform({-}2, 2));}
\NormalTok{  \}}
\NormalTok{  return v.map(d =\textgreater{} (\{val: d\}));}
\NormalTok{\}}

\NormalTok{// 2. Ovládací prvky}
\NormalTok{viewof distType02 = Inputs.select(}
\NormalTok{  ["Normální", "Zešikmené", "Uniformní"], }
\NormalTok{  \{label: "Typ rozdělení"\}}
\NormalTok{)}

\NormalTok{viewof binCount = Inputs.range([2, 200], \{value: 20, step: 1, label: "Počet bins"\})}

\NormalTok{// 3. Graf}
\NormalTok{Plot.plot(\{}
\NormalTok{  height: 400,}
\NormalTok{  grid: true,}
\NormalTok{  marks: [}
\NormalTok{    // Histogram (Density)}
\NormalTok{    Plot.rectY(data02, Plot.binX(\{y: "count"\}, \{x: "val", thresholds: binCount, fill: "steelblue", fillOpacity: 0.7\})),}
    
\NormalTok{  // dodgeY vytvoří "mrak" bodů, anchor: "bottom" je drží u spodní hrany}
\NormalTok{    Plot.dot(data02, Plot.dodgeY(\{}
\NormalTok{      x: "val", }
\NormalTok{      anchor: "bottom", }
\NormalTok{      padding: 0.5, // Tímto zvětšuješ/zmenšuješ rozestupy mezi body}
\NormalTok{      r: 2, }
\NormalTok{      fill: "black",}
\NormalTok{      fillOpacity: 0.4}
\NormalTok{    \})),}
    
\NormalTok{    Plot.ruleY([0])}
\NormalTok{  ],}
\NormalTok{  y: \{}
\NormalTok{    label: "Četnost / Hustota",}
\NormalTok{    tickFormat: "d"}
\NormalTok{  \},}
\NormalTok{  x: \{}
\NormalTok{    label: "Hodnota x"}
\NormalTok{  \}}
\NormalTok{\})}
\end{Highlighting}
\end{Shaded}

\subsection{šikmost}\label{ux161ikmost}

Jak číselně odlišit data, tkerá jsou symetrická od dat, která symetrická
nejsou? Pokud využijeme odchylky, pak musíme \textbf{zdůraznit větší
odchylky} a \textbf{potlačit malé odchylky}, a zároveň \textbf{zachovat
znaménko} odchylky. Abychom získali míru která je nezávislá na poloze a
variabilitě konkrétních dat, musíme data nejprve
\textbf{standardizovat}, tedy převést na

\[z_i=\frac{(x_i-\bar{x})}{sd}\]

Proměnná \(z\) má průměr nula a směrodatnou odchylku jedna nezávisle na
tom, z jakých dat pochází (viz kapitola o aritmetice dat). hodnota
\(z_i\) naám tak přímo ukazuje standardizovanou odchylku dat od průměru.
Pro zjištění zešikmení využijeme třetí mocninu - ta potlačuje hodnoty
menší než \(|x_i|<1\) ,zvýrazňuje hodnoty \(|x_i|>1\) a přitom zachovává
stejné znaménko, takže budeme umět rozlišit, zda jsou data zešikmená
napravo nebo nalevo. \textbf{Šikmost} (\emph{skewness, skew}) tedy
počítáme pomocí vzorce

\[ 
  skew(X)=\frac{\sum z_i^3}{n}
\]

Když je skew záporné znamená to, že větší odchylky jsou vlevo a datsa
jsou nahuštěná vpravo, když je šikmost kladná, znamená to že větší
odchylky jsou doprava a data jsou nahuštěná vlevo. skew kolem hodnoty 0
znamená symetrické rozdělení

\begin{figure}

\centering{

\begin{Shaded}
\begin{Highlighting}[]
\NormalTok{//| echo: false}
\NormalTok{// 1. Ovládací prvek}
\NormalTok{viewof distType03 = Inputs.select(}
\NormalTok{  ["Normální", "Doprava (kladně) zešikmené", "Doleva (záporně) zešikmené", "Uniformní"], }
\NormalTok{  \{label: "Typ rozdělení"\}}
\NormalTok{)}

\NormalTok{// 2. Generování dat}
\NormalTok{data03 = \{}
\NormalTok{  const n = 512;}
\NormalTok{  let v;}
\NormalTok{  if (distType03 === "Normální") \{}
\NormalTok{    v = Array.from(\{length: n\}, d3.randomNormal(0,1));}
\NormalTok{  \} else if (distType03 === "Doprava (kladně) zešikmené") \{}
\NormalTok{    v = Array.from(\{length: n\}, () =\textgreater{} d3.randomLogNormal(0, 0.8)());}
\NormalTok{  \} else if (distType03 === "Doleva (záporně) zešikmené") \{}
\NormalTok{    const gen = d3.randomLogNormal(0, 0.8);}
\NormalTok{    v = Array.from(\{length: n\}, () =\textgreater{} 4 {-} gen()); }
\NormalTok{  \} else if (distType03 === "Uniformní") \{}
\NormalTok{    v = Array.from(\{length: n\}, d3.randomUniform({-}2, 2));}
\NormalTok{  \}}
\NormalTok{  return v.map(d =\textgreater{} (\{val: d\}));}
\NormalTok{\}}

\NormalTok{// 3. Výpočet statistik (včetně Sturgese a momentů)}
\NormalTok{stats03 = \{}
\NormalTok{  const values = data03.map(d =\textgreater{} d.val);}
\NormalTok{  const n = values.length;}
\NormalTok{  const mean = d3.mean(values);}
\NormalTok{  const sd = d3.deviation(values);}
\NormalTok{  const median = d3.median(values);}
\NormalTok{  const skew = d3.sum(values, v =\textgreater{} Math.pow((v {-} mean) / sd, 3)) / n;}
\NormalTok{  const kurt = (d3.sum(values, v =\textgreater{} Math.pow((v {-} mean) / sd, 4)) / n) {-} 3;}
\NormalTok{  const sturges = Math.ceil(Math.log2(n) + 1);}
  
\NormalTok{// Dynamický text porovnání}
\NormalTok{  let relation = "";}
\NormalTok{  if (distType03 === "Normální" || distType03 === "Uniformní") relation = "průměr ≈ medián";}
\NormalTok{  else if (distType03 === "Doprava (kladně) zešikmené") relation = "průměr \textgreater{} medián";}
\NormalTok{  else if (distType03 === "Doleva (záporně) zešikmené") relation = "průměr \textless{} medián";  }

\NormalTok{  return \{ mean, median, skew, kurt, sturges, relation \};  }
\NormalTok{\}}

\NormalTok{// 4. Graf s integrovanou legendou v Plot.text}
\NormalTok{Plot.plot(\{}
\NormalTok{  height: 450,}
\NormalTok{  grid: true,}
\NormalTok{  marks: [}
\NormalTok{    // Histogram}
\NormalTok{    Plot.rectY(data03, Plot.binX(\{y: "count"\}, \{}
\NormalTok{      x: "val", }
\NormalTok{      thresholds: stats03.sturges, }
\NormalTok{      fill: "\#f0f0f0", }
\NormalTok{      stroke: "\#ccc"}
\NormalTok{    \})),}
    
\NormalTok{    // Svislé čáry}
\NormalTok{    Plot.ruleX([stats03.mean], \{stroke: "steelblue", strokeWidth: 2\}),}
\NormalTok{    Plot.ruleX([stats03.median], \{stroke: "darkred", strokeWidth: 2\}),}

\NormalTok{    // Dynamická legenda pomocí Plot.text}
\NormalTok{    // Souřadnice x a y je nutné ladit podle rozsahu tvých dat (zde x: vlevo, y: nahoře)}
\NormalTok{    Plot.text([\textasciigrave{}Průměr: $\{stats03.mean.toFixed(3)\}\textasciigrave{}], \{frameAnchor: "top{-}left", dx: 10, dy: 20, fill: "steelblue", fontWeight: "bold", fontSize: 14\}),}
\NormalTok{    Plot.text([\textasciigrave{}Medián: $\{stats03.median.toFixed(3)\}\textasciigrave{}], \{frameAnchor: "top{-}left", dx: 10, dy: 40, fill: "darkred", fontWeight: "bold", fontSize: 14\}),}
\NormalTok{    Plot.text([\textasciigrave{}Šikmost (z³): $\{stats03.skew.toFixed(3)\}\textasciigrave{}], \{frameAnchor: "top{-}left", dx: 10, dy: 60, fontWeight: "bold", fontSize: 14\}),}
\NormalTok{    Plot.text([stats03.relation], \{frameAnchor: "top{-}left", dx: 10, dy: 80, fontSize: 14, fontWeight: "bold"\}),    }

\NormalTok{    // Mrak bodů}
\NormalTok{    Plot.dot(data03, Plot.dodgeY(\{}
\NormalTok{      x: "val", }
\NormalTok{      anchor: "bottom", }
\NormalTok{      padding: 0.5, }
\NormalTok{      r: 2, }
\NormalTok{      fill: "black",}
\NormalTok{      fillOpacity: 0.2}
\NormalTok{    \})),}
    
\NormalTok{    Plot.ruleY([0])}
\NormalTok{  ],}
\NormalTok{  y: \{ label: "Četnost" \},}
\NormalTok{  x: \{ label: "Hodnota (x)" \}}
\NormalTok{\})}
\end{Highlighting}
\end{Shaded}

}

\caption{\label{fig-skewness}Vztah mezi šikmostí průměrem, mediánem a
tvarem distribuce.}

\end{figure}%

\subsection{Špičatost}\label{ux161piux10datost}

Jak číselně odlišit data, která mají většinu hodnot blízko průměru od
hodnot, která mají rozprostření dat volnější a vzdálenější hodnoty
nejsou vyjímkou (\emph{heavy tails, těžké ocasy})? Protože zešikmení
máme popsané pomocí skew, nezáleží nám nyní na směru ochylky, musíme
však velmi zvýraznit odlehlé hodnoty. K tomu nám poslouží čtvrtá mocnina
\(z\), pomocí které určujeme \textbf{špičatost} (\emph{kurtoisis,
kurt}), kterou určíme pomocí vzorce

\[
  kurt(X) = \frac{\sum z_i^4}{n}
\]

Protože se jako základní distribuce bere \textbf{normální, gaussova,
zvonová} která má \(kurt=3\), uvádí se často vzorec pro \emph{excess
kurtosis} \(kurt_{excess}=\frac{\sum z_i^4}{n}-3\) který pak lze
intepretovat jako odchýlení špičatosti od normálního rozdělení.
Podrobněji o normálním rozdělení viz samostatná kapitola.

Následující graf ukazuje vliv velikosti odchylek na jejich ``příspěvek''
do průměru, rozptylu, šikmosti a špičatosti. Na ose x je hodnota \(z\),
na jednotlivých křivkách velikosti \(|z|,z^2,z^3,z^4\), které tvoří
příspěvek do příslušných momentových metrik.

Pro \(k={2,3,4}\) platí: \(|z_i|<1\Longrightarrow |z_i|^k<|z_i|\), takže
jejich vliv je potlačen. Pro hodnoty
\(|z_i|>1\Longrightarrow |z_i|^k > |z_i|\), takže jejich vliv je výrazně
posílen.

\begin{figure}

\centering{

\begin{Shaded}
\begin{Highlighting}[]
\NormalTok{//| echo: false}
\NormalTok{viewof z = Inputs.range([{-}2, 2], \{label: "Odchylka (z{-}skóre)", step: 0.05, value: 1.5\})}

\NormalTok{Plot.plot(\{}
\NormalTok{  height: 800,}
\NormalTok{  y: \{domain: [{-}8, 8], grid: true, label: "Váha v momentu (f(z))"\},}
\NormalTok{  x: \{domain: [{-}2, 2], grid: true, label: "Vzdálenost od průměru (z)"\},}
\NormalTok{  style: \{fontSize: "14px"\},}
\NormalTok{  marks: [}
\NormalTok{    // Vodorovná osa (nula) {-} opraveno na ruleY}
\NormalTok{    Plot.ruleY([0], \{strokeWidth: 1.5\}),}
    
\NormalTok{    // Svislé pomocné čáry na jednotkových odchylkách}
\NormalTok{    Plot.ruleX([{-}1, 1], \{strokeOpacity: 0.2, strokeDasharray: "4,4"\}),}

\NormalTok{    // MAD |z| {-} Lineární}
\NormalTok{    Plot.line(d3.range({-}2, 2.05, 0.05), \{x: d =\textgreater{} d, y: d =\textgreater{} Math.abs(d), stroke: "steelblue", dash: "none", strokeWidth: 2, label: "Průměr (|z|¹)"\}),}
\NormalTok{    // Rozptyl z\^{}2}
\NormalTok{    Plot.line(d3.range({-}2, 2.05, 0.05), \{x: d =\textgreater{} d, y: d =\textgreater{} d**2, stroke: "\#7A5B8C", dash: "8,4", strokeWidth: 2, label: "Rozptyl (z²)"\}),}
\NormalTok{    // Šikmost z\^{}3}
\NormalTok{    Plot.line(d3.range({-}2, 2.05, 0.05), \{x: d =\textgreater{} d, y: d =\textgreater{} d**3, stroke: "\#AF3363", dash: "4,3", strokeWidth: 2, label: "Šikmost (z³)"\}),}
\NormalTok{    // Špičatost z\^{}4}
\NormalTok{    Plot.line(d3.range({-}2, 2.05, 0.05), \{x: d =\textgreater{} d, y: d =\textgreater{} d**4, stroke: "darkred", dash: "1,2", strokeWidth: 2,  label: "Špičatost (z⁴)"\}),}

\NormalTok{    // Interaktivní tečky (body na křivkách podle posuvníku)}
\NormalTok{    Plot.dot([z], \{x: d =\textgreater{} d, y: d =\textgreater{} Math.abs(d), fill: "steelblue", r: 4\}),}
\NormalTok{    Plot.dot([z], \{x: d =\textgreater{} d, y: d =\textgreater{} d**2, fill: "\#7A5B8C", r: 4\}),}
\NormalTok{    Plot.dot([z], \{x: d =\textgreater{} d, y: d =\textgreater{} d**3, fill: "\#AF3363", r: 4\}),}
\NormalTok{    Plot.dot([z], \{x: d =\textgreater{} d, y: d =\textgreater{} d**4, fill: "darkred", r: 6\}),}

\NormalTok{    // Dynamická legenda v rohu}
\NormalTok{    Plot.text([\textasciigrave{}pro průměr: |z|= $\{Math.abs(z).toFixed(2)\}\textasciigrave{}], \{x: {-}1.9, y: 7.5, fill: "steelblue", textAnchor: "start"\}),}
\NormalTok{    Plot.text([\textasciigrave{}pro rozptyl: z²= $\{(z**2).toFixed(2)\}\textasciigrave{}], \{x: {-}1.9, y: 6.5, fill: "\#7A5B8C", textAnchor: "start"\}),}
\NormalTok{    Plot.text([\textasciigrave{}pro šikmost: z³= $\{(z**3).toFixed(2)\}\textasciigrave{}], \{x: {-}1.9, y: 5.5, fill: "\#AF3363", textAnchor: "start"\}),}
\NormalTok{    Plot.text([\textasciigrave{}pro špičatost: z⁴= $\{(z**4).toFixed(2)\}\textasciigrave{}], \{x: {-}1.9, y: 4.5, fill: "darkred", textAnchor: "start"\})}
\NormalTok{  ]}
\NormalTok{\})}
\end{Highlighting}
\end{Shaded}

}

\caption{\label{fig-skewness}Vliv vzdálenosti od průměru na velikost
započítávané hodnoty do momentových charakteristik mad, sd, skew a kurt.
Osa \(x\) obsahuje standardizovaná data \(z_i\), osa y transformovaná
data \(z_i^k\) pro všechny čtyři momentové charakteristiky.}

\end{figure}%

A jak to konkrétně vypadá pro různě rozprostřená data ukazuje
následující graf. Pro \(k={2,3,4}\) platí:
\(|z_i|<1\Longrightarrow |z_i|^k<|z_i|\), takže jejich vliv je potlačen.
V následujícím grafu jsou to data v modře podsvícené oblasti. Pro
hodnoty \(|z_i|>1\Longrightarrow |z_i|^k > |z_i|\), takže jejich vliv je
výrazně posílen.

\begin{figure}

\centering{

\begin{Shaded}
\begin{Highlighting}[]
\NormalTok{//| echo: false}
\NormalTok{// 1. Vstupy}
\NormalTok{viewof moment = Inputs.select(}
\NormalTok{  new Map([}
\NormalTok{    ["MAD (Průměrná odchylka) |z|", "abs"],}
\NormalTok{    ["Rozptyl (Kvadratická) z²", "var"],}
\NormalTok{    ["Šikmost (Kubická) z³", "skew"],}
\NormalTok{    ["Špičatost (Kvadraticá) z⁴", "kurt"]}
\NormalTok{  ]), }
\NormalTok{  \{label: "Zvolte statistický moment:"\}}
\NormalTok{)}

\NormalTok{viewof distType = Inputs.radio(}
\NormalTok{  ["Normální", "Těžké ocasy", "Šikmé vpravo", "Šikmé vlevo"], }
\NormalTok{  \{value: "Normální", label: "Rozdělení dat:"\}}
\NormalTok{)}

\NormalTok{// Slovník pro popisky výsledku}
\NormalTok{momentLabels = (\{}
\NormalTok{  abs: "Průměrná odchylka (MAD)",}
\NormalTok{  var: "Rozptyl (Var)",}
\NormalTok{  skew: "Šikmost (Skewness)",}
\NormalTok{  kurt: "Špičatost (Kurtosis)"}
\NormalTok{\})}

\NormalTok{// 2. Dynamický rozsah osy Y podle zvoleného momentu}
\NormalTok{yLimit = \{}
\NormalTok{  if (moment === "abs") return 2.5;}
\NormalTok{  if (moment === "var") return 4.5;}
\NormalTok{  if (moment === "skew") return 8.5;}
\NormalTok{  return 16.5; // pro z\^{}3 a z\^{}4}
\NormalTok{\}}



\NormalTok{// 3. Definice dat (beze změny)}
\NormalTok{dataMap = new Map([}
\NormalTok{  ["Normální", [{-}1.5, {-}0.6, {-}0.3, {-}0.2, {-}0.1, 0.1, 0.2, 0.3, 0.6, 1.5]],}
\NormalTok{  ["Těžké ocasy", [{-}2.0, {-}0.5, {-}0.3, {-}0.1, 0, 0, 0.1, 0.3, 0.5, 2.0]],}
\NormalTok{  ["Šikmé vpravo", [{-}0.8, {-}0.7, {-}0.6, {-}0.5, {-}0.4, {-}0.3, {-}0.2, 0.5, 1.2, 1.8]],}
\NormalTok{  ["Šikmé vlevo", [{-}1.8, {-}1.2, {-}0.5, 0.2, 0.3, 0.4, 0.5, 0.6, 0.7, 0.8]]  }
\NormalTok{])}


\NormalTok{currentData = dataMap.get(distType)}

\NormalTok{functionTransform = (z) =\textgreater{} \{}
\NormalTok{  if (moment === "abs") return Math.abs(z);}
\NormalTok{  if (moment === "var") return z**2;}
\NormalTok{  if (moment === "skew") return z**3;}
\NormalTok{  if (moment === "kurt") return z**4;}
\NormalTok{\}}

\NormalTok{results = currentData.map(functionTransform)}
\NormalTok{finalValue = d3.mean(results)}

\NormalTok{// 4. Graf s responzivním y domain}
\NormalTok{Plot.plot(\{}
\NormalTok{  height: 600,}
\NormalTok{  width: width,}
\NormalTok{  y: \{}
\NormalTok{    domain: moment === "skew" ? [{-}yLimit, yLimit] : [moment === "abs" ? {-}1.5 : {-}1, yLimit], }
\NormalTok{    grid: true, }
\NormalTok{    label: "Transformovaná hodnota (y)"}
\NormalTok{  \},}

\NormalTok{  x: \{domain: [{-}3, 3], grid: true, label: "Původní z{-}skóre (x)"\},}
\NormalTok{  marks: [}
\NormalTok{    // KRABICE TLUMIVOSTI: x i y v rozsahu [{-}1, 1]}
\NormalTok{    Plot.rect([{-}1], \{}
\NormalTok{      x1:{-}1, x2: 1, y1: {-}1, y2: 1, }
\NormalTok{      fill: "steelblue", fillOpacity: 0.15,}
\NormalTok{      stroke: "steelblue",}
\NormalTok{      strokeDasharray: "2,2",}
\NormalTok{      strokeOpacity: 0.3}
\NormalTok{    \}),}
\NormalTok{    Plot.frame(\{stroke: "\#ccc", strokeDasharray: "4,4"\}),}
    
\NormalTok{    // hlavní osy }
\NormalTok{    Plot.ruleY([0], \{strokeOpacity: 0.8\}),}
\NormalTok{    Plot.ruleX([0], \{strokeOpacity: 0.5\}),}

\NormalTok{    // Svislé referenční čáry na bodech zlomu}
\NormalTok{    Plot.ruleX([{-}1, 1], \{stroke: "\#666", strokeDasharray: "4,4", strokeWidth: 1\}),}
    
\NormalTok{    // Vodorovná referenční čára y=1 (bod, kde se všechny křivky potkají)}
\NormalTok{    Plot.ruleY([1], \{stroke: "\#666", strokeDasharray: "4,4", strokeWidth: 1, strokeOpacity: 1\}),}
\NormalTok{    Plot.ruleY([{-}1], \{stroke: "\#666", strokeDasharray: "4,4", strokeWidth: 1, strokeOpacity: 1\}),}

\NormalTok{    // Křivka funkce}
\NormalTok{    Plot.line(d3.range({-}3, 3.1, 0.1), \{}
\NormalTok{      x: d =\textgreater{} d, }
\NormalTok{      y: d =\textgreater{} functionTransform(d), }
\NormalTok{      stroke: "\#ccc", strokeWidth: 4}
\NormalTok{    \}),}

\NormalTok{    // Projekce X {-}\textgreater{} Křivka {-}\textgreater{} Y}
\NormalTok{    Plot.ruleX(currentData, \{x: d =\textgreater{} d, y1: 0, y2: d =\textgreater{} functionTransform(d), stroke: "steelblue", strokeOpacity: 0.5\}),}
\NormalTok{    Plot.ruleY(results, \{y: d =\textgreater{} d, x1: (d, i) =\textgreater{} currentData[i], x2: 0, stroke: "steelblue", strokeDasharray: "2,2", strokeOpacity: 0.5\}),}

\NormalTok{    // Body}
\NormalTok{    Plot.dot(currentData, \{x: d =\textgreater{} d, y: 0, fill: "steelblue", r: 6\}), // Data na X}
\NormalTok{    Plot.dot(currentData, \{x: d =\textgreater{} d, y: d =\textgreater{} functionTransform(d), fill: "steelblue", r: 3\}), // Na křivce}
\NormalTok{    Plot.dot(results, \{x: 0, y: d =\textgreater{} d, fill: "steelblue", r: 4\}), // Na ose Y}

\NormalTok{    // Modrý diamant {-} Průměr (Výsledný moment)}
\NormalTok{    Plot.dot([finalValue], \{x: 0, y: d =\textgreater{} d, fill: "darkred", symbol: "circle", r: 10, stroke: "white", strokeWidth: 2\}),}
    
\NormalTok{    // Dynamický popisek výsledku}
\NormalTok{    Plot.text([finalValue], \{}
\NormalTok{      x: 0.2, y: d =\textgreater{} d, }
\NormalTok{      text: d =\textgreater{} \textasciigrave{} $\{momentLabels[moment]\} = $\{d.toFixed(3)\}\textasciigrave{},}
\NormalTok{      textAnchor: "start", fill: "darkred", fontWeight: "bold", fontSize: 16}
\NormalTok{    \})}
\NormalTok{  ]}
\NormalTok{\})}
\end{Highlighting}
\end{Shaded}

}

\caption{\label{fig-moments}Momentové charakteristiky různých rozdělení.
Osa \(x\) obsahuje standardizovaná data \(z_i\), osa y transformovaná
data \(z_i^k\) a jejich průměr. Modrá oblast je oblast útlumu, data z
této oblasti mají potlačený vliv a data mimo naopak posílený.}

\end{figure}%

\section*{Klíčové poznatky}\label{kluxedux10dovuxe9-poznatky-1}
\addcontentsline{toc}{section}{Klíčové poznatky}

\markright{Klíčové poznatky}

{\def\LTcaptype{none} % do not increment counter
\begin{longtable}[]{@{}
  >{\raggedright\arraybackslash}p{(\linewidth - 2\tabcolsep) * \real{0.5000}}
  >{\raggedright\arraybackslash}p{(\linewidth - 2\tabcolsep) * \real{0.5000}}@{}}
\toprule\noalign{}
\endhead
\bottomrule\noalign{}
\endlastfoot
průměr & \(\bar{x}=\frac{\sum x_i}{n}\) \\
směrodatná odchylka &
\(sd=\sqrt{\frac{1}{n}\sum (x_i - \bar{x})^2}=\sqrt{var}\) \\
rozptyl & \(var=\frac{1}{n}\sum (x_i - \bar{x})^2=sd^2\) \\
standardizace & \(z_i=\frac{x_i - \bar{x}}{sd}\) \\
šikmost & \(skew=\frac{1}{n}\sum z_i^3\) \\
špičatost & \(kurt=\frac{1}{n}\sum z_i^4\), nebo
\(kurt_{excess}=\frac{1}{n}\sum z_i^4-3\) \\
\end{longtable}
}

\chapter{Vlastnosti metrik a aritmetika
dat}\label{vlastnosti-metrik-a-aritmetika-dat}

V této kapitole si ukážeme, jak reagují jednotlivé statistiky na to když

\begin{enumerate}
\def\labelenumi{\arabic{enumi}.}
\tightlist
\item
  k datům přičteme stejnou hodnotu
\item
  když data vynásobíme nějakým číslem
\item
  když přidáme odlehlou hodnotu
\end{enumerate}

\section{Aritmetika dat}\label{aritmetika-dat}

Zůstaňme u příkladu z předchozí lekce: máme deset domácností a
sledujeme, jaký mají počet motorových vozidel (automobilů, motocyklů,
elektrokoloběžek). Získaná data po uspořádání od nejmenšího po největší
jsou \(X=\{1, 2, 2, 2, 3, 3, 4, 5, 8, 11\}\). O nich víme, že je lze
popsat následujícími charakteristikami (zaokrouhlenými na dvě desetinná
místa):

{\def\LTcaptype{none} % do not increment counter
\begin{longtable}[]{@{}ll@{}}
\toprule\noalign{}
statistika & \(X\) \\
\midrule\noalign{}
\endhead
\bottomrule\noalign{}
\endlastfoot
\(\bar{x}=m\) & 4.10 \\
\(mds\) & 2.00 \\
\(mdn\) & 3.00 \\
\((min, max)\) & (1.00, 11.00) \\
\((1Q, 3Q)\) & (2.00, 4.75) \\
\((LF, UF)\) & (-2.13, 8.88) \\
\(range\) & 10.00 \\
\(IQR\) & 2.75 \\
\(var, sd\) & 8.89, 2.98 \\
\end{longtable}
}

Zkusme vytvořit nové soubory dat tak, že nejprve přičteme hodnotu 2 ke
každému číslu (\(X+2\)) a nebo ho vynásobíme (\(X\cdot 2\))

{\def\LTcaptype{none} % do not increment counter
\begin{longtable}[]{@{}llll@{}}
\toprule\noalign{}
statistika & \(X\) & \(X+2\) & \(X\cdot 2\) \\
\midrule\noalign{}
\endhead
\bottomrule\noalign{}
\endlastfoot
\(\bar{x}=m\) & 4.10 & 6.10 & 8.20 \\
\(mds\) & 2.00 & 4.00 & 4.00 \\
\(mdn\) & 3.00 & 5.00 & 6.00 \\
\((min, max)\) & (1.00, 11.00) & (3.00, 11.00) & (2.00, 22.00) \\
\((1Q, 3Q)\) & (2.00, 4.75) & (5.00, 6.75) & (4.00, 9.50) \\
\((LF, UF)\) & (-2.13, 8.88) & (-0.13, 10.33) & (-4.25, 17.75) \\
\(range\) & 10.00 & 10.00 & 20.00 \\
\(IQR\) & 2.75 & 2.75 & 5.50 \\
\(var, sd\) & 8.89, 2.98 & 8.89, 2.98 & 35.56, 5.96 \\
\end{longtable}
}

Můžete pozorovat, že při přičtení se míry polohy se všechny posunuly od
dva, zatímco míry variability \textbf{range, IQR, var, sd} zůstaly
stejné. Při násobení se míry polohy vynásobily dvěma, míry variability
\textbf{range, iqr, sd} se také vynásobily hodnotou \(a=2\), zatímco
\textbf{var} se vynásobila hodnotou \(a^2=4\).

\begin{tcolorbox}[enhanced jigsaw, arc=.35mm, bottomrule=.15mm, bottomtitle=1mm, breakable, colback=white, colbacktitle=quarto-callout-note-color!10!white, colframe=quarto-callout-note-color-frame, coltitle=black, left=2mm, leftrule=.75mm, opacityback=0, opacitybacktitle=0.6, rightrule=.15mm, title=\textcolor{quarto-callout-note-color}{\faInfo}\hspace{0.5em}{Poznámka}, titlerule=0mm, toprule=.15mm, toptitle=1mm]

Pokud manažer firmy všem přidá stejnou hodnotu (\(x+1000\)), všichni si
polepší a rozdíly zůstanou stejné. Pokud všem navýší plat o 10\%
(\(x\cdot 1.01)\), všichni si polepší ale také se zvětší rozdíly -
\emph{procentuální navýšení příjmů rozevírá nůžky mezi bohatšími a
chušími}.

Pokud stát požaduje od všech stejnou částku (např. televizní poplatky),
všem se sníží rozpočet stejně \(x-200\). Pokud však nastaví platby např
procentuálně podle příjmu \(x-0.01\cdot x\), budou vyskopříjmoví platit
více než nízkopříjmoví - \emph{procentuální snžení je trest za úspěch}.

\end{tcolorbox}

\section{Transformace dat}\label{transformace-dat}

Zkusme naše data \textbf{standardizovat} tak, abychom eliminovali posun
a zároveň upravili variabilitu. Cílem je, aby data s rozdílnou polohou a
rozdílnou variabilitou byla porovnateln (např počet vozidel v rodině a
celkový roční příjem rodiny). Jednou z tradičních standardizací je
převod hodnot \(x_i\) na \(z-skór\) pomocí vzorce

\[z_i=\frac{x_i-\bar{x}}{sd}\]

Z skór tedy získáme nejprve posunem dat o průměr a poté vydělením
směrodatnou odchylkou. Získané proměnná \(z\) pak má průměr nula a
směrodatnou odchylku 1. Tato standardizace je vhodná pro data alespoň
přibližně symetrická kolem průměru, jinými slovy kdy je průměr blízko
mediánu,\(\bar{x}-mdn\to 0\).

Variantou odolnou vůči odlehlým hodotám je \textbf{robustní
standardizace}m která každé hodnotě \(x_i\) přiřadí hodnotu

\[x_r=\frac{x_i-mdn}{iqr}\]

Robustní standarizaci tedy získáme posunem dat o medián a vydělením
mezikvartilovým rozpětím. Získaná proměnná \(x_r\) pak má medán nula a
iqr jedna. Rubustní standardizace je vhodná pro data zešikmená, kdy
většina hodnot je na jedné straně od průměru, jinými slovy medián a
průměr jsou daleko.

Jinou bežně používanou trasformací je \textbf{normalizace} rozsahu,
které každé hodnotě \(x\) přiřadí hodnotu

\[x_{n}=\frac{x-min}{range}\]

Hodnoty tedy nejprve posuneme o minimum a pak vydělíme celkovým
rozsahem. Takto Rozsah proměnné \(x_n\) bude \((min,max)=(0,1)\), tedy
\(range=1\). Tato transformace je vhodná a někdy i nutná pro algoritmy,
kde by hodnoty nad jedna mohly vrazně ovlivnit schopnost nalezení
výsledku.

Existují i další transformace, které jsou založené na rozdělení dat, ale
těmi se nyní zabývat nebudeme.

{\def\LTcaptype{none} % do not increment counter
\begin{longtable}[]{@{}
  >{\raggedright\arraybackslash}p{(\linewidth - 8\tabcolsep) * \real{0.2000}}
  >{\raggedright\arraybackslash}p{(\linewidth - 8\tabcolsep) * \real{0.2000}}
  >{\raggedright\arraybackslash}p{(\linewidth - 8\tabcolsep) * \real{0.2000}}
  >{\raggedright\arraybackslash}p{(\linewidth - 8\tabcolsep) * \real{0.2000}}
  >{\raggedright\arraybackslash}p{(\linewidth - 8\tabcolsep) * \real{0.2000}}@{}}
\toprule\noalign{}
\begin{minipage}[b]{\linewidth}\raggedright
statistika
\end{minipage} & \begin{minipage}[b]{\linewidth}\raggedright
\(X\)
\end{minipage} & \begin{minipage}[b]{\linewidth}\raggedright
\(z(X)\)
\end{minipage} & \begin{minipage}[b]{\linewidth}\raggedright
\(x_r\)
\end{minipage} & \begin{minipage}[b]{\linewidth}\raggedright
\(x_n\)
\end{minipage} \\
\midrule\noalign{}
\endhead
\bottomrule\noalign{}
\endlastfoot
\(\bar{x}=m\) & 4.10 & \textbf{0.00} & 0.40 & 0.31 \\
\(mds\) & 2.00 & -0.70 & -0.36 & 0.10 \\
\(mdn\) & 3.00 & -0.37 & \textbf{0.00} & 0.20 \\
\((min, max)\) & (1.00, 11.00) & (-1.04, 2.32) & (-0.73, 2.91) &
\textbf{(0.00, 1.00)} \\
\((1Q, 3Q)\) & (2.00, 4.75) & (-0.70, 2.20) & (-0.36, 0.64) & (0.10,
0.38) \\
\((LF, UF)\) & (-2.13, 8.88) & (-2.09, 1.60) & (-1.86, 2.14) & (-0.31,
0.79) \\
\(range\) & 10.00 & 3.35 & 3.64 & \textbf{1.00} \\
\(IQR\) & 2.75 & 0.92 & \textbf{1.00} & 0.28 \\
\(var, sd\) & 8.89, 2.98 & \textbf{1.00, 1.00} & 1.18, 1.08 & 0.09,
0.03 \\
\end{longtable}
}

\begin{Shaded}
\begin{Highlighting}[]
\NormalTok{//| echo: false}
\NormalTok{data = [1, 2, 2, 2, 3, 3, 4, 5, 8, 11].map(d =\textgreater{} (\{ hodnota: d \}))}

\NormalTok{// Statistiky pro čáry a legendu}
\NormalTok{stats = [}
\NormalTok{  \{ label: "Modus", val: 2, color: "green", dash: "2,2", yOffset: 15 \},}
\NormalTok{  \{ label: "Medián", val: 3, color: "blue", dash: "2.2", yOffset: 35 \},}
\NormalTok{  \{ label: "Průměr", val: 4.1, color: "red", dash: "2.2", yOffset: 55 \}}
\NormalTok{]}

\NormalTok{Plot.plot(\{}
\NormalTok{  width: 600,}
\NormalTok{  height: 200,}
\NormalTok{  grid: true,}
\NormalTok{  x: \{  domain: [0, 12], }
\NormalTok{        ticks: d3.range(1, 12), }
\NormalTok{        label: "Hodnota →" \},}
\NormalTok{  y: \{  axis: null, }
\NormalTok{        domain: [0, 3] \}, }
\NormalTok{  marks: [}
\NormalTok{    // Tečky nad sebou}
\NormalTok{    Plot.dot(data, Plot.stackY(\{ x: "hodnota", fill: "gray", stroke: "white", r: 6 \})),}

\NormalTok{    // Svislé čáry na pozicích statistik}
\NormalTok{    stats.map(s =\textgreater{} Plot.ruleX([s.val], \{ stroke: s.color, strokeWidth: 1, strokeDasharray: s.dash \})),}

\NormalTok{    // Popisky v pravém horním rohu pod sebou}
\NormalTok{    stats.map(s =\textgreater{} }
\NormalTok{      Plot.text([0], \{ // [0] je zde jen zástupný, pozici určí frameAnchor}
\NormalTok{        text: [\textasciigrave{}$\{s.label\}: $\{s.val\}\textasciigrave{}],}
\NormalTok{        frameAnchor: "top{-}right",}
\NormalTok{        dx: {-}20,        // odsazení od pravého okraje}
\NormalTok{        dy: s.yOffset,  // vertikální rozestup mezi popisky}
\NormalTok{        fill: s.color,}
\NormalTok{        // fontWeight: "bold",}
\NormalTok{        fontSize: 14,}
\NormalTok{        textAnchor: "end"}
\NormalTok{      \})}
\NormalTok{    )}
\NormalTok{  ]}
\NormalTok{\})}
\end{Highlighting}
\end{Shaded}

\chapter{Deskriptivní statistiky z tabulky
dat}\label{deskriptivnuxed-statistiky-z-tabulky-dat}

Představte si, že jsme provedli průzkum spokojenosti zákazníků s novou
službou, kde zákazníci hodnotili spokojenost na škále od 1 (velmi
nespokojený) do 10 (velmi spokojený). Zde je souhrn dat z odpovědí 200
náhodně vybraných zákazníků:

\begin{Shaded}
\begin{Highlighting}[]
\NormalTok{//| echo: false}
\NormalTok{// Definice dat: 200 uživatelů, zašikmeno k vyšším hodnotám (1{-}10)}
\NormalTok{data\_base = [}
\NormalTok{  \{x: 1,  n: 2\},}
\NormalTok{  \{x: 2,  n: 3\},}
\NormalTok{  \{x: 3,  n: 5\},}
\NormalTok{  \{x: 4,  n: 10\},}
\NormalTok{  \{x: 5,  n: 15\},}
\NormalTok{  \{x: 6,  n: 25\},}
\NormalTok{  \{x: 7,  n: 30\},}
\NormalTok{  \{x: 8,  n: 40\},}
\NormalTok{  \{x: 9,  n: 45\},}
\NormalTok{  \{x: 10, n: 25\}}
\NormalTok{]}
\end{Highlighting}
\end{Shaded}

\begin{Shaded}
\begin{Highlighting}[]
\NormalTok{//| echo: false}
\NormalTok{//| label: tbl{-}spokojenost{-}kumulativni}
\NormalTok{//| tbl{-}cap: "Tabulka četností s kumulativní relativní četností (n=200)"}

\NormalTok{// Výpočet relativních a kumulativních četností}
\NormalTok{table\_final = \{}
\NormalTok{  let acc = 0;}
\NormalTok{  return data\_base.map(d =\textgreater{} \{}
\NormalTok{    const f = d.n / 200;}
\NormalTok{    acc += f;}
\NormalTok{    return \{}
\NormalTok{      x: d.x,}
\NormalTok{      n: d.n,}
\NormalTok{      f: f.toFixed(3),}
\NormalTok{      F: acc.toFixed(3)}
\NormalTok{    \};}
\NormalTok{  \});}
\NormalTok{\}}

\NormalTok{// Vykreslení tabulky}
\NormalTok{Inputs.table(table\_final, \{}
\NormalTok{  columns: ["x", "n", "f", "F"],}
\NormalTok{  header: \{}
\NormalTok{    x: "Hodnocení (xᵢ)",}
\NormalTok{    n: "Četnost (nᵢ)",}
\NormalTok{    f: "Relat. četnost (pᵢ)",}
\NormalTok{    F: "Kumul. relat. četnost (Fᵢ)"}
\NormalTok{  \}}
\NormalTok{\})}
\end{Highlighting}
\end{Shaded}

Pojďme tato data zobrazit graficky v podobně histogramu, tedy diagramu
četností

\begin{Shaded}
\begin{Highlighting}[]
\NormalTok{//| echo: false}

\NormalTok{// Histogram}
\NormalTok{Plot.plot(\{}
\NormalTok{  height: 250,}
\NormalTok{  //x: \{label: "Hodnocení (xᵢ)", domain: [0.5, 10.5], ticks: d3.range(1, 11)\},}
\NormalTok{  //y: \{label: "Četnost (nᵢ)", grid: true\},}
\NormalTok{  //marks: [}
\NormalTok{  //  Plot.barY(data\_base, \{x: "x", y: "n", fill: "\#0275d8", inset: 0.5\}),}
\NormalTok{  //  Plot.ruleY([0])}
\NormalTok{  //]}
\NormalTok{  marks: [}
\NormalTok{    Plot.barY(data\_base, \{}
\NormalTok{      x: "x", }
\NormalTok{      y: "n", }
\NormalTok{      fill: "steelblue",}
\NormalTok{      inset: 0.5 }
\NormalTok{    \}),}
\NormalTok{    Plot.ruleY([0])}
\NormalTok{  ]  }
\NormalTok{\})}
\end{Highlighting}
\end{Shaded}

\section{Ukazatele polohy: kvantily}\label{ukazatele-polohy-kvantily}

\textbf{Kvantil} je \textbf{hodnota, pod kterou leží určité procento
dat}. Důležité kvantily mají svá spěcifická označení: \textbf{medián}
pro polovinu, \textbf{kvartily} pro čtvrtiny, \textbf{decily} pro
desetiny a \textbf{percentily} pro setiny. Proto platí že např medián =
druhý kvartil = pátý deci = padesátý percentil.

\textbf{Min} je nejmenší, \textbf{max} největší hodnota v datech.

Pokud budeme hledat kvantily v tabulce, zaměříme se na sloupec
\emph{Kumulativní (Fᵢ)} ve kterém hledáme, kdy poprvé dosáhneme dané
hodnoty.

\begin{enumerate}
\def\labelenumi{\arabic{enumi}.}
\tightlist
\item
  \textbf{Dolní kvartil (\(Q_1\)):} hledáme hodnotu \textbf{0,25}. V
  tabulce vidíme, že pro \(x=5\) je \(F(5)=0,175\) menší než \(0.5\),
  ale \(F(6)=0,300\). Takže \(Q_1 = 6\).
\item
  \textbf{Medián (\(Q_2\)):} stručněji: \(F(7)= 0,450, F(8)=0.650\)
  Protože \(F(7) < 0.5 < F(8)\), je \(Q2=8\).
\item
  \textbf{Horní kvartil (\(Q_3\)):}
  \(F(8)=0.650, F(9)=0.875 \longrightarrow Q_3 = 9\).
\item
  \textbf{\(Q_3\) pomocí interpolací:}
  \(k=1 + (10-1)\cdot 0.75 = 7.75, F(8)<7.75<F(9)\longrightarrow Q_3= 8+(8-7) \cdot 0.75=8.75\)
\end{enumerate}

Různé způsoby výpočtu percentilů mohou vést k různým, ale podobným
výsledkům.

\section{Kvantilové ukazatele variability: jak moc se zákazníci
(ne)shodují?}\label{kvantilovuxe9-ukazatele-variability-jak-moc-se-zuxe1kaznuxedci-neshodujuxed}

Když známe tyto body, můžeme popsat, jak moc jsou data ``rozptýlená'':

\begin{enumerate}
\def\labelenumi{\arabic{enumi}.}
\setcounter{enumi}{3}
\tightlist
\item
  \textbf{Rozpětí} (\emph{range}): rozdíl mezi nejlepším a nejhorším
  \(max-min\), v našich datech \(range = 10 - 1 = 9\).
\item
  \textbf{Mezikvartilové rozpětí} (\emph{inter-quartile-range, IQR}):
  rozdíl mezi třemi čtvrtinami a jednou čtvrtinou, \(Q_3-Q_1\), tedy
  \(IQR = 9 - 6 = 3\).
\end{enumerate}

\begin{tcolorbox}[enhanced jigsaw, arc=.35mm, bottomrule=.15mm, bottomtitle=1mm, breakable, colback=white, colbacktitle=quarto-callout-note-color!10!white, colframe=quarto-callout-note-color-frame, coltitle=black, left=2mm, leftrule=.75mm, opacityback=0, opacitybacktitle=0.6, rightrule=.15mm, title=\textcolor{quarto-callout-note-color}{\faInfo}\hspace{0.5em}{Poznámka}, titlerule=0mm, toprule=.15mm, toptitle=1mm]

\textbf{Range} ani min,max se nemusí shodovat z rozsahem škály.
Například škála 1-10 může mít v datech min=5 a max=7.

\textbf{IQR} nám říká, v jakém rozmezí se pohybují hodnoty ``střední
poloviny'', tedy když ignorujeme čtvrtinu nejmenších a čtvrtinu
největších hodnot.

\end{tcolorbox}

\section{Aritmetický průměr a těžiště
dat}\label{aritmetickuxfd-prux16fmux11br-a-tux11bux17eiux161tux11b-dat}

Nejznámější hodnotou popisující soubor dat je bezesporu průměr, přesněji
artimetický průměr. Ten lze vypočít tak, že sečteme všechny hodnoty a
vydělíme je počtem hodnot, \(\bar{x} = \frac{\sum (x_i)}{N}\). Pokud
máme data v tabulce, pak tam některé hodnoty máme vícekrát, např u nás

\(\bar{x} = \frac{1+1+2+2+2+3+\dots}{200}\), což můžeme lehce upravit
pomocí součinu na

\(\bar{x} = \frac{1\cdot 2+2\cdot 3+3\cdot 5+\dots}{200}\), velký zlomek
ještě rozdělit na samostatné zlomky

\(\bar{x} = \frac{1\cdot 2}{200}+\frac{2\cdot 3}{200}+\frac{3\cdot 5}{200}+\dots\)
a na závěr ještě upravit na relativní četnosti

\(\bar{x} = 1\cdot \frac{2}{200}+2\cdot \frac{3}{200}+3\cdot\frac{5}{200}+\dots\)

Získali jsme tak vzorec který je uváděný jako \textbf{vážený průměr},
kde každou hodnotu ``vážíme'' její četností \(n_i\) nebo relativní
četností \(p_i\):

\[\bar{x} = \frac{\sum x_i\cdot n_i}{N}=\sum x_i\cdot \frac{n_i}{N}=\sum x_i\cdot p_i\]

Pro naše data pomocí četností
\[ \bar{x} =\frac{\sum x_i\cdot n_i}{N} = \frac{1\cdot 2 + 2\cdot 3 + 3\cdot 5 + \dots }{200} = 7.56 \]

nebo pomocí relativních četností
\[\bar{x} =\sum x_i\cdot p_i = 1\cdot 0.010 + 2\cdot 0.025 + 3\cdot 0.050 + \dots = 7.56\]

\begin{itemize}
\tightlist
\item
  \textbf{Průměr} se chová jako \textbf{těžiště}. Ti nespokojení
  zákazníci na levé straně (známky 1, 2, 3) fungují jako závaží na
  dlouhém rameni páky. I když je jich málo, svou „vzdáleností`` od
  středu táhnou průměr doleva.
\end{itemize}

\begin{Shaded}
\begin{Highlighting}[]
\NormalTok{//| echo: false}
\NormalTok{//| label: fig{-}tiziste{-}oprava}

\NormalTok{// 1. Spočítáme průměr přímo zde, aby byla buňka soběstačná}
\NormalTok{mean\_vypocet = d3.sum(data\_base, d =\textgreater{} d.x * (d.n / 200))}

\NormalTok{// 2. Vizualizace páky}
\NormalTok{Plot.plot(\{}
\NormalTok{  height: 100,}
\NormalTok{  width: 600,}
\NormalTok{  marginLeft: 50,}
\NormalTok{  x: \{ domain: [1, 10], grid: false \},}
\NormalTok{  y: \{ axis: null, domain: [0, 0.4] \},}
\NormalTok{  marks: [}
\NormalTok{    // Vodorovná osa páky (deska)}
\NormalTok{    Plot.ruleY([0.1], \{x1: 1, x2: 10, strokeWidth: 2, stroke: "\#666"\}),}
    
\NormalTok{    // Závaží (body) {-} jejich plocha odpovídá počtu lidí}
\NormalTok{    Plot.dot(data\_base, \{}
\NormalTok{      x: "x", }
\NormalTok{      y: 0.1, }
\NormalTok{      r: d =\textgreater{} d.n / 10, }
\NormalTok{      fill: d =\textgreater{} d.x \textgreater{} mean\_vypocet ? "\#4e79a7" : "\#e15759", // modrá vpravo, červená vlevo od průměru}
\NormalTok{      stroke: "white"}
\NormalTok{    \}),}
    
\NormalTok{    // Těžiště (Podstavec / Fulcrum)}
\NormalTok{    Plot.areaY([}
\NormalTok{        \{x: mean\_vypocet {-} 0.2, y: 0\}, }
\NormalTok{        \{x: mean\_vypocet, y: 0.1\}, }
\NormalTok{        \{x: mean\_vypocet + 0.2, y: 0\}}
\NormalTok{      ], \{}
\NormalTok{      x: "x", }
\NormalTok{      y: "y", }
\NormalTok{      fill: "\#666"}
\NormalTok{    \}),}
    
\NormalTok{    // Text s hodnotou průměru}
\NormalTok{    Plot.text([mean\_vypocet], \{}
\NormalTok{      x: mean\_vypocet, }
\NormalTok{      y: 0.1, }
\NormalTok{      text: d =\textgreater{} \textasciigrave{}x̄ = $\{d.toFixed(2)\}\textasciigrave{}, }
\NormalTok{      dy: 25, }
\NormalTok{      fontWeight: "bold",}
\NormalTok{      fill: "\#666"}
\NormalTok{    \})}
\NormalTok{  ]}
\NormalTok{\})}
\end{Highlighting}
\end{Shaded}

\section{variabilita odvozená z
průměru}\label{variabilita-odvozenuxe1-z-prux16fmux11bru}

připomeňme, že rozptyl z dat získáme pomocí vzorce

\[\sigma^2 = \frac{\sum (x_i - \bar{x})^2}{n}\]

pokud budeme mt data vážená vektorem relativních četností
\(\mathbf{p}\), pak

\[ 
  \sigma^2 = \sum (x_i - \bar{x})^2\cdot p_i
\]

\begin{Shaded}
\begin{Highlighting}[]
\NormalTok{//| echo: false}

\NormalTok{// 1. Konstanty a ovládací prvky}
\NormalTok{source\_n = 200}
\NormalTok{viewof extra\_x = Inputs.range([1, 10], \{step: 1, value: 1, label: "Poloha extra bodu (x)"\})}
\NormalTok{viewof extra\_n = Inputs.range([0, 20], \{step: 1, value: 5, label: "Počet přidaných osob"\})}

\NormalTok{// 2. Výpočty statistik}
\NormalTok{dynamic\_data = [}
\NormalTok{  ...data\_base,}
\NormalTok{  \{x: extra\_x, n: extra\_n\}}
\NormalTok{]}
\NormalTok{total\_n = source\_n + extra\_n}
\NormalTok{current\_mean = (d3.sum(data\_base, d =\textgreater{} d.x * d.n) + (extra\_x * extra\_n)) / total\_n}

\NormalTok{current\_median = \{}
\NormalTok{  const sorted = dynamic\_data.flatMap(d =\textgreater{} Array(Math.round(d.n)).fill(d.x)).sort((a, b) =\textgreater{} a {-} b);}
\NormalTok{  return d3.median(sorted);}
\NormalTok{\}}

\NormalTok{// 3. Vizualizace}
\NormalTok{Plot.plot(\{}
\NormalTok{  height: 100,}
\NormalTok{  width: 600,}
\NormalTok{  marginLeft: 50,}
\NormalTok{  x: \{ domain: [1, 10], grid: false \},}
\NormalTok{  y: \{ axis: null, domain: [0, 0.4] \},}
\NormalTok{  marks: [}
\NormalTok{    // Vodorovná osa páky}
\NormalTok{    Plot.ruleY([0.1], \{x1: 1, x2: 10, strokeWidth: 4, stroke: "\#666"\}),}
    
\NormalTok{    // VRSTVA 1: Původní data (VELIKOST JEN Z data\_base)}
\NormalTok{    Plot.dot(data\_base, \{}
\NormalTok{      x: "x", }
\NormalTok{      y: 0.1, }
\NormalTok{      r: d =\textgreater{} d.n / 5, }
\NormalTok{      fill: d =\textgreater{} d.x \textgreater{} current\_mean ? "\#4e79a7" : "\#e15759",}
\NormalTok{      fillOpacity: 0.8}
\NormalTok{    \}),}
    
\NormalTok{    // VRSTVA 2: Extra bod (POSUNUTÝ NAHORU A DYNAMICKÝ)}
\NormalTok{    // y: 0.2 ho zvedne nad úroveň původních teček}
\NormalTok{    Plot.dot([\{x: extra\_x, n: extra\_n\}], \{}
\NormalTok{      x: "x", }
\NormalTok{      y: 0.2, }
\NormalTok{      r: extra\_n/2 , }
\NormalTok{      fill: "\#666",}
\NormalTok{      fillOpacity: 1,}
\NormalTok{      stroke: "white",}
\NormalTok{      strokeWidth: 2}
\NormalTok{    \}),}
    
\NormalTok{    // Svislá čára pro MEDIÁN}
\NormalTok{    Plot.ruleX([current\_median], \{stroke: "\#666", strokeWidth: 2, strokeDasharray: "4,4"\}),}
\NormalTok{    Plot.text([current\_median], \{x: current\_median, y: {-}0.2, text: d =\textgreater{} \textasciigrave{}Medián: $\{d\}\textasciigrave{}, dy: {-}10, fill: "\#666", fontWeight: "bold"\}),}

\NormalTok{    // Těžiště (Průměr)}
\NormalTok{    Plot.areaY([}
\NormalTok{        \{x: current\_mean {-} 0.2, y: 0\}, }
\NormalTok{        \{x: current\_mean, y: 0.1\}, }
\NormalTok{        \{x: current\_mean + 0.2, y: 0\}}
\NormalTok{      ], \{ x: "x", y: "y", fill: "\#666" \}),}
    
\NormalTok{    Plot.text([current\_mean], \{}
\NormalTok{      x: current\_mean, y: 0.08, }
\NormalTok{      text: d =\textgreater{} \textasciigrave{}x̄ = $\{d.toFixed(2)\}\textasciigrave{}, }
\NormalTok{      dy: 20, fill: "\#666", fontWeight: "bold"}
\NormalTok{    \})}
\NormalTok{  ]}
\NormalTok{\})}
\end{Highlighting}
\end{Shaded}

\chapter{```qmd}\label{qmd}

title: ``Diskrétní distribuce'' lang: cs format: html: css:
hag\_custom.scss toc: true toc-depth: 3 number-sections: true
number-depth: 3 ---

\chapter{Diskrétní distribuce}\label{diskruxe9tnuxed-distribuce}

V reálném světě se často setkáváme s událostmi, jejichž výsledky jsou
počitatelné a nabývají pouze určitých, izolovaných hodnot. Ať už jde o
počet hodů mincí, počet zákazníků, kteří přijdou do obchodu za hodinu,
nebo počet vadných výrobků v dávce, tyto situace popisujeme pomocí
\textbf{diskrétních náhodných veličin}. Tato kapitola se zaměří na
klíčové diskrétní rozdělení pravděpodobnosti, která nám umožňují
modelovat a analyzovat takové jevy.

Porozumění diskrétním distribucím je základním stavebním kamenem
statistiky. Umožňuje nám nejen předvídat budoucí události, ale také
kvantifikovat nejistotu a dělat informovaná rozhodnutí v nejrůznějších
oblastech -- od řízení kvality a financí po medicínu a sociální vědy.

V této kapitole se podrobně seznámíme s nejpoužívanějšími diskrétními
distribucemi: \textbf{Binomickým}, \textbf{Poissonovým},
\textbf{Geometrickým} a \textbf{Hypergeometrickým} rozdělením. Pro
každou z nich si vysvětlíme její princip, podmínky použití a klíčové
charakteristiky, jako je pravděpodobnostní funkce, kumulativní funkce,
střední hodnota a rozptyl. Důraz bude kladen na intuitivní pochopení a
praktické aplikace, které si demonstrujeme v nástrojích jako Excel, R,
Python a Observable JS.

\section{Základy diskrétních
distribucí}\label{zuxe1klady-diskruxe9tnuxedch-distribucuxed}

Předtím, než se ponoříme do konkrétních typů distribucí, je důležité si
ujasnit základní pojmy, které jsou společné pro všechny diskrétní
náhodné veličiny.

\subsection{Diskrétní náhodná
veličina}\label{diskruxe9tnuxed-nuxe1hodnuxe1-veliux10dina}

{\textbf{Diskrétní náhodná veličina (DNČ)}} je taková náhodná veličina,
jejíž hodnoty jsou počitatelné a obvykle celočíselné. Lze je vyjmenovat
(např. 0, 1, 2, \ldots). Nemůže nabývat žádné hodnoty mezi dvěma
sousedními hodnotami (např. nemůže být 1.5 hlavy při hodu mincí).

\textbf{Příklady diskrétních náhodných veličin:} * Počet hlav při 10
hodech mincí. * Počet aut, která projedou určitým úsekem silnice za
hodinu. * Počet vadných výrobků ve vzorku 50 kusů. * Počet hodů kostkou
potřebných k dosažení první šestky.

\subsection{Pravděpodobnostní funkce
(PMF)}\label{pravdux11bpodobnostnuxed-funkce-pmf}

{\textbf{Pravděpodobnostní funkce (Probability Mass Function - PMF)}},
značena jako \(P(X=k)\) nebo \(f(k)\), přiřazuje každé možné hodnotě
\(k\) diskrétní náhodné veličiny \(X\) pravděpodobnost, s jakou tato
hodnota nastane.

\textbf{Klíčové vlastnosti PMF:} 1. Pravděpodobnost pro jakoukoli
hodnotu \(k\) musí být nezáporná: \(P(X=k) \geq 0\). 2. Součet všech
pravděpodobností pro všechny možné hodnoty musí být roven 1:
\(\sum P(X=k) = 1\).

\textbf{Intuitivní interpretace:} PMF nám říká, ``jak pravděpodobné je,
že náhodná veličina \(X\) nabude přesně hodnoty \(k\).'' Graficky se
často zobrazuje jako sloupcový graf, kde výška sloupce nad hodnotou
\(k\) odpovídá \(P(X=k)\).

\subsection{Kumulativní distribuční funkce
(CDF)}\label{kumulativnuxed-distribuux10dnuxed-funkce-cdf}

{\textbf{Kumulativní distribuční funkce (Cumulative Distribution
Function - CDF)}}, značena jako \(F(x)\), udává pravděpodobnost, že
náhodná veličina \(X\) nabude hodnoty menší nebo rovné \(x\). Pro
diskrétní veličiny je to součet pravděpodobností PMF až do bodu \(x\).

\[ F(x) = P(X \le x) = \sum_{k \le x} P(X=k) \]

\textbf{Klíčové vlastnosti CDF:} 1. Hodnoty CDF jsou mezi 0 a 1:
\(0 \le F(x) \le 1\). 2. CDF je neklesající funkce: pokud \(x_1 < x_2\),
pak \(F(x_1) \le F(x_2)\). 3. \(F(-\infty) = 0\) a \(F(\infty) = 1\).

\textbf{Intuitivní interpretace:} CDF nám říká, ``jak pravděpodobné je,
že náhodná veličina \(X\) nepřekročí hodnotu \(x\).'' Graficky je CDF
pro diskrétní veličiny schodovitá funkce, kde skoky nastávají v možných
hodnotách \(k\) a jejich výška odpovídá \(P(X=k)\).

\subsection{Střední hodnota (Mean)}\label{stux159ednuxed-hodnota-mean}

{\textbf{Střední hodnota (Expected Value)}}, značena jako \(E[X]\) nebo
\(\mu\), je míra centrální tendence diskrétní náhodné veličiny. Je to
vážený průměr všech možných hodnot, kde váhami jsou jejich
pravděpodobnosti.

\[ E[X] = \sum_{k} k \cdot P(X=k) \]

\textbf{Intuitivní interpretace:} Střední hodnota představuje
``dlouhodobý průměr'' nebo ``očekávanou hodnotu'' náhodné veličiny,
pokud bychom experiment opakovali nekonečněkrát. Není nutné, aby střední
hodnota byla jednou z možných hodnot, které náhodná veličina může nabýt
(např. průměrný počet dětí v rodině může být 1.8, přestože nemůžete mít
1.8 dítěte).

\subsection{Rozptyl (Variance)}\label{rozptyl-variance}

{\textbf{Rozptyl (Variance)}}, značena jako \(Var[X]\) nebo
\(\sigma^2\), je míra rozptýlenosti hodnot náhodné veličiny kolem její
střední hodnoty. Udává, jak moc se jednotlivé hodnoty liší od průměru.

\[ Var[X] = E[(X - E[X])^2] = \sum_{k} (k - E[X])^2 \cdot P(X=k) \]
Často se používá i ekvivalentní vzorec: \[ Var[X] = E[X^2] - (E[X])^2 \]
kde \(E[X^2] = \sum_{k} k^2 \cdot P(X=k)\).

\textbf{Intuitivní interpretace:} Vysoký rozptyl znamená, že hodnoty
jsou široce rozprostřeny kolem střední hodnoty, zatímco nízký rozptyl
naznačuje, že hodnoty jsou soustředěny blízko střední hodnoty.
\textbf{Směrodatná odchylka} (\(\sigma = \sqrt{Var[X]}\)) je pak míra
rozptýlení ve stejných jednotkách jako náhodná veličina.

\section{Binomické rozdělení}\label{binomickuxe9-rozdux11blenuxed}

Binomické rozdělení je jedním z nejzákladnějších a nejčastěji
používaných diskrétních rozdělení. Modeluje počet úspěchů v sérii
nezávislých pokusů.

\subsection{Definice a podmínky}\label{definice-a-podmuxednky}

Náhodná veličina \(X\) má {\textbf{binomické rozdělení}} (značíme
\(X \sim Bin(n, p)\)), pokud splňuje následující podmínky: 1. Existuje
pevný počet \(n\) \textbf{nezávislých} pokusů (opakování experimentu).
2. Každý pokus má pouze dva možné výsledky: \textbf{úspěch} nebo
\textbf{neúspěch}. 3. Pravděpodobnost úspěchu \(p\) je
\textbf{konstantní} pro každý pokus. 4. Pravděpodobnost neúspěchu je
\(1-p\). 5. \(X\) je počet úspěchů v těchto \(n\) pokusech.

\textbf{Intuitivní interpretace:} Představte si, že opakujete určitý
experiment \(n\)-krát, například hod mincí, testování výrobku na vadu,
nebo pokus o vyřešení úkolu. Binomické rozdělení nám odpovídá na otázku:
``Kolikrát z \(n\) pokusů se mi podaří dosáhnout úspěchu, když
pravděpodobnost úspěchu je \(p\)?''

\textbf{Příklad:} * Počet šestek při 10 hodech kostkou (\(n=10\),
\(p=1/6\)). * Počet studentů, kteří úspěšně složí zkoušku ze skupiny 20
studentů, když pravděpodobnost složení je 0.7 (\(n=20\), \(p=0.7\)).

\subsection{Pravděpodobnostní funkce
(PMF)}\label{pravdux11bpodobnostnuxed-funkce-pmf-1}

Pravděpodobnost, že nastane přesně \(k\) úspěchů v \(n\) pokusech, je
dána vzorcem:

\[ P(X=k) = \binom{n}{k} p^k (1-p)^{n-k} \quad \text{pro } k=0, 1, \dots, n \]

kde \(\binom{n}{k} = \frac{n!}{k!(n-k)!}\) je binomický koeficient,
který udává počet způsobů, jak vybrat \(k\) úspěchů z \(n\) pokusů.

\subsection{Kumulativní distribuční funkce
(CDF)}\label{kumulativnuxed-distribuux10dnuxed-funkce-cdf-1}

CDF pro binomické rozdělení je součet PMF pro všechny hodnoty menší nebo
rovné \(x\):

\[ F(x) = P(X \le x) = \sum_{i=0}^{x} \binom{n}{i} p^i (1-p)^{n-i} \]

\subsection{Střední hodnota a
rozptyl}\label{stux159ednuxed-hodnota-a-rozptyl}

\begin{itemize}
\tightlist
\item
  \textbf{Střední hodnota (očekávaný počet úspěchů):}
  \[ E[X] = n \cdot p \]
\item
  \textbf{Rozptyl:} \[ Var[X] = n \cdot p \cdot (1-p) \]
\end{itemize}

\subsection{Praktické ukázky a
vizualizace}\label{praktickuxe9-ukuxe1zky-a-vizualizace}

Předpokládejme, že máme \textbf{10 nezávislých pokusů} a pravděpodobnost
úspěchu v každém pokusu je \textbf{\(p=0.3\)}.

\subsubsection*{Excel}\label{excel}

\begin{tcolorbox}[enhanced jigsaw, arc=.35mm, bottomrule=.15mm, bottomtitle=1mm, breakable, colback=white, colbacktitle=quarto-callout-note-color!10!white, colframe=quarto-callout-note-color-frame, coltitle=black, left=2mm, leftrule=.75mm, opacityback=0, opacitybacktitle=0.6, rightrule=.15mm, title=\textcolor{quarto-callout-note-color}{\faInfo}\hspace{0.5em}{Poznámka}, titlerule=0mm, toprule=.15mm, toptitle=1mm]

\# Příklad 1.1: Pravděpodobnostní a kumulativní funkce v Excelu
\{.unnumbered .unlisted\} Vypočítejte pravděpodobnost, že nastanou: a)
Přesně 3 úspěchy. b) Nejvýše 3 úspěchy. c) Určete nejmenší počet úspěchů
\(k\), pro který je kumulativní pravděpodobnost \(P(X \le k) \ge 0.5\).

\begin{center}\rule{0.5\linewidth}{0.5pt}\end{center}

\#\#\#\# Řešení

\begin{itemize}
\item
  \textbf{a) Přesně 3 úspěchy (\(P(X=3)\)):}
  \texttt{=BINOM.DIST(3;\ 10;\ 0.3;\ FALSE)} Výsledek: \texttt{0.2668}
  (přibližně 26.68\%)
\item
  \textbf{b) Nejvýše 3 úspěchy (\(P(X \le 3)\)):}
  \texttt{=BINOM.DIST(3;\ 10;\ 0.3;\ TRUE)} Výsledek: \texttt{0.6496}
  (přibližně 64.96\%)
\item
  \textbf{c) Nejmenší \(k\) pro \(P(X \le k) \ge 0.5\) (kvantil):}
  \texttt{=BINOM.INV(10;\ 0.3;\ 0.5)} Výsledek: \texttt{3} (Nejmenší
  počet úspěchů je 3. To znamená, že \(P(X \le 3) \ge 0.5\), konkrétně
  0.6496, zatímco \(P(X \le 2)\) by bylo menší než 0.5).
\end{itemize}

\end{tcolorbox}

\begin{tcolorbox}[enhanced jigsaw, arc=.35mm, bottomrule=.15mm, bottomtitle=1mm, breakable, colback=white, colbacktitle=quarto-callout-note-color!10!white, colframe=quarto-callout-note-color-frame, coltitle=black, left=2mm, leftrule=.75mm, opacityback=0, opacitybacktitle=0.6, rightrule=.15mm, title=\textcolor{quarto-callout-note-color}{\faInfo}\hspace{0.5em}{Poznámka}, titlerule=0mm, toprule=.15mm, toptitle=1mm]

\# Příklad 1.2: Střední hodnota a rozptyl v Excelu \{.unnumbered
.unlisted\} Vypočítejte střední hodnotu a rozptyl pro binomické
rozdělení s parametry \(n=10\) a \(p=0.3\).

\begin{center}\rule{0.5\linewidth}{0.5pt}\end{center}

\#\#\#\# Řešení * \textbf{Střední hodnota (\(E[X] = n \cdot p\)):}
\texttt{=10\ *\ 0.3} Výsledek: \texttt{3}

\begin{itemize}
\tightlist
\item
  \textbf{Rozptyl (\(Var[X] = n \cdot p \cdot (1-p)\)):}
  \texttt{=10\ *\ 0.3\ *\ (1\ -\ 0.3)} Výsledek: \texttt{2.1}
\end{itemize}

\end{tcolorbox}

\subsubsection*{R}\label{r}

\begin{tcolorbox}[enhanced jigsaw, arc=.35mm, bottomrule=.15mm, bottomtitle=1mm, breakable, colback=white, colbacktitle=quarto-callout-note-color!10!white, colframe=quarto-callout-note-color-frame, coltitle=black, left=2mm, leftrule=.75mm, opacityback=0, opacitybacktitle=0.6, rightrule=.15mm, title=\textcolor{quarto-callout-note-color}{\faInfo}\hspace{0.5em}{Poznámka}, titlerule=0mm, toprule=.15mm, toptitle=1mm]

\# Příklad 1.3: Pravděpodobnostní a kumulativní funkce v R \{.unnumbered
.unlisted\} Vypočítejte pravděpodobnost, že nastanou: a) Přesně 3
úspěchy. b) Nejvýše 3 úspěchy. c) Určete nejmenší počet úspěchů \(k\),
pro který je kumulativní pravděpodobnost \(P(X \le k) \ge 0.5\).

\begin{center}\rule{0.5\linewidth}{0.5pt}\end{center}

\#\#\#\# Řešení

\begin{Shaded}
\begin{Highlighting}[]
\NormalTok{n }\OtherTok{\textless{}{-}} \DecValTok{10}
\NormalTok{p }\OtherTok{\textless{}{-}} \FloatTok{0.3}

\CommentTok{\# a) P(X=3) {-} Pravděpodobnostní funkce (PMF)}
\NormalTok{pmf\_3 }\OtherTok{\textless{}{-}} \FunctionTok{dbinom}\NormalTok{(}\DecValTok{3}\NormalTok{, }\AttributeTok{size =}\NormalTok{ n, }\AttributeTok{prob =}\NormalTok{ p)}
\FunctionTok{print}\NormalTok{(}\FunctionTok{paste}\NormalTok{(}\StringTok{"P(X=3) ="}\NormalTok{, }\FunctionTok{round}\NormalTok{(pmf\_3, }\DecValTok{4}\NormalTok{)))}

\CommentTok{\# b) P(X\textless{}=3) {-} Kumulativní distribuční funkce (CDF)}
\NormalTok{cdf\_3 }\OtherTok{\textless{}{-}} \FunctionTok{pbinom}\NormalTok{(}\DecValTok{3}\NormalTok{, }\AttributeTok{size =}\NormalTok{ n, }\AttributeTok{prob =}\NormalTok{ p)}
\FunctionTok{print}\NormalTok{(}\FunctionTok{paste}\NormalTok{(}\StringTok{"P(X\textless{}=3) ="}\NormalTok{, }\FunctionTok{round}\NormalTok{(cdf\_3, }\DecValTok{4}\NormalTok{)))}

\CommentTok{\# c) Kvantil pro P(X\textless{}=k) \textgreater{}= 0.5}
\NormalTok{quantile\_50 }\OtherTok{\textless{}{-}} \FunctionTok{qbinom}\NormalTok{(}\FloatTok{0.5}\NormalTok{, }\AttributeTok{size =}\NormalTok{ n, }\AttributeTok{prob =}\NormalTok{ p)}
\FunctionTok{print}\NormalTok{(}\FunctionTok{paste}\NormalTok{(}\StringTok{"Nejmenší k pro P(X\textless{}=k) \textgreater{}= 0.5 ="}\NormalTok{, quantile\_50))}
\end{Highlighting}
\end{Shaded}

\end{tcolorbox}

\begin{tcolorbox}[enhanced jigsaw, arc=.35mm, bottomrule=.15mm, bottomtitle=1mm, breakable, colback=white, colbacktitle=quarto-callout-note-color!10!white, colframe=quarto-callout-note-color-frame, coltitle=black, left=2mm, leftrule=.75mm, opacityback=0, opacitybacktitle=0.6, rightrule=.15mm, title=\textcolor{quarto-callout-note-color}{\faInfo}\hspace{0.5em}{Poznámka}, titlerule=0mm, toprule=.15mm, toptitle=1mm]

\# Příklad 1.4: Střední hodnota a rozptyl v R \{.unnumbered .unlisted\}
Vypočítejte střední hodnotu a rozptyl pro binomické rozdělení s
parametry \(n=10\) a \(p=0.3\).

\begin{center}\rule{0.5\linewidth}{0.5pt}\end{center}

\#\#\#\# Řešení

\begin{Shaded}
\begin{Highlighting}[]
\NormalTok{n }\OtherTok{\textless{}{-}} \DecValTok{10}
\NormalTok{p }\OtherTok{\textless{}{-}} \FloatTok{0.3}

\CommentTok{\# Střední hodnota E[X] = n * p}
\NormalTok{mean\_binom }\OtherTok{\textless{}{-}}\NormalTok{ n }\SpecialCharTok{*}\NormalTok{ p}
\FunctionTok{print}\NormalTok{(}\FunctionTok{paste}\NormalTok{(}\StringTok{"Střední hodnota ="}\NormalTok{, mean\_binom))}

\CommentTok{\# Rozptyl Var[X] = n * p * (1{-}p)}
\NormalTok{var\_binom }\OtherTok{\textless{}{-}}\NormalTok{ n }\SpecialCharTok{*}\NormalTok{ p }\SpecialCharTok{*}\NormalTok{ (}\DecValTok{1} \SpecialCharTok{{-}}\NormalTok{ p)}
\FunctionTok{print}\NormalTok{(}\FunctionTok{paste}\NormalTok{(}\StringTok{"Rozptyl ="}\NormalTok{, var\_binom))}
\end{Highlighting}
\end{Shaded}

\end{tcolorbox}

\subsubsection*{Python (SciPy)}\label{python-scipy}

\begin{tcolorbox}[enhanced jigsaw, arc=.35mm, bottomrule=.15mm, bottomtitle=1mm, breakable, colback=white, colbacktitle=quarto-callout-note-color!10!white, colframe=quarto-callout-note-color-frame, coltitle=black, left=2mm, leftrule=.75mm, opacityback=0, opacitybacktitle=0.6, rightrule=.15mm, title=\textcolor{quarto-callout-note-color}{\faInfo}\hspace{0.5em}{Poznámka}, titlerule=0mm, toprule=.15mm, toptitle=1mm]

\# Příklad 1.5: Pravděpodobnostní a kumulativní funkce v Pythonu
\{.unnumbered .unlisted\} Vypočítejte pravděpodobnost, že nastanou: a)
Přesně 3 úspěchy. b) Nejvýše 3 úspěchy. c) Určete nejmenší počet úspěchů
\(k\), pro který je kumulativní pravděpodobnost \(P(X \le k) \ge 0.5\).

\begin{center}\rule{0.5\linewidth}{0.5pt}\end{center}

\#\#\#\# Řešení

\begin{Shaded}
\begin{Highlighting}[]
\ImportTok{from}\NormalTok{ scipy.stats }\ImportTok{import}\NormalTok{ binom}

\NormalTok{n }\OperatorTok{=} \DecValTok{10}
\NormalTok{p }\OperatorTok{=} \FloatTok{0.3}

\CommentTok{\# a) P(X=3) {-} Pravděpodobnostní funkce (PMF)}
\NormalTok{pmf\_3 }\OperatorTok{=}\NormalTok{ binom.pmf(}\DecValTok{3}\NormalTok{, n, p)}
\BuiltInTok{print}\NormalTok{(}\SpecialStringTok{f"P(X=3) = }\SpecialCharTok{\{}\NormalTok{pmf\_3}\SpecialCharTok{:.4f\}}\SpecialStringTok{"}\NormalTok{)}

\CommentTok{\# b) P(X\textless{}=3) {-} Kumulativní distribuční funkce (CDF)}
\NormalTok{cdf\_3 }\OperatorTok{=}\NormalTok{ binom.cdf(}\DecValTok{3}\NormalTok{, n, p)}
\BuiltInTok{print}\NormalTok{(}\SpecialStringTok{f"P(X\textless{}=3) = }\SpecialCharTok{\{}\NormalTok{cdf\_3}\SpecialCharTok{:.4f\}}\SpecialStringTok{"}\NormalTok{)}

\CommentTok{\# c) Kvantil pro P(X\textless{}=k) \textgreater{}= 0.5}
\CommentTok{\# PPF (Percent Point Function) je inverzní k CDF, vrací kvantil}
\NormalTok{quantile\_50 }\OperatorTok{=}\NormalTok{ binom.ppf(}\FloatTok{0.5}\NormalTok{, n, p)}
\BuiltInTok{print}\NormalTok{(}\SpecialStringTok{f"Nejmenší k pro P(X\textless{}=k) \textgreater{}= 0.5 = }\SpecialCharTok{\{}\BuiltInTok{int}\NormalTok{(quantile\_50)}\SpecialCharTok{\}}\SpecialStringTok{"}\NormalTok{)}
\end{Highlighting}
\end{Shaded}

\end{tcolorbox}

\begin{tcolorbox}[enhanced jigsaw, arc=.35mm, bottomrule=.15mm, bottomtitle=1mm, breakable, colback=white, colbacktitle=quarto-callout-note-color!10!white, colframe=quarto-callout-note-color-frame, coltitle=black, left=2mm, leftrule=.75mm, opacityback=0, opacitybacktitle=0.6, rightrule=.15mm, title=\textcolor{quarto-callout-note-color}{\faInfo}\hspace{0.5em}{Poznámka}, titlerule=0mm, toprule=.15mm, toptitle=1mm]

\# Příklad 1.6: Střední hodnota a rozptyl v Pythonu \{.unnumbered
.unlisted\} Vypočítejte střední hodnotu a rozptyl pro binomické
rozdělení s parametry \(n=10\) a \(p=0.3\).

\begin{center}\rule{0.5\linewidth}{0.5pt}\end{center}

\#\#\#\# Řešení

\begin{Shaded}
\begin{Highlighting}[]
\ImportTok{from}\NormalTok{ scipy.stats }\ImportTok{import}\NormalTok{ binom}

\NormalTok{n }\OperatorTok{=} \DecValTok{10}
\NormalTok{p }\OperatorTok{=} \FloatTok{0.3}

\CommentTok{\# Střední hodnota a rozptyl lze získat přímo}
\NormalTok{mean\_binom, var\_binom }\OperatorTok{=}\NormalTok{ binom.stats(n, p, moments}\OperatorTok{=}\StringTok{\textquotesingle{}mv\textquotesingle{}}\NormalTok{)}
\BuiltInTok{print}\NormalTok{(}\SpecialStringTok{f"Střední hodnota = }\SpecialCharTok{\{}\NormalTok{mean\_binom}\SpecialCharTok{:.1f\}}\SpecialStringTok{"}\NormalTok{)}
\BuiltInTok{print}\NormalTok{(}\SpecialStringTok{f"Rozptyl = }\SpecialCharTok{\{}\NormalTok{var\_binom}\SpecialCharTok{:.1f\}}\SpecialStringTok{"}\NormalTok{)}
\end{Highlighting}
\end{Shaded}

\end{tcolorbox}

\subsubsection*{Observable JS}\label{observable-js}

\begin{tcolorbox}[enhanced jigsaw, arc=.35mm, bottomrule=.15mm, bottomtitle=1mm, breakable, colback=white, colbacktitle=quarto-callout-note-color!10!white, colframe=quarto-callout-note-color-frame, coltitle=black, left=2mm, leftrule=.75mm, opacityback=0, opacitybacktitle=0.6, rightrule=.15mm, title=\textcolor{quarto-callout-note-color}{\faInfo}\hspace{0.5em}{Poznámka}, titlerule=0mm, toprule=.15mm, toptitle=1mm]

\# Příklad 1.7: Vizualizace binomického rozdělení v Observable JS
\{.unnumbered .unlisted\} Vizualizujte PMF a CDF binomického rozdělení s
parametry \(n=10\) a \(p=0.3\). Umožněte interaktivně měnit
pravděpodobnost \(p\) a pozorujte změny distribuce. Zobrazte také
kvantil pro zadanou kumulativní pravděpodobnost.

\begin{center}\rule{0.5\linewidth}{0.5pt}\end{center}

\#\#\#\# Řešení

\begin{Shaded}
\begin{Highlighting}[]
\NormalTok{// Parametry}
\NormalTok{n = 10}
\NormalTok{p\_input = view(Inputs.range([0, 1], \{ step: 0.01, label: "Pravděpodobnost úspěchu (p)", value: 0.3 \}))}

\NormalTok{// Kvantil k zobrazení}
\NormalTok{target\_prob\_input = view(Inputs.range([0, 1], \{ step: 0.01, label: "Cílová kumulativní pravděpodobnost", value: 0.5 \}))}

\NormalTok{// Vypočet PMF a CDF}
\NormalTok{data = Array.from(\{ length: n + 1 \}, (\_, k) =\textgreater{} \{}
\NormalTok{  const pmf = dbinom\_js(k, n, p\_input);}
\NormalTok{  return \{ k, pmf \};}
\NormalTok{\});}

\NormalTok{// Pomocná funkce pro dbinom (PMF) {-} implementace}
\NormalTok{function dbinom\_js(k, n, p) \{}
\NormalTok{  if (k \textless{} 0 || k \textgreater{} n || !Number.isInteger(k)) return 0;}
\NormalTok{  function combinations(n, k) \{}
\NormalTok{    if (k \textless{} 0 || k \textgreater{} n) return 0;}
\NormalTok{    if (k === 0 || k === n) return 1;}
\NormalTok{    if (k \textgreater{} n / 2) k = n {-} k;}
\NormalTok{    let res = 1;}
\NormalTok{    for (let i = 1; i \textless{}= k; ++i) \{}
\NormalTok{      res = res * (n {-} i + 1) / i;}
\NormalTok{    \}}
\NormalTok{    return res;}
\NormalTok{  \}}
\NormalTok{  return combinations(n, k) * Math.pow(p, k) * Math.pow(1 {-} p, n {-} k);}
\NormalTok{\}}

\NormalTok{// Výpočet CDF pro vizualizaci}
\NormalTok{data\_cdf = data.map((d, i, arr) =\textgreater{} \{}
\NormalTok{  const cdf = arr.slice(0, i + 1).reduce((sum, item) =\textgreater{} sum + item.pmf, 0);}
\NormalTok{  return \{ k: d.k, cdf: cdf \};}
\NormalTok{\});}

\NormalTok{// Najít kvantil (nejmenší k, kde CDF \textgreater{}= target\_prob\_input)}
\NormalTok{quantile\_k = data\_cdf.find(d =\textgreater{} d.cdf \textgreater{}= target\_prob\_input.value)?.k ?? 0;}

\NormalTok{Plot.plot(\{}
\NormalTok{  width: width,}
\NormalTok{  height: 400,}
\NormalTok{  title: "Binomické rozdělení (PMF a CDF)",}
\NormalTok{  x: \{ domain: [{-}0.5, n + 0.5], label: "Počet úspěchů (k)" \},}
\NormalTok{  y: \{ domain: [0, 1], label: "Pravděpodobnost", tickFormat: "\%" \},}
\NormalTok{  grid: true,}
\NormalTok{  color: \{}
\NormalTok{    domain: ["PMF", "CDF", "Kvantil"],}
\NormalTok{    range: ["steelblue", "orange", "grey"],}
\NormalTok{    legend: true}
\NormalTok{  \},}
\NormalTok{  marks: [}
\NormalTok{    Plot.barY(data, \{ x: "k", y: "pmf", fill: "steelblue", tip: true, title: d =\textgreater{} \textasciigrave{}P(X=$\{d.k\}) = $\{d.pmf.toFixed(4)\}\textasciigrave{} \}),}
\NormalTok{    Plot.step(data\_cdf, \{ x: "k", y: "cdf", stroke: "orange", marker: "circle", tip: true, title: d =\textgreater{} \textasciigrave{}P(X\textless{}=$\{d.k\}) = $\{d.cdf.toFixed(4)\}\textasciigrave{} \}),}
\NormalTok{    // Zobrazení kvantilu}
\NormalTok{    Plot.ruleY([target\_prob\_input.value], \{ stroke: "grey", strokeDasharray: "4 4", label: \textasciigrave{}Cílová pravděpodobnost ($\{(target\_prob\_input.value * 100).toFixed(0)\}\%)\textasciigrave{} \}),}
\NormalTok{    Plot.ruleX([quantile\_k], \{ stroke: "grey", strokeDasharray: "4 4", label: \textasciigrave{}Kvantil ($\{quantile\_k\})\textasciigrave{} \}),}
\NormalTok{    Plot.text([}
\NormalTok{        \{ k: quantile\_k, prob: target\_prob\_input.value \}}
\NormalTok{      ],}
\NormalTok{      \{}
\NormalTok{        x: "k",}
\NormalTok{        y: "prob",}
\NormalTok{        text: \textasciigrave{}k=$\{quantile\_k\}\textasciigrave{},}
\NormalTok{        dy: {-}10,}
\NormalTok{        dx: {-}10,}
\NormalTok{        fill: "grey",}
\NormalTok{        fontSize: 12}
\NormalTok{      \}}
\NormalTok{    )}
\NormalTok{  ]}
\NormalTok{\})}
\end{Highlighting}
\end{Shaded}

\end{tcolorbox}

\section{Poissonovo rozdělení}\label{poissonovo-rozdux11blenuxed}

Poissonovo rozdělení je ideální pro modelování počtu událostí, které
nastanou v pevném časovém intervalu nebo prostorovém úseku, když jsou
tyto události vzácné a nezávislé.

\subsection{Definice a podmínky}\label{definice-a-podmuxednky-1}

Náhodná veličina \(X\) má {\textbf{Poissonovo rozdělení}} (značíme
\(X \sim Po(\lambda)\)), pokud splňuje následující podmínky: 1. Události
nastávají s \textbf{konstantní průměrnou četností} v daném intervalu
(časovém nebo prostorovém). Tato průměrná četnost je parametr
\(\lambda\) (lambda). 2. Události jsou \textbf{nezávislé}. Nastání jedné
události neovlivňuje pravděpodobnost nastání jiné. 3. Pravděpodobnost,
že nastane více než jedna událost v krátkém intervalu, je zanedbatelná.
4. \(X\) je počet událostí v daném intervalu.

\textbf{Intuitivní interpretace:} Představte si, že sledujete určitý jev
v průběhu času nebo prostoru. Poissonovo rozdělení vám pomůže odpovědět
na otázku: ``Kolikrát se tento jev objeví v daném intervalu, když víme,
jak často se průměrně objevuje?'' Je často používáno pro modelování
\emph{vzácných událostí}.

\textbf{Příklad:} * Počet telefonních hovorů, které call centrum obdrží
za hodinu. * Počet autonehod na určitém úseku silnice za měsíc. * Počet
tiskových chyb v knize. * Počet radioaktivních rozpadů v určitém časovém
okně.

\subsection{Pravděpodobnostní funkce
(PMF)}\label{pravdux11bpodobnostnuxed-funkce-pmf-2}

Pravděpodobnost, že nastane přesně \(k\) událostí v daném intervalu, je
dána vzorcem:

\[ P(X=k) = \frac{\lambda^k e^{-\lambda}}{k!} \quad \text{pro } k=0, 1, 2, \dots \]

kde \(e\) je Eulerovo číslo (základ přirozeného logaritmu, přibližně
2.71828) a \(\lambda\) je průměrný počet událostí.

\subsection{Kumulativní distribuční funkce
(CDF)}\label{kumulativnuxed-distribuux10dnuxed-funkce-cdf-2}

CDF pro Poissonovo rozdělení je součet PMF pro všechny hodnoty menší
nebo rovné \(x\):

\[ F(x) = P(X \le x) = \sum_{i=0}^{x} \frac{\lambda^i e^{-\lambda}}{i!} \]

\subsection{Střední hodnota a
rozptyl}\label{stux159ednuxed-hodnota-a-rozptyl-1}

Zajímavostí Poissonova rozdělení je, že jeho střední hodnota i rozptyl
jsou rovny parametru \(\lambda\).

\begin{itemize}
\tightlist
\item
  \textbf{Střední hodnota (očekávaný počet událostí):}
  \[ E[X] = \lambda \]
\item
  \textbf{Rozptyl:} \[ Var[X] = \lambda \]
\end{itemize}

\subsection{Praktické ukázky a
vizualizace}\label{praktickuxe9-ukuxe1zky-a-vizualizace-1}

Předpokládejme, že průměrný počet zákazníků, kteří přijdou do obchodu za
hodinu, je \textbf{\(\lambda=5\)}.

\subsubsection*{Excel}\label{excel-1}

\begin{tcolorbox}[enhanced jigsaw, arc=.35mm, bottomrule=.15mm, bottomtitle=1mm, breakable, colback=white, colbacktitle=quarto-callout-note-color!10!white, colframe=quarto-callout-note-color-frame, coltitle=black, left=2mm, leftrule=.75mm, opacityback=0, opacitybacktitle=0.6, rightrule=.15mm, title=\textcolor{quarto-callout-note-color}{\faInfo}\hspace{0.5em}{Poznámka}, titlerule=0mm, toprule=.15mm, toptitle=1mm]

\# Příklad 2.1: Pravděpodobnostní a kumulativní funkce v Excelu
\{.unnumbered .unlisted\} Vypočítejte pravděpodobnost, že za hodinu
přijde do obchodu: a) Přesně 3 zákazníci. b) Nejvýše 5 zákazníků. c)
Určete nejmenší počet zákazníků \(k\), pro který je kumulativní
pravděpodobnost \(P(X \le k) \ge 0.9\).

\begin{center}\rule{0.5\linewidth}{0.5pt}\end{center}

\#\#\#\# Řešení

\begin{itemize}
\item
  \textbf{a) Přesně 3 zákazníci (\(P(X=3)\)):}
  \texttt{=POISSON.DIST(3;\ 5;\ FALSE)} Výsledek: \texttt{0.1404}
  (přibližně 14.04\%)
\item
  \textbf{b) Nejvýše 5 zákazníků (\(P(X \le 5)\)):}
  \texttt{=POISSON.DIST(5;\ 5;\ TRUE)} Výsledek: \texttt{0.6160}
  (přibližně 61.60\%)
\item
  \textbf{c) Nejmenší \(k\) pro \(P(X \le k) \ge 0.9\) (kvantil):} V
  Excelu není přímá funkce pro kvantil Poissonova rozdělení. Museli
  bychom iterovat přes hodnoty \(k\) a počítat
  \texttt{POISSON.DIST(k;\ 5;\ TRUE)}, dokud nenajdeme první \(k\), pro
  které je hodnota \(\ge 0.9\).

  Při \(k=7\): \texttt{=POISSON.DIST(7;\ 5;\ TRUE)} je \texttt{0.8666}.
  Při \(k=8\): \texttt{=POISSON.DIST(8;\ 5;\ TRUE)} je \texttt{0.9319}.
  Tedy nejmenší \(k\) je \texttt{8}.
\end{itemize}

\end{tcolorbox}

\begin{tcolorbox}[enhanced jigsaw, arc=.35mm, bottomrule=.15mm, bottomtitle=1mm, breakable, colback=white, colbacktitle=quarto-callout-note-color!10!white, colframe=quarto-callout-note-color-frame, coltitle=black, left=2mm, leftrule=.75mm, opacityback=0, opacitybacktitle=0.6, rightrule=.15mm, title=\textcolor{quarto-callout-note-color}{\faInfo}\hspace{0.5em}{Poznámka}, titlerule=0mm, toprule=.15mm, toptitle=1mm]

\# Příklad 2.2: Střední hodnota a rozptyl v Excelu \{.unnumbered
.unlisted\} Vypočítejte střední hodnotu a rozptyl pro Poissonovo
rozdělení s parametrem \(\lambda=5\).

\begin{center}\rule{0.5\linewidth}{0.5pt}\end{center}

\#\#\#\# Řešení * \textbf{Střední hodnota (\(E[X] = \lambda\)):}
\texttt{=5} Výsledek: \texttt{5}

\begin{itemize}
\tightlist
\item
  \textbf{Rozptyl (\(Var[X] = \lambda\)):} \texttt{=5} Výsledek:
  \texttt{5}
\end{itemize}

\end{tcolorbox}

\subsubsection*{R}\label{r-1}

\begin{tcolorbox}[enhanced jigsaw, arc=.35mm, bottomrule=.15mm, bottomtitle=1mm, breakable, colback=white, colbacktitle=quarto-callout-note-color!10!white, colframe=quarto-callout-note-color-frame, coltitle=black, left=2mm, leftrule=.75mm, opacityback=0, opacitybacktitle=0.6, rightrule=.15mm, title=\textcolor{quarto-callout-note-color}{\faInfo}\hspace{0.5em}{Poznámka}, titlerule=0mm, toprule=.15mm, toptitle=1mm]

\# Příklad 2.3: Pravděpodobnostní a kumulativní funkce v R \{.unnumbered
.unlisted\} Vypočítejte pravděpodobnost, že za hodinu přijde do obchodu:
a) Přesně 3 zákazníci. b) Nejvýše 5 zákazníků. c) Určete nejmenší počet
zákazníků \(k\), pro který je kumulativní pravděpodobnost
\(P(X \le k) \ge 0.9\).

\begin{center}\rule{0.5\linewidth}{0.5pt}\end{center}

\#\#\#\# Řešení

\begin{Shaded}
\begin{Highlighting}[]
\NormalTok{lambda }\OtherTok{\textless{}{-}} \DecValTok{5}

\CommentTok{\# a) P(X=3) {-} Pravděpodobnostní funkce (PMF)}
\NormalTok{pmf\_3\_poisson }\OtherTok{\textless{}{-}} \FunctionTok{dpois}\NormalTok{(}\DecValTok{3}\NormalTok{, }\AttributeTok{lambda =}\NormalTok{ lambda)}
\FunctionTok{print}\NormalTok{(}\FunctionTok{paste}\NormalTok{(}\StringTok{"P(X=3) ="}\NormalTok{, }\FunctionTok{round}\NormalTok{(pmf\_3\_poisson, }\DecValTok{4}\NormalTok{)))}

\CommentTok{\# b) P(X\textless{}=5) {-} Kumulativní distribuční funkce (CDF)}
\NormalTok{cdf\_5\_poisson }\OtherTok{\textless{}{-}} \FunctionTok{ppois}\NormalTok{(}\DecValTok{5}\NormalTok{, }\AttributeTok{lambda =}\NormalTok{ lambda)}
\FunctionTok{print}\NormalTok{(}\FunctionTok{paste}\NormalTok{(}\StringTok{"P(X\textless{}=5) ="}\NormalTok{, }\FunctionTok{round}\NormalTok{(cdf\_5\_poisson, }\DecValTok{4}\NormalTok{)))}

\CommentTok{\# c) Kvantil pro P(X\textless{}=k) \textgreater{}= 0.9}
\NormalTok{quantile\_90\_poisson }\OtherTok{\textless{}{-}} \FunctionTok{qpois}\NormalTok{(}\FloatTok{0.9}\NormalTok{, }\AttributeTok{lambda =}\NormalTok{ lambda)}
\FunctionTok{print}\NormalTok{(}\FunctionTok{paste}\NormalTok{(}\StringTok{"Nejmenší k pro P(X\textless{}=k) \textgreater{}= 0.9 ="}\NormalTok{, quantile\_90\_poisson))}
\end{Highlighting}
\end{Shaded}

\end{tcolorbox}

\begin{tcolorbox}[enhanced jigsaw, arc=.35mm, bottomrule=.15mm, bottomtitle=1mm, breakable, colback=white, colbacktitle=quarto-callout-note-color!10!white, colframe=quarto-callout-note-color-frame, coltitle=black, left=2mm, leftrule=.75mm, opacityback=0, opacitybacktitle=0.6, rightrule=.15mm, title=\textcolor{quarto-callout-note-color}{\faInfo}\hspace{0.5em}{Poznámka}, titlerule=0mm, toprule=.15mm, toptitle=1mm]

\# Příklad 2.4: Střední hodnota a rozptyl v R \{.unnumbered .unlisted\}
Vypočítejte střední hodnotu a rozptyl pro Poissonovo rozdělení s
parametrem \(\lambda=5\).

\begin{center}\rule{0.5\linewidth}{0.5pt}\end{center}

\#\#\#\# Řešení

\begin{Shaded}
\begin{Highlighting}[]
\NormalTok{lambda }\OtherTok{\textless{}{-}} \DecValTok{5}

\CommentTok{\# Střední hodnota E[X] = lambda}
\NormalTok{mean\_poisson }\OtherTok{\textless{}{-}}\NormalTok{ lambda}
\FunctionTok{print}\NormalTok{(}\FunctionTok{paste}\NormalTok{(}\StringTok{"Střední hodnota ="}\NormalTok{, mean\_poisson))}

\CommentTok{\# Rozptyl Var[X] = lambda}
\NormalTok{var\_poisson }\OtherTok{\textless{}{-}}\NormalTok{ lambda}
\FunctionTok{print}\NormalTok{(}\FunctionTok{paste}\NormalTok{(}\StringTok{"Rozptyl ="}\NormalTok{, var\_poisson))}
\end{Highlighting}
\end{Shaded}

\end{tcolorbox}

\subsubsection*{Python (SciPy)}\label{python-scipy-1}

\begin{tcolorbox}[enhanced jigsaw, arc=.35mm, bottomrule=.15mm, bottomtitle=1mm, breakable, colback=white, colbacktitle=quarto-callout-note-color!10!white, colframe=quarto-callout-note-color-frame, coltitle=black, left=2mm, leftrule=.75mm, opacityback=0, opacitybacktitle=0.6, rightrule=.15mm, title=\textcolor{quarto-callout-note-color}{\faInfo}\hspace{0.5em}{Poznámka}, titlerule=0mm, toprule=.15mm, toptitle=1mm]

\# Příklad 2.5: Pravděpodobnostní a kumulativní funkce v Pythonu
\{.unnumbered .unlisted\} Vypočítejte pravděpodobnost, že za hodinu
přijde do obchodu: a) Přesně 3 zákazníci. b) Nejvýše 5 zákazníků. c)
Určete nejmenší počet zákazníků \(k\), pro který je kumulativní
pravděpodobnost \(P(X \le k) \ge 0.9\).

\begin{center}\rule{0.5\linewidth}{0.5pt}\end{center}

\#\#\#\# Řešení

\begin{Shaded}
\begin{Highlighting}[]
\ImportTok{from}\NormalTok{ scipy.stats }\ImportTok{import}\NormalTok{ poisson}

\NormalTok{lambda\_param }\OperatorTok{=} \DecValTok{5}

\CommentTok{\# a) P(X=3) {-} Pravděpodobnostní funkce (PMF)}
\NormalTok{pmf\_3\_poisson }\OperatorTok{=}\NormalTok{ poisson.pmf(}\DecValTok{3}\NormalTok{, lambda\_param)}
\BuiltInTok{print}\NormalTok{(}\SpecialStringTok{f"P(X=3) = }\SpecialCharTok{\{}\NormalTok{pmf\_3\_poisson}\SpecialCharTok{:.4f\}}\SpecialStringTok{"}\NormalTok{)}

\CommentTok{\# b) P(X\textless{}=5) {-} Kumulativní distribuční funkce (CDF)}
\NormalTok{cdf\_5\_poisson }\OperatorTok{=}\NormalTok{ poisson.cdf(}\DecValTok{5}\NormalTok{, lambda\_param)}
\BuiltInTok{print}\NormalTok{(}\SpecialStringTok{f"P(X\textless{}=5) = }\SpecialCharTok{\{}\NormalTok{cdf\_5\_poisson}\SpecialCharTok{:.4f\}}\SpecialStringTok{"}\NormalTok{)}

\CommentTok{\# c) Kvantil pro P(X\textless{}=k) \textgreater{}= 0.9}
\NormalTok{quantile\_90\_poisson }\OperatorTok{=}\NormalTok{ poisson.ppf(}\FloatTok{0.9}\NormalTok{, lambda\_param)}
\BuiltInTok{print}\NormalTok{(}\SpecialStringTok{f"Nejmenší k pro P(X\textless{}=k) \textgreater{}= 0.9 = }\SpecialCharTok{\{}\BuiltInTok{int}\NormalTok{(quantile\_90\_poisson)}\SpecialCharTok{\}}\SpecialStringTok{"}\NormalTok{)}
\end{Highlighting}
\end{Shaded}

\end{tcolorbox}

\begin{tcolorbox}[enhanced jigsaw, arc=.35mm, bottomrule=.15mm, bottomtitle=1mm, breakable, colback=white, colbacktitle=quarto-callout-note-color!10!white, colframe=quarto-callout-note-color-frame, coltitle=black, left=2mm, leftrule=.75mm, opacityback=0, opacitybacktitle=0.6, rightrule=.15mm, title=\textcolor{quarto-callout-note-color}{\faInfo}\hspace{0.5em}{Poznámka}, titlerule=0mm, toprule=.15mm, toptitle=1mm]

\# Příklad 2.6: Střední hodnota a rozptyl v Pythonu \{.unnumbered
.unlisted\} Vypočítejte střední hodnotu a rozptyl pro Poissonovo
rozdělení s parametrem \(\lambda=5\).

\begin{center}\rule{0.5\linewidth}{0.5pt}\end{center}

\#\#\#\# Řešení

\begin{Shaded}
\begin{Highlighting}[]
\ImportTok{from}\NormalTok{ scipy.stats }\ImportTok{import}\NormalTok{ poisson}

\NormalTok{lambda\_param }\OperatorTok{=} \DecValTok{5}

\CommentTok{\# Střední hodnota a rozptyl lze získat přímo}
\NormalTok{mean\_poisson, var\_poisson }\OperatorTok{=}\NormalTok{ poisson.stats(lambda\_param, moments}\OperatorTok{=}\StringTok{\textquotesingle{}mv\textquotesingle{}}\NormalTok{)}
\BuiltInTok{print}\NormalTok{(}\SpecialStringTok{f"Střední hodnota = }\SpecialCharTok{\{}\NormalTok{mean\_poisson}\SpecialCharTok{:.1f\}}\SpecialStringTok{"}\NormalTok{)}
\BuiltInTok{print}\NormalTok{(}\SpecialStringTok{f"Rozptyl = }\SpecialCharTok{\{}\NormalTok{var\_poisson}\SpecialCharTok{:.1f\}}\SpecialStringTok{"}\NormalTok{)}
\end{Highlighting}
\end{Shaded}

\end{tcolorbox}

\subsubsection*{Observable JS}\label{observable-js-1}

\begin{tcolorbox}[enhanced jigsaw, arc=.35mm, bottomrule=.15mm, bottomtitle=1mm, breakable, colback=white, colbacktitle=quarto-callout-note-color!10!white, colframe=quarto-callout-note-color-frame, coltitle=black, left=2mm, leftrule=.75mm, opacityback=0, opacitybacktitle=0.6, rightrule=.15mm, title=\textcolor{quarto-callout-note-color}{\faInfo}\hspace{0.5em}{Poznámka}, titlerule=0mm, toprule=.15mm, toptitle=1mm]

\# Příklad 2.7: Vizualizace Poissonova rozdělení v Observable JS
\{.unnumbered .unlisted\} Vizualizujte PMF a CDF Poissonova rozdělení s
parametrem \(\lambda=5\). Umožněte interaktivně měnit parametr
\(\lambda\) a pozorujte změny distribuce. Zobrazte také kvantil pro
zadanou kumulativní pravděpodobnost.

\begin{center}\rule{0.5\linewidth}{0.5pt}\end{center}

\#\#\#\# Řešení

\begin{Shaded}
\begin{Highlighting}[]
\NormalTok{// Parametry}
\NormalTok{lambda\_input = view(Inputs.range([0.1, 10], \{ step: 0.1, label: "Průměrná četnost (λ)", value: 5 \}))}

\NormalTok{// Kvantil k zobrazení}
\NormalTok{target\_prob\_input\_poisson = view(Inputs.range([0, 1], \{ step: 0.01, label: "Cílová kumulativní pravděpodobnost", value: 0.9 \}))}

\NormalTok{// Vypočet PMF a CDF}
\NormalTok{max\_k = Math.ceil(lambda\_input.value + 3 * Math.sqrt(lambda\_input.value)); // Rozumný rozsah pro vizualizaci}
\NormalTok{if (max\_k \textless{} 10) max\_k = 10;}
\NormalTok{data\_poisson = Array.from(\{ length: max\_k + 1 \}, (\_, k) =\textgreater{} \{}
\NormalTok{  const pmf = dpois\_js(k, lambda\_input);}
\NormalTok{  return \{ k, pmf \};}
\NormalTok{\});}

\NormalTok{// Pomocná funkce pro dpois (PMF) {-} implementace}
\NormalTok{function dpois\_js(k, lambda) \{}
\NormalTok{  if (k \textless{} 0 || !Number.isInteger(k)) return 0;}
\NormalTok{  return (Math.pow(lambda, k) * Math.exp({-}lambda)) / factorial(k);}
\NormalTok{\}}

\NormalTok{function factorial(n) \{}
\NormalTok{  if (n \textless{} 0) return NaN;}
\NormalTok{  if (n === 0) return 1;}
\NormalTok{  let res = 1;}
\NormalTok{  for (let i = 1; i \textless{}= n; i++) \{}
\NormalTok{    res *= i;}
\NormalTok{  \}}
\NormalTok{  return res;}
\NormalTok{\}}

\NormalTok{// Výpočet CDF pro vizualizaci}
\NormalTok{data\_cdf\_poisson = data\_poisson.map((d, i, arr) =\textgreater{} \{}
\NormalTok{  const cdf = arr.slice(0, i + 1).reduce((sum, item) =\textgreater{} sum + item.pmf, 0);}
\NormalTok{  return \{ k: d.k, cdf: cdf \};}
\NormalTok{\});}

\NormalTok{// Najít kvantil (nejmenší k, kde CDF \textgreater{}= target\_prob\_input\_poisson)}
\NormalTok{quantile\_k\_poisson = data\_cdf\_poisson.find(d =\textgreater{} d.cdf \textgreater{}= target\_prob\_input\_poisson.value)?.k ?? 0;}

\NormalTok{Plot.plot(\{}
\NormalTok{  width: width,}
\NormalTok{  height: 400,}
\NormalTok{  title: "Poissonovo rozdělení (PMF a CDF)",}
\NormalTok{  x: \{ domain: [{-}0.5, max\_k + 0.5], label: "Počet událostí (k)" \},}
\NormalTok{  y: \{ domain: [0, 1], label: "Pravděpodobnost", tickFormat: "\%" \},}
\NormalTok{  grid: true,}
\NormalTok{  color: \{}
\NormalTok{    domain: ["PMF", "CDF", "Kvantil"],}
\NormalTok{    range: ["darkred", "lightgreen", "grey"],}
\NormalTok{    legend: true}
\NormalTok{  \},}
\NormalTok{  marks: [}
\NormalTok{    Plot.barY(data\_poisson, \{ x: "k", y: "pmf", fill: "darkred", tip: true, title: d =\textgreater{} \textasciigrave{}P(X=$\{d.k\}) = $\{d.pmf.toFixed(4)\}\textasciigrave{} \}),}
\NormalTok{    Plot.step(data\_cdf\_poisson, \{ x: "k", y: "cdf", stroke: "lightgreen", marker: "circle", tip: true, title: d =\textgreater{} \textasciigrave{}P(X\textless{}=$\{d.k\}) = $\{d.cdf.toFixed(4)\}\textasciigrave{} \}),}
\NormalTok{    // Zobrazení kvantilu}
\NormalTok{    Plot.ruleY([target\_prob\_input\_poisson.value], \{ stroke: "grey", strokeDasharray: "4 4", label: \textasciigrave{}Cílová pravděpodobnost ($\{(target\_prob\_input\_poisson.value * 100).toFixed(0)\}\%)\textasciigrave{} \}),}
\NormalTok{    Plot.ruleX([quantile\_k\_poisson], \{ stroke: "grey", strokeDasharray: "4 4", label: \textasciigrave{}Kvantil ($\{quantile\_k\_poisson\})\textasciigrave{} \}),}
\NormalTok{    Plot.text([}
\NormalTok{        \{ k: quantile\_k\_poisson, prob: target\_prob\_input\_poisson.value \}}
\NormalTok{      ],}
\NormalTok{      \{}
\NormalTok{        x: "k",}
\NormalTok{        y: "prob",}
\NormalTok{        text: \textasciigrave{}k=$\{quantile\_k\_poisson\}\textasciigrave{},}
\NormalTok{        dy: {-}10,}
\NormalTok{        dx: {-}10,}
\NormalTok{        fill: "grey",}
\NormalTok{        fontSize: 12}
\NormalTok{      \}}
\NormalTok{    )}
\NormalTok{  ]}
\NormalTok{\})}
\end{Highlighting}
\end{Shaded}

\end{tcolorbox}

\section{Geometrické rozdělení}\label{geometrickuxe9-rozdux11blenuxed}

Geometrické rozdělení modeluje počet nezávislých pokusů potřebných k
dosažení prvního úspěchu.

\subsection{Definice a podmínky}\label{definice-a-podmuxednky-2}

Náhodná veličina \(X\) má {\textbf{geometrické rozdělení}} (značíme
\(X \sim Geom(p)\)), pokud splňuje následující podmínky: 1. Pokusy jsou
\textbf{nezávislé}. 2. Každý pokus má pouze dva možné výsledky:
\textbf{úspěch} nebo \textbf{neúspěch}. 3. Pravděpodobnost úspěchu \(p\)
je \textbf{konstantní} pro každý pokus. 4. \(X\) je počet pokusů
\textbf{až do prvního úspěchu} (včetně tohoto úspěšného pokusu). Hodnoty
\(X\) mohou být \(1, 2, 3, \dots\). (Někdy je definováno jako počet
\emph{neúspěchů} před prvním úspěchem, pak \(X\) začíná od 0. V této
kapitole budeme držet definici, kde \(X\) je počet pokusů až do prvního
úspěchu, tedy \(X \ge 1\).)

\textbf{Intuitivní interpretace:} Představte si, že opakujete určitý
pokus, dokud nedosáhnete prvního úspěchu. Geometrické rozdělení vám
pomůže odpovědět na otázku: ``Kolik pokusů budu potřebovat, než se mi
poprvé podaří dosáhnout úspěchu, když pravděpodobnost úspěchu je
\(p\)?''

\textbf{Příklad:} * Počet hodů mincí potřebných k získání první hlavy. *
Počet zkoušek, které student musí absolvovat, než poprvé složí test (při
konstantní pravděpodobnosti složení). * Počet telefonátů, které musíte
uskutečnit, než se poprvé dovoláte k zákazníkovi.

\subsection{Pravděpodobnostní funkce
(PMF)}\label{pravdux11bpodobnostnuxed-funkce-pmf-3}

Pravděpodobnost, že první úspěch nastane přesně při \(k\)-tém pokusu, je
dána vzorcem:

\[ P(X=k) = (1-p)^{k-1} p \quad \text{pro } k=1, 2, 3, \dots \]

Tento vzorec logicky vyplývá z toho, že pro \(k\)-tý pokus jako první
úspěšný muselo předcházet \(k-1\) neúspěšných pokusů a pak jeden úspěšný
pokus.

\subsection{Kumulativní distribuční funkce
(CDF)}\label{kumulativnuxed-distribuux10dnuxed-funkce-cdf-3}

CDF pro geometrické rozdělení udává pravděpodobnost, že první úspěch
nastane nejpozději při \(x\)-tém pokusu:

\[ F(x) = P(X \le x) = \sum_{i=1}^{x} (1-p)^{i-1} p = 1 - (1-p)^x \]

Tento zjednodušený tvar je užitečný, protože \(P(X \le x)\) je doplňkem
k pravděpodobnosti, že první úspěch nastane až po \(x\)-tém pokusu (tj.
všechny první \(x\) pokusů byly neúspěšné).

\subsection{Střední hodnota a
rozptyl}\label{stux159ednuxed-hodnota-a-rozptyl-2}

\begin{itemize}
\tightlist
\item
  \textbf{Střední hodnota (očekávaný počet pokusů do prvního úspěchu):}
  \[ E[X] = \frac{1}{p} \]
\item
  \textbf{Rozptyl:} \[ Var[X] = \frac{1-p}{p^2} \]
\end{itemize}

\subsection{Praktické ukázky a
vizualizace}\label{praktickuxe9-ukuxe1zky-a-vizualizace-2}

Předpokládejme, že pravděpodobnost úspěchu v každém pokusu je
\textbf{\(p=0.2\)}.

\subsubsection*{Excel}\label{excel-2}

\begin{tcolorbox}[enhanced jigsaw, arc=.35mm, bottomrule=.15mm, bottomtitle=1mm, breakable, colback=white, colbacktitle=quarto-callout-note-color!10!white, colframe=quarto-callout-note-color-frame, coltitle=black, left=2mm, leftrule=.75mm, opacityback=0, opacitybacktitle=0.6, rightrule=.15mm, title=\textcolor{quarto-callout-note-color}{\faInfo}\hspace{0.5em}{Poznámka}, titlerule=0mm, toprule=.15mm, toptitle=1mm]

\# Příklad 3.1: Pravděpodobnostní a kumulativní funkce v Excelu
\{.unnumbered .unlisted\} Vypočítejte pravděpodobnost, že první úspěch
nastane: a) Přesně při 4. pokusu. b) Nejpozději při 5. pokusu. c) Určete
nejmenší počet pokusů \(k\), pro který je kumulativní pravděpodobnost
\(P(X \le k) \ge 0.8\).

\begin{center}\rule{0.5\linewidth}{0.5pt}\end{center}

\#\#\#\# Řešení

\begin{itemize}
\item
  \textbf{a) Přesně při 4. pokusu (\(P(X=4)\)):} V Excelu není přímá
  funkce \texttt{GEOM.DIST} pro PMF geometrického rozdělení ve
  standardní verzi pro \(X \ge 1\). Můžeme použít vzorec:
  \((1-p)^{k-1}p\). \texttt{=(1-0.2)\^{}(4-1)\ *\ 0.2}
  \texttt{=(0.8)\^{}3\ *\ 0.2} Výsledek: \texttt{0.1024} (přibližně
  10.24\%)
\item
  \textbf{b) Nejpozději při 5. pokusu (\(P(X \le 5)\)):} Můžeme použít
  vzorec CDF: \(1-(1-p)^x\). \texttt{=1\ -\ (1-0.2)\^{}5}
  \texttt{=1\ -\ (0.8)\^{}5} Výsledek: \texttt{0.67232} (přibližně
  67.23\%)
\item
  \textbf{c) Nejmenší \(k\) pro \(P(X \le k) \ge 0.8\) (kvantil):} Opět,
  v Excelu není přímá funkce. Budeme iterovat nebo použijeme inverzní
  vzorec \(k = \frac{\log(1 - P)}{\log(1 - p)}\). Pro \(P=0.8, p=0.2\):
  \texttt{=LN(1-0.8)\ /\ LN(1-0.2)} \texttt{=LN(0.2)\ /\ LN(0.8)}
  Výsledek: \texttt{7.21} Protože \(k\) musí být celé číslo a
  \(P(X \le k)\) je neklesající, vezmeme nejbližší větší celé číslo.
  Nejmenší \(k\) je \texttt{8}. Ověření:
  \(P(X \le 7) = 1 - (0.8)^7 = 0.7902\)
  \(P(X \le 8) = 1 - (0.8)^8 = 0.8322\) Tedy nejmenší \(k\) je
  \texttt{8}.
\end{itemize}

\end{tcolorbox}

\begin{tcolorbox}[enhanced jigsaw, arc=.35mm, bottomrule=.15mm, bottomtitle=1mm, breakable, colback=white, colbacktitle=quarto-callout-note-color!10!white, colframe=quarto-callout-note-color-frame, coltitle=black, left=2mm, leftrule=.75mm, opacityback=0, opacitybacktitle=0.6, rightrule=.15mm, title=\textcolor{quarto-callout-note-color}{\faInfo}\hspace{0.5em}{Poznámka}, titlerule=0mm, toprule=.15mm, toptitle=1mm]

\# Příklad 3.2: Střední hodnota a rozptyl v Excelu \{.unnumbered
.unlisted\} Vypočítejte střední hodnotu a rozptyl pro geometrické
rozdělení s parametrem \(p=0.2\).

\begin{center}\rule{0.5\linewidth}{0.5pt}\end{center}

\#\#\#\# Řešení * \textbf{Střední hodnota (\(E[X] = 1/p\)):}
\texttt{=1\ /\ 0.2} Výsledek: \texttt{5}

\begin{itemize}
\tightlist
\item
  \textbf{Rozptyl (\(Var[X] = (1-p)/p^2\)):}
  \texttt{=(1-0.2)\ /\ (0.2\^{}2)} \texttt{=0.8\ /\ 0.04} Výsledek:
  \texttt{20}
\end{itemize}

\end{tcolorbox}

\subsubsection*{R}\label{r-2}

\begin{tcolorbox}[enhanced jigsaw, arc=.35mm, bottomrule=.15mm, bottomtitle=1mm, breakable, colback=white, colbacktitle=quarto-callout-note-color!10!white, colframe=quarto-callout-note-color-frame, coltitle=black, left=2mm, leftrule=.75mm, opacityback=0, opacitybacktitle=0.6, rightrule=.15mm, title=\textcolor{quarto-callout-note-color}{\faInfo}\hspace{0.5em}{Poznámka}, titlerule=0mm, toprule=.15mm, toptitle=1mm]

\# Příklad 3.3: Pravděpodobnostní a kumulativní funkce v R \{.unnumbered
.unlisted\} Vypočítejte pravděpodobnost, že první úspěch nastane: a)
Přesně při 4. pokusu. b) Nejpozději při 5. pokusu. c) Určete nejmenší
počet pokusů \(k\), pro který je kumulativní pravděpodobnost
\(P(X \le k) \ge 0.8\).

\emph{Poznámka: Funkce \texttt{dgeom} a \texttt{pgeom} v R počítají
počet \textbf{neúspěchů} před prvním úspěchem. Tedy pokud chceme
\(P(X=k)\) pro počet pokusů, musíme použít \(dgeom(k-1, prob)\). Podobně
pro kumulativní funkci.}

\begin{center}\rule{0.5\linewidth}{0.5pt}\end{center}

\#\#\#\# Řešení

\begin{Shaded}
\begin{Highlighting}[]
\NormalTok{p }\OtherTok{\textless{}{-}} \FloatTok{0.2}

\CommentTok{\# a) P(X=4) {-} Pravděpodobnostní funkce (PMF) pro počet pokusů}
\CommentTok{\# dgeom(k{-}1, prob) kde k{-}1 je počet neúspěchů}
\NormalTok{pmf\_4\_geom }\OtherTok{\textless{}{-}} \FunctionTok{dgeom}\NormalTok{(}\DecValTok{4{-}1}\NormalTok{, }\AttributeTok{prob =}\NormalTok{ p)}
\FunctionTok{print}\NormalTok{(}\FunctionTok{paste}\NormalTok{(}\StringTok{"P(X=4) ="}\NormalTok{, }\FunctionTok{round}\NormalTok{(pmf\_4\_geom, }\DecValTok{4}\NormalTok{)))}

\CommentTok{\# b) P(X\textless{}=5) {-} Kumulativní distribuční funkce (CDF) pro počet pokusů}
\CommentTok{\# pgeom(k{-}1, prob)}
\NormalTok{cdf\_5\_geom }\OtherTok{\textless{}{-}} \FunctionTok{pgeom}\NormalTok{(}\DecValTok{5{-}1}\NormalTok{, }\AttributeTok{prob =}\NormalTok{ p)}
\FunctionTok{print}\NormalTok{(}\FunctionTok{paste}\NormalTok{(}\StringTok{"P(X\textless{}=5) ="}\NormalTok{, }\FunctionTok{round}\NormalTok{(cdf\_5\_geom, }\DecValTok{4}\NormalTok{)))}

\CommentTok{\# c) Kvantil pro P(X\textless{}=k) \textgreater{}= 0.8}
\CommentTok{\# qgeom vrací počet neúspěchů, takže k = qgeom(P, prob) + 1}
\NormalTok{quantile\_80\_geom }\OtherTok{\textless{}{-}} \FunctionTok{qgeom}\NormalTok{(}\FloatTok{0.8}\NormalTok{, }\AttributeTok{prob =}\NormalTok{ p) }\SpecialCharTok{+} \DecValTok{1}
\FunctionTok{print}\NormalTok{(}\FunctionTok{paste}\NormalTok{(}\StringTok{"Nejmenší k pro P(X\textless{}=k) \textgreater{}= 0.8 ="}\NormalTok{, quantile\_80\_geom))}
\end{Highlighting}
\end{Shaded}

\end{tcolorbox}

\begin{tcolorbox}[enhanced jigsaw, arc=.35mm, bottomrule=.15mm, bottomtitle=1mm, breakable, colback=white, colbacktitle=quarto-callout-note-color!10!white, colframe=quarto-callout-note-color-frame, coltitle=black, left=2mm, leftrule=.75mm, opacityback=0, opacitybacktitle=0.6, rightrule=.15mm, title=\textcolor{quarto-callout-note-color}{\faInfo}\hspace{0.5em}{Poznámka}, titlerule=0mm, toprule=.15mm, toptitle=1mm]

\# Příklad 3.4: Střední hodnota a rozptyl v R \{.unnumbered .unlisted\}
Vypočítejte střední hodnotu a rozptyl pro geometrické rozdělení s
parametrem \(p=0.2\).

\begin{center}\rule{0.5\linewidth}{0.5pt}\end{center}

\#\#\#\# Řešení

\begin{Shaded}
\begin{Highlighting}[]
\NormalTok{p }\OtherTok{\textless{}{-}} \FloatTok{0.2}

\CommentTok{\# Střední hodnota E[X] = 1/p}
\NormalTok{mean\_geom }\OtherTok{\textless{}{-}} \DecValTok{1} \SpecialCharTok{/}\NormalTok{ p}
\FunctionTok{print}\NormalTok{(}\FunctionTok{paste}\NormalTok{(}\StringTok{"Střední hodnota ="}\NormalTok{, mean\_geom))}

\CommentTok{\# Rozptyl Var[X] = (1{-}p)/p\^{}2}
\NormalTok{var\_geom }\OtherTok{\textless{}{-}}\NormalTok{ (}\DecValTok{1} \SpecialCharTok{{-}}\NormalTok{ p) }\SpecialCharTok{/}\NormalTok{ (p}\SpecialCharTok{\^{}}\DecValTok{2}\NormalTok{)}
\FunctionTok{print}\NormalTok{(}\FunctionTok{paste}\NormalTok{(}\StringTok{"Rozptyl ="}\NormalTok{, var\_geom))}
\end{Highlighting}
\end{Shaded}

\end{tcolorbox}

\subsubsection*{Python (SciPy)}\label{python-scipy-2}

\begin{tcolorbox}[enhanced jigsaw, arc=.35mm, bottomrule=.15mm, bottomtitle=1mm, breakable, colback=white, colbacktitle=quarto-callout-note-color!10!white, colframe=quarto-callout-note-color-frame, coltitle=black, left=2mm, leftrule=.75mm, opacityback=0, opacitybacktitle=0.6, rightrule=.15mm, title=\textcolor{quarto-callout-note-color}{\faInfo}\hspace{0.5em}{Poznámka}, titlerule=0mm, toprule=.15mm, toptitle=1mm]

\# Příklad 3.5: Pravděpodobnostní a kumulativní funkce v Pythonu
\{.unnumbered .unlisted\} Vypočítejte pravděpodobnost, že první úspěch
nastane: a) Přesně při 4. pokusu. b) Nejpozději při 5. pokusu. c) Určete
nejmenší počet pokusů \(k\), pro který je kumulativní pravděpodobnost
\(P(X \le k) \ge 0.8\).

\emph{Poznámka: \texttt{scipy.stats.geom} má dvě konvence pro PMF:
\texttt{geom.pmf(k,\ p)} pro počet neúspěchů před prvním úspěchem
(\texttt{k} od 0) a \texttt{geom.pmf(k,\ p,\ loc=1)} pro počet pokusů až
do prvního úspěchu (\texttt{k} od 1). Budeme používat konvenci
\texttt{loc=1} pro srozumitelnost.}

\begin{center}\rule{0.5\linewidth}{0.5pt}\end{center}

\#\#\#\# Řešení

\begin{Shaded}
\begin{Highlighting}[]
\ImportTok{from}\NormalTok{ scipy.stats }\ImportTok{import}\NormalTok{ geom}

\NormalTok{p }\OperatorTok{=} \FloatTok{0.2}

\CommentTok{\# a) P(X=4) {-} Pravděpodobnostní funkce (PMF)}
\CommentTok{\# loc=1 znamená, že k je počet pokusů, začínající od 1}
\NormalTok{pmf\_4\_geom }\OperatorTok{=}\NormalTok{ geom.pmf(}\DecValTok{4}\NormalTok{, p, loc}\OperatorTok{=}\DecValTok{1}\NormalTok{)}
\BuiltInTok{print}\NormalTok{(}\SpecialStringTok{f"P(X=4) = }\SpecialCharTok{\{}\NormalTok{pmf\_4\_geom}\SpecialCharTok{:.4f\}}\SpecialStringTok{"}\NormalTok{)}

\CommentTok{\# b) P(X\textless{}=5) {-} Kumulativní distribuční funkce (CDF)}
\NormalTok{cdf\_5\_geom }\OperatorTok{=}\NormalTok{ geom.cdf(}\DecValTok{5}\NormalTok{, p, loc}\OperatorTok{=}\DecValTok{1}\NormalTok{)}
\BuiltInTok{print}\NormalTok{(}\SpecialStringTok{f"P(X\textless{}=5) = }\SpecialCharTok{\{}\NormalTok{cdf\_5\_geom}\SpecialCharTok{:.4f\}}\SpecialStringTok{"}\NormalTok{)}

\CommentTok{\# c) Kvantil pro P(X\textless{}=k) \textgreater{}= 0.8}
\NormalTok{quantile\_80\_geom }\OperatorTok{=}\NormalTok{ geom.ppf(}\FloatTok{0.8}\NormalTok{, p, loc}\OperatorTok{=}\DecValTok{1}\NormalTok{)}
\BuiltInTok{print}\NormalTok{(}\SpecialStringTok{f"Nejmenší k pro P(X\textless{}=k) \textgreater{}= 0.8 = }\SpecialCharTok{\{}\BuiltInTok{int}\NormalTok{(quantile\_80\_geom)}\SpecialCharTok{\}}\SpecialStringTok{"}\NormalTok{)}
\end{Highlighting}
\end{Shaded}

\end{tcolorbox}

\begin{tcolorbox}[enhanced jigsaw, arc=.35mm, bottomrule=.15mm, bottomtitle=1mm, breakable, colback=white, colbacktitle=quarto-callout-note-color!10!white, colframe=quarto-callout-note-color-frame, coltitle=black, left=2mm, leftrule=.75mm, opacityback=0, opacitybacktitle=0.6, rightrule=.15mm, title=\textcolor{quarto-callout-note-color}{\faInfo}\hspace{0.5em}{Poznámka}, titlerule=0mm, toprule=.15mm, toptitle=1mm]

\# Příklad 3.6: Střední hodnota a rozptyl v Pythonu \{.unnumbered
.unlisted\} Vypočítejte střední hodnotu a rozptyl pro geometrické
rozdělení s parametrem \(p=0.2\).

\begin{center}\rule{0.5\linewidth}{0.5pt}\end{center}

\#\#\#\# Řešení

\begin{Shaded}
\begin{Highlighting}[]
\ImportTok{from}\NormalTok{ scipy.stats }\ImportTok{import}\NormalTok{ geom}

\NormalTok{p }\OperatorTok{=} \FloatTok{0.2}

\CommentTok{\# Střední hodnota a rozptyl lze získat přímo}
\CommentTok{\# Funkce geom.stats() vrací hodnoty pro "počet neúspěchů před prvním úspěchem"}
\CommentTok{\# Tedy pro "počet pokusů do prvního úspěchu" (loc=1) je střední hodnota (1/p) a rozptyl ((1{-}p)/p\^{}2)}
\CommentTok{\# Musíme je spočítat ručně, nebo vědět, že geom.stats(p) vrací 1/p {-} 1 a (1{-}p)/p\^{}2}
\NormalTok{mean\_geom }\OperatorTok{=} \DecValTok{1} \OperatorTok{/}\NormalTok{ p}
\NormalTok{var\_geom }\OperatorTok{=}\NormalTok{ (}\DecValTok{1} \OperatorTok{{-}}\NormalTok{ p) }\OperatorTok{/}\NormalTok{ (p}\OperatorTok{**}\DecValTok{2}\NormalTok{)}

\BuiltInTok{print}\NormalTok{(}\SpecialStringTok{f"Střední hodnota = }\SpecialCharTok{\{}\NormalTok{mean\_geom}\SpecialCharTok{:.1f\}}\SpecialStringTok{"}\NormalTok{)}
\BuiltInTok{print}\NormalTok{(}\SpecialStringTok{f"Rozptyl = }\SpecialCharTok{\{}\NormalTok{var\_geom}\SpecialCharTok{:.1f\}}\SpecialStringTok{"}\NormalTok{)}
\end{Highlighting}
\end{Shaded}

\end{tcolorbox}

\subsubsection*{Observable JS}\label{observable-js-2}

\begin{tcolorbox}[enhanced jigsaw, arc=.35mm, bottomrule=.15mm, bottomtitle=1mm, breakable, colback=white, colbacktitle=quarto-callout-note-color!10!white, colframe=quarto-callout-note-color-frame, coltitle=black, left=2mm, leftrule=.75mm, opacityback=0, opacitybacktitle=0.6, rightrule=.15mm, title=\textcolor{quarto-callout-note-color}{\faInfo}\hspace{0.5em}{Poznámka}, titlerule=0mm, toprule=.15mm, toptitle=1mm]

\# Příklad 3.7: Vizualizace geometrického rozdělení v Observable JS
\{.unnumbered .unlisted\} Vizualizujte PMF a CDF geometrického rozdělení
s parametrem \(p=0.2\). Umožněte interaktivně měnit pravděpodobnost
\(p\) a pozorujte změny distribuce. Zobrazte také kvantil pro zadanou
kumulativní pravděpodobnost.

\begin{center}\rule{0.5\linewidth}{0.5pt}\end{center}

\#\#\#\# Řešení

\begin{Shaded}
\begin{Highlighting}[]
\NormalTok{// Parametry}
\NormalTok{p\_geom\_input = view(Inputs.range([0.01, 0.99], \{ step: 0.01, label: "Pravděpodobnost úspěchu (p)", value: 0.2 \}))}

\NormalTok{// Kvantil k zobrazení}
\NormalTok{target\_prob\_input\_geom = view(Inputs.range([0, 1], \{ step: 0.01, label: "Cílová kumulativní pravděpodobnost", value: 0.8 \}))}

\NormalTok{// Vypočet PMF a CDF}
\NormalTok{max\_k\_geom = Math.ceil(1/p\_geom\_input.value + 3 * Math.sqrt((1{-}p\_geom\_input.value)/p\_geom\_input.value**2));}
\NormalTok{if (max\_k\_geom \textless{} 10) max\_k\_geom = 10;}
\NormalTok{data\_geom = Array.from(\{ length: max\_k\_geom \}, (\_, i) =\textgreater{} \{}
\NormalTok{  const k = i + 1; // k od 1}
\NormalTok{  const pmf = dgeom\_js(k, p\_geom\_input);}
\NormalTok{  return \{ k, pmf \};}
\NormalTok{\});}

\NormalTok{// Pomocná funkce pro dgeom (PMF) {-} implementace pro X \textgreater{}= 1}
\NormalTok{function dgeom\_js(k, p) \{}
\NormalTok{  if (k \textless{} 1 || !Number.isInteger(k)) return 0;}
\NormalTok{  return Math.pow(1 {-} p, k {-} 1) * p;}
\NormalTok{\}}

\NormalTok{// Výpočet CDF pro vizualizaci}
\NormalTok{data\_cdf\_geom = data\_geom.map((d, i, arr) =\textgreater{} \{}
\NormalTok{  const cdf = arr.slice(0, i + 1).reduce((sum, item) =\textgreater{} sum + item.pmf, 0);}
\NormalTok{  return \{ k: d.k, cdf: cdf \};}
\NormalTok{\});}

\NormalTok{// Najít kvantil (nejmenší k, kde CDF \textgreater{}= target\_prob\_input\_geom)}
\NormalTok{quantile\_k\_geom = data\_cdf\_geom.find(d =\textgreater{} d.cdf \textgreater{}= target\_prob\_input\_geom.value)?.k ?? 1;}

\NormalTok{Plot.plot(\{}
\NormalTok{  width: width,}
\NormalTok{  height: 400,}
\NormalTok{  title: "Geometrické rozdělení (PMF a CDF)",}
\NormalTok{  x: \{ domain: [0.5, max\_k\_geom + 0.5], label: "Počet pokusů do prvního úspěchu (k)" \},}
\NormalTok{  y: \{ domain: [0, 1], label: "Pravděpodobnost", tickFormat: "\%" \},}
\NormalTok{  grid: true,}
\NormalTok{  color: \{}
\NormalTok{    domain: ["PMF", "CDF", "Kvantil"],}
\NormalTok{    range: ["darkgreen", "purple", "grey"],}
\NormalTok{    legend: true}
\NormalTok{  \},}
\NormalTok{  marks: [}
\NormalTok{    Plot.barY(data\_geom, \{ x: "k", y: "pmf", fill: "darkgreen", tip: true, title: d =\textgreater{} \textasciigrave{}P(X=$\{d.k\}) = $\{d.pmf.toFixed(4)\}\textasciigrave{} \}),}
\NormalTok{    Plot.step(data\_cdf\_geom, \{ x: "k", y: "cdf", stroke: "purple", marker: "circle", tip: true, title: d =\textgreater{} \textasciigrave{}P(X\textless{}=$\{d.k\}) = $\{d.cdf.toFixed(4)\}\textasciigrave{} \}),}
\NormalTok{    // Zobrazení kvantilu}
\NormalTok{    Plot.ruleY([target\_prob\_input\_geom.value], \{ stroke: "grey", strokeDasharray: "4 4", label: \textasciigrave{}Cílová pravděpodobnost ($\{(target\_prob\_input\_geom.value * 100).toFixed(0)\}\%)\textasciigrave{} \}),}
\NormalTok{    Plot.ruleX([quantile\_k\_geom], \{ stroke: "grey", strokeDasharray: "4 4", label: \textasciigrave{}Kvantil ($\{quantile\_k\_geom\})\textasciigrave{} \}),}
\NormalTok{    Plot.text([}
\NormalTok{        \{ k: quantile\_k\_geom, prob: target\_prob\_input\_geom.value \}}
\NormalTok{      ],}
\NormalTok{      \{}
\NormalTok{        x: "k",}
\NormalTok{        y: "prob",}
\NormalTok{        text: \textasciigrave{}k=$\{quantile\_k\_geom\}\textasciigrave{},}
\NormalTok{        dy: {-}10,}
\NormalTok{        dx: {-}10,}
\NormalTok{        fill: "grey",}
\NormalTok{        fontSize: 12}
\NormalTok{      \}}
\NormalTok{    )}
\NormalTok{  ]}
\NormalTok{\})}
\end{Highlighting}
\end{Shaded}

\end{tcolorbox}

\section{Hypergeometrické
rozdělení}\label{hypergeometrickuxe9-rozdux11blenuxed}

Hypergeometrické rozdělení popisuje pravděpodobnost získání určitého
počtu úspěchů při výběru bez vracení z konečné populace.

\subsection{Definice a podmínky}\label{definice-a-podmuxednky-3}

Náhodná veličina \(X\) má {\textbf{hypergeometrické rozdělení}} (značíme
\(X \sim HyperGeom(N, K, n)\)), pokud splňuje následující podmínky: 1.
Existuje \textbf{konečná populace} o velikosti \(N\). 2. V populaci je
\(K\) prvků typu ``úspěch'' a \(N-K\) prvků typu ``neúspěch''. 3. Ze
populace je náhodně \textbf{vybrán vzorek o velikosti \(n\)}. 4. Výběr
probíhá \textbf{bez vracení}. 5. \(X\) je počet prvků typu ``úspěch'' ve
vybraném vzorku.

\textbf{Intuitivní interpretace:} Představte si, že máte uzavřenou
nádobu s kuličkami různých barev. Hypergeometrické rozdělení vám pomůže
odpovědět na otázku: ``Kolik kuliček určité barvy vytáhnu, když si z
nádoby vyberu \(n\) kuliček, aniž bych je vracel zpět?'' Klíčový rozdíl
oproti binomickému rozdělení je \textbf{výběr bez vracení}, což znamená,
že pravděpodobnost úspěchu se mění s každým vytaženým prvkem.

\textbf{Příklad:} * Ve třídě je 30 studentů, z toho 10 dívek a 20
chlapců. Vybereme náhodně 5 studentů. Počet dívek ve vybrané skupině.
(\(N=30\), \(K=10\), \(n=5\)) * V dodávce 100 výrobků je 5 vadných.
Vybereme vzorek 10 výrobků. Počet vadných výrobků ve vzorku. (\(N=100\),
\(K=5\), \(n=10\))

\subsection{Pravděpodobnostní funkce
(PMF)}\label{pravdux11bpodobnostnuxed-funkce-pmf-4}

Pravděpodobnost, že ve vzorku velikosti \(n\) bude přesně \(k\) úspěchů,
je dána vzorcem:

\[ P(X=k) = \frac{\binom{K}{k} \binom{N-K}{n-k}}{\binom{N}{n}} \]

kde: * \(N\) je celková velikost populace. * \(K\) je celkový počet
úspěchů v populaci. * \(n\) je velikost vybraného vzorku. * \(k\) je
počet úspěchů ve vzorku. * \(\binom{a}{b}\) je binomický koeficient,
který se vypočítá jako \(\frac{a!}{b!(a-b)!}\).

Přípustné hodnoty pro \(k\) jsou
\(\max(0, n-(N-K)) \le k \le \min(n, K)\).

\subsection{Kumulativní distribuční funkce
(CDF)}\label{kumulativnuxed-distribuux10dnuxed-funkce-cdf-4}

CDF pro hypergeometrické rozdělení je součet PMF pro všechny hodnoty
menší nebo rovné \(x\):

\[ F(x) = P(X \le x) = \sum_{i=\max(0, n-(N-K))}^{x} \frac{\binom{K}{i} \binom{N-K}{n-i}}{\binom{N}{n}} \]

\subsection{Střední hodnota a
rozptyl}\label{stux159ednuxed-hodnota-a-rozptyl-3}

\begin{itemize}
\tightlist
\item
  \textbf{Střední hodnota (očekávaný počet úspěchů ve vzorku):}
  \[ E[X] = n \frac{K}{N} \] Všimněte si podobnosti s binomickým
  rozdělením, kde \(p = K/N\).
\item
  \textbf{Rozptyl:}
  \[ Var[X] = n \frac{K}{N} \frac{N-K}{N} \frac{N-n}{N-1} \] Poslední
  člen \(\frac{N-n}{N-1}\) je tzv. \textbf{korekční faktor pro konečnou
  populaci (finite population correction factor)}. Ten zohledňuje, že
  výběr probíhá bez vracení. Pokud by \(n\) bylo malé ve srovnání s
  \(N\), tento faktor by se blížil 1 a rozptyl by se podobal binomickému
  rozdělení.
\end{itemize}

\subsection{Praktické ukázky a
vizualizace}\label{praktickuxe9-ukuxe1zky-a-vizualizace-3}

Předpokládejme, že máme populaci \(N=20\) výrobků, z nichž \(K=5\) je
vadných. Vybereme vzorek velikosti \(n=6\).

\subsubsection*{Excel}\label{excel-3}

\begin{tcolorbox}[enhanced jigsaw, arc=.35mm, bottomrule=.15mm, bottomtitle=1mm, breakable, colback=white, colbacktitle=quarto-callout-note-color!10!white, colframe=quarto-callout-note-color-frame, coltitle=black, left=2mm, leftrule=.75mm, opacityback=0, opacitybacktitle=0.6, rightrule=.15mm, title=\textcolor{quarto-callout-note-color}{\faInfo}\hspace{0.5em}{Poznámka}, titlerule=0mm, toprule=.15mm, toptitle=1mm]

\# Příklad 4.1: Pravděpodobnostní a kumulativní funkce v Excelu
\{.unnumbered .unlisted\} Vypočítejte pravděpodobnost, že ve vzorku 6
výrobků bude: a) Přesně 1 vadný výrobek. b) Nejvýše 2 vadné výrobky. c)
Určete nejmenší počet vadných výrobků \(k\), pro který je kumulativní
pravděpodobnost \(P(X \le k) \ge 0.5\).

\begin{center}\rule{0.5\linewidth}{0.5pt}\end{center}

\#\#\#\# Řešení

\begin{itemize}
\item
  \textbf{a) Přesně 1 vadný výrobek (\(P(X=1)\)):}
  \texttt{=HYPGEOM.DIST(1;\ 6;\ 5;\ 20;\ FALSE)} Výsledek:
  \texttt{0.4312} (přibližně 43.12\%)
\item
  \textbf{b) Nejvýše 2 vadné výrobky (\(P(X \le 2)\)):}
  \texttt{=HYPGEOM.DIST(2;\ 6;\ 5;\ 20;\ TRUE)} Výsledek:
  \texttt{0.9324} (přibližně 93.24\%)
\item
  \textbf{c) Nejmenší \(k\) pro \(P(X \le k) \ge 0.5\) (kvantil):} V
  Excelu není přímá funkce pro kvantil. Budeme iterovat. \(P(X \le 0)\)
  = \texttt{HYPGEOM.DIST(0;\ 6;\ 5;\ 20;\ TRUE)} = \texttt{0.2185}
  \(P(X \le 1)\) = \texttt{HYPGEOM.DIST(1;\ 6;\ 5;\ 20;\ TRUE)} =
  \texttt{0.6497} Nejmenší \(k\) pro \(P(X \le k) \ge 0.5\) je tedy
  \texttt{1}.
\end{itemize}

\end{tcolorbox}

\begin{tcolorbox}[enhanced jigsaw, arc=.35mm, bottomrule=.15mm, bottomtitle=1mm, breakable, colback=white, colbacktitle=quarto-callout-note-color!10!white, colframe=quarto-callout-note-color-frame, coltitle=black, left=2mm, leftrule=.75mm, opacityback=0, opacitybacktitle=0.6, rightrule=.15mm, title=\textcolor{quarto-callout-note-color}{\faInfo}\hspace{0.5em}{Poznámka}, titlerule=0mm, toprule=.15mm, toptitle=1mm]

\# Příklad 4.2: Střední hodnota a rozptyl v Excelu \{.unnumbered
.unlisted\} Vypočítejte střední hodnotu a rozptyl pro hypergeometrické
rozdělení s parametry \(N=20\), \(K=5\) a \(n=6\).

\begin{center}\rule{0.5\linewidth}{0.5pt}\end{center}

\#\#\#\# Řešení * \textbf{Střední hodnota (\(E[X] = n \cdot (K/N)\)):}
\texttt{=6\ *\ (5\ /\ 20)} Výsledek: \texttt{1.5}

\begin{itemize}
\tightlist
\item
  \textbf{Rozptyl
  (\(Var[X] = n \cdot (K/N) \cdot ((N-K)/N) \cdot ((N-n)/(N-1))\)):}
  \texttt{=6\ *\ (5/20)\ *\ ((20-5)/20)\ *\ ((20-6)/(20-1))}
  \texttt{=6\ *\ 0.25\ *\ (15/20)\ *\ (14/19)}
  \texttt{=6\ *\ 0.25\ *\ 0.75\ *\ 0.736842} Výsledek: \texttt{0.8289}
\end{itemize}

\end{tcolorbox}

\subsubsection*{R}\label{r-3}

\begin{tcolorbox}[enhanced jigsaw, arc=.35mm, bottomrule=.15mm, bottomtitle=1mm, breakable, colback=white, colbacktitle=quarto-callout-note-color!10!white, colframe=quarto-callout-note-color-frame, coltitle=black, left=2mm, leftrule=.75mm, opacityback=0, opacitybacktitle=0.6, rightrule=.15mm, title=\textcolor{quarto-callout-note-color}{\faInfo}\hspace{0.5em}{Poznámka}, titlerule=0mm, toprule=.15mm, toptitle=1mm]

\# Příklad 4.3: Pravděpodobnostní a kumulativní funkce v R \{.unnumbered
.unlisted\} Vypočítejte pravděpodobnost, že ve vzorku 6 výrobků bude: a)
Přesně 1 vadný výrobek. b) Nejvýše 2 vadné výrobky. c) Určete nejmenší
počet vadných výrobků \(k\), pro který je kumulativní pravděpodobnost
\(P(X \le k) \ge 0.5\).

\emph{Parametry v R: \texttt{dhyper(x,\ m,\ n,\ k)}: \texttt{x} = počet
úspěchů ve vzorku, \texttt{m} = počet úspěchů v populaci (\(K\)),
\texttt{n} = počet neúspěchů v populaci (\(N-K\)), \texttt{k} = velikost
vzorku (\(n\)).}

\begin{center}\rule{0.5\linewidth}{0.5pt}\end{center}

\#\#\#\# Řešení

\begin{Shaded}
\begin{Highlighting}[]
\NormalTok{N }\OtherTok{\textless{}{-}} \DecValTok{20} \CommentTok{\# Celková populace}
\NormalTok{K }\OtherTok{\textless{}{-}} \DecValTok{5}  \CommentTok{\# Počet úspěchů v populaci (vadných)}
\NormalTok{n }\OtherTok{\textless{}{-}} \DecValTok{6}  \CommentTok{\# Velikost vzorku}

\CommentTok{\# Parametry pro R funkce}
\NormalTok{m\_r }\OtherTok{\textless{}{-}}\NormalTok{ K }\CommentTok{\# počet bílých kuliček (úspěchů) v urně}
\NormalTok{n\_r }\OtherTok{\textless{}{-}}\NormalTok{ N }\SpecialCharTok{{-}}\NormalTok{ K }\CommentTok{\# počet černých kuliček (neúspěchů) v urně}
\NormalTok{k\_r }\OtherTok{\textless{}{-}}\NormalTok{ n }\CommentTok{\# počet tažených kuliček (velikost vzorku)}

\CommentTok{\# a) P(X=1) {-} Pravděpodobnostní funkce (PMF)}
\NormalTok{pmf\_1\_hyper }\OtherTok{\textless{}{-}} \FunctionTok{dhyper}\NormalTok{(}\DecValTok{1}\NormalTok{, }\AttributeTok{m =}\NormalTok{ m\_r, }\AttributeTok{n =}\NormalTok{ n\_r, }\AttributeTok{k =}\NormalTok{ k\_r)}
\FunctionTok{print}\NormalTok{(}\FunctionTok{paste}\NormalTok{(}\StringTok{"P(X=1) ="}\NormalTok{, }\FunctionTok{round}\NormalTok{(pmf\_1\_hyper, }\DecValTok{4}\NormalTok{)))}

\CommentTok{\# b) P(X\textless{}=2) {-} Kumulativní distribuční funkce (CDF)}
\NormalTok{cdf\_2\_hyper }\OtherTok{\textless{}{-}} \FunctionTok{phyper}\NormalTok{(}\DecValTok{2}\NormalTok{, }\AttributeTok{m =}\NormalTok{ m\_r, }\AttributeTok{n =}\NormalTok{ n\_r, }\AttributeTok{k =}\NormalTok{ k\_r)}
\FunctionTok{print}\NormalTok{(}\FunctionTok{paste}\NormalTok{(}\StringTok{"P(X\textless{}=2) ="}\NormalTok{, }\FunctionTok{round}\NormalTok{(cdf\_2\_hyper, }\DecValTok{4}\NormalTok{)))}

\CommentTok{\# c) Kvantil pro P(X\textless{}=k) \textgreater{}= 0.5}
\NormalTok{quantile\_50\_hyper }\OtherTok{\textless{}{-}} \FunctionTok{qhyper}\NormalTok{(}\FloatTok{0.5}\NormalTok{, }\AttributeTok{m =}\NormalTok{ m\_r, }\AttributeTok{n =}\NormalTok{ n\_r, }\AttributeTok{k =}\NormalTok{ k\_r)}
\FunctionTok{print}\NormalTok{(}\FunctionTok{paste}\NormalTok{(}\StringTok{"Nejmenší k pro P(X\textless{}=k) \textgreater{}= 0.5 ="}\NormalTok{, quantile\_50\_hyper))}
\end{Highlighting}
\end{Shaded}

\end{tcolorbox}

\begin{tcolorbox}[enhanced jigsaw, arc=.35mm, bottomrule=.15mm, bottomtitle=1mm, breakable, colback=white, colbacktitle=quarto-callout-note-color!10!white, colframe=quarto-callout-note-color-frame, coltitle=black, left=2mm, leftrule=.75mm, opacityback=0, opacitybacktitle=0.6, rightrule=.15mm, title=\textcolor{quarto-callout-note-color}{\faInfo}\hspace{0.5em}{Poznámka}, titlerule=0mm, toprule=.15mm, toptitle=1mm]

\# Příklad 4.4: Střední hodnota a rozptyl v R \{.unnumbered .unlisted\}
Vypočítejte střední hodnotu a rozptyl pro hypergeometrické rozdělení s
parametry \(N=20\), \(K=5\) a \(n=6\).

\begin{center}\rule{0.5\linewidth}{0.5pt}\end{center}

\#\#\#\# Řešení

\begin{Shaded}
\begin{Highlighting}[]
\NormalTok{N }\OtherTok{\textless{}{-}} \DecValTok{20}
\NormalTok{K }\OtherTok{\textless{}{-}} \DecValTok{5}
\NormalTok{n }\OtherTok{\textless{}{-}} \DecValTok{6}

\CommentTok{\# Střední hodnota E[X] = n * (K/N)}
\NormalTok{mean\_hyper }\OtherTok{\textless{}{-}}\NormalTok{ n }\SpecialCharTok{*}\NormalTok{ (K }\SpecialCharTok{/}\NormalTok{ N)}
\FunctionTok{print}\NormalTok{(}\FunctionTok{paste}\NormalTok{(}\StringTok{"Střední hodnota ="}\NormalTok{, mean\_hyper))}

\CommentTok{\# Rozptyl Var[X] = n * (K/N) * ((N{-}K)/N) * ((N{-}n)/(N{-}1))}
\NormalTok{var\_hyper }\OtherTok{\textless{}{-}}\NormalTok{ n }\SpecialCharTok{*}\NormalTok{ (K }\SpecialCharTok{/}\NormalTok{ N) }\SpecialCharTok{*}\NormalTok{ ((N }\SpecialCharTok{{-}}\NormalTok{ K) }\SpecialCharTok{/}\NormalTok{ N) }\SpecialCharTok{*}\NormalTok{ ((N }\SpecialCharTok{{-}}\NormalTok{ n) }\SpecialCharTok{/}\NormalTok{ (N }\SpecialCharTok{{-}} \DecValTok{1}\NormalTok{))}
\FunctionTok{print}\NormalTok{(}\FunctionTok{paste}\NormalTok{(}\StringTok{"Rozptyl ="}\NormalTok{, }\FunctionTok{round}\NormalTok{(var\_hyper, }\DecValTok{4}\NormalTok{)))}
\end{Highlighting}
\end{Shaded}

\end{tcolorbox}

\subsubsection*{Python (SciPy)}\label{python-scipy-3}

\begin{tcolorbox}[enhanced jigsaw, arc=.35mm, bottomrule=.15mm, bottomtitle=1mm, breakable, colback=white, colbacktitle=quarto-callout-note-color!10!white, colframe=quarto-callout-note-color-frame, coltitle=black, left=2mm, leftrule=.75mm, opacityback=0, opacitybacktitle=0.6, rightrule=.15mm, title=\textcolor{quarto-callout-note-color}{\faInfo}\hspace{0.5em}{Poznámka}, titlerule=0mm, toprule=.15mm, toptitle=1mm]

\# Příklad 4.5: Pravděpodobnostní a kumulativní funkce v Pythonu
\{.unnumbered .unlisted\} Vypočítejte pravděpodobnost, že ve vzorku 6
výrobků bude: a) Přesně 1 vadný výrobek. b) Nejvýše 2 vadné výrobky. c)
Určete nejmenší počet vadných výrobků \(k\), pro který je kumulativní
pravděpodobnost \(P(X \le k) \ge 0.5\).

\emph{Parametry v SciPy: \texttt{hypergeom.pmf(k,\ M,\ n,\ N)}:
\texttt{k} = počet úspěchů ve vzorku, \texttt{M} = celková populace
(\(N\)), \texttt{n} = počet úspěchů v populaci (\(K\)), \texttt{N} =
velikost vzorku (\(n\)).}

\begin{center}\rule{0.5\linewidth}{0.5pt}\end{center}

\#\#\#\# Řešení

\begin{Shaded}
\begin{Highlighting}[]
\ImportTok{from}\NormalTok{ scipy.stats }\ImportTok{import}\NormalTok{ hypergeom}

\NormalTok{N\_pop }\OperatorTok{=} \DecValTok{20}  \CommentTok{\# Celková velikost populace (M v SciPy)}
\NormalTok{K\_success }\OperatorTok{=} \DecValTok{5} \CommentTok{\# Počet úspěchů v populaci (n v SciPy)}
\NormalTok{n\_sample }\OperatorTok{=} \DecValTok{6} \CommentTok{\# Velikost vzorku (N v SciPy)}

\CommentTok{\# a) P(X=1) {-} Pravděpodobnostní funkce (PMF)}
\NormalTok{pmf\_1\_hyper }\OperatorTok{=}\NormalTok{ hypergeom.pmf(}\DecValTok{1}\NormalTok{, N\_pop, K\_success, n\_sample)}
\BuiltInTok{print}\NormalTok{(}\SpecialStringTok{f"P(X=1) = }\SpecialCharTok{\{}\NormalTok{pmf\_1\_hyper}\SpecialCharTok{:.4f\}}\SpecialStringTok{"}\NormalTok{)}

\CommentTok{\# b) P(X\textless{}=2) {-} Kumulativní distribuční funkce (CDF)}
\NormalTok{cdf\_2\_hyper }\OperatorTok{=}\NormalTok{ hypergeom.cdf(}\DecValTok{2}\NormalTok{, N\_pop, K\_success, n\_sample)}
\BuiltInTok{print}\NormalTok{(}\SpecialStringTok{f"P(X\textless{}=2) = }\SpecialCharTok{\{}\NormalTok{cdf\_2\_hyper}\SpecialCharTok{:.4f\}}\SpecialStringTok{"}\NormalTok{)}

\CommentTok{\# c) Kvantil pro P(X\textless{}=k) \textgreater{}= 0.5}
\NormalTok{quantile\_50\_hyper }\OperatorTok{=}\NormalTok{ hypergeom.ppf(}\FloatTok{0.5}\NormalTok{, N\_pop, K\_success, n\_sample)}
\BuiltInTok{print}\NormalTok{(}\SpecialStringTok{f"Nejmenší k pro P(X\textless{}=k) \textgreater{}= 0.5 = }\SpecialCharTok{\{}\BuiltInTok{int}\NormalTok{(quantile\_50\_hyper)}\SpecialCharTok{\}}\SpecialStringTok{"}\NormalTok{)}
\end{Highlighting}
\end{Shaded}

\end{tcolorbox}

\begin{tcolorbox}[enhanced jigsaw, arc=.35mm, bottomrule=.15mm, bottomtitle=1mm, breakable, colback=white, colbacktitle=quarto-callout-note-color!10!white, colframe=quarto-callout-note-color-frame, coltitle=black, left=2mm, leftrule=.75mm, opacityback=0, opacitybacktitle=0.6, rightrule=.15mm, title=\textcolor{quarto-callout-note-color}{\faInfo}\hspace{0.5em}{Poznámka}, titlerule=0mm, toprule=.15mm, toptitle=1mm]

\# Příklad 4.6: Střední hodnota a rozptyl v Pythonu \{.unnumbered
.unlisted\} Vypočítejte střední hodnotu a rozptyl pro hypergeometrické
rozdělení s parametry \(N=20\), \(K=5\) a \(n=6\).

\begin{center}\rule{0.5\linewidth}{0.5pt}\end{center}

\#\#\#\# Řešení

\begin{Shaded}
\begin{Highlighting}[]
\ImportTok{from}\NormalTok{ scipy.stats }\ImportTok{import}\NormalTok{ hypergeom}

\NormalTok{N\_pop }\OperatorTok{=} \DecValTok{20}
\NormalTok{K\_success }\OperatorTok{=} \DecValTok{5}
\NormalTok{n\_sample }\OperatorTok{=} \DecValTok{6}

\CommentTok{\# Střední hodnota a rozptyl lze získat přímo}
\NormalTok{mean\_hyper, var\_hyper }\OperatorTok{=}\NormalTok{ hypergeom.stats(N\_pop, K\_success, n\_sample, moments}\OperatorTok{=}\StringTok{\textquotesingle{}mv\textquotesingle{}}\NormalTok{)}
\BuiltInTok{print}\NormalTok{(}\SpecialStringTok{f"Střední hodnota = }\SpecialCharTok{\{}\NormalTok{mean\_hyper}\SpecialCharTok{:.1f\}}\SpecialStringTok{"}\NormalTok{)}
\BuiltInTok{print}\NormalTok{(}\SpecialStringTok{f"Rozptyl = }\SpecialCharTok{\{}\NormalTok{var\_hyper}\SpecialCharTok{:.4f\}}\SpecialStringTok{"}\NormalTok{)}
\end{Highlighting}
\end{Shaded}

\end{tcolorbox}

\subsubsection*{Observable JS}\label{observable-js-3}

\begin{tcolorbox}[enhanced jigsaw, arc=.35mm, bottomrule=.15mm, bottomtitle=1mm, breakable, colback=white, colbacktitle=quarto-callout-note-color!10!white, colframe=quarto-callout-note-color-frame, coltitle=black, left=2mm, leftrule=.75mm, opacityback=0, opacitybacktitle=0.6, rightrule=.15mm, title=\textcolor{quarto-callout-note-color}{\faInfo}\hspace{0.5em}{Poznámka}, titlerule=0mm, toprule=.15mm, toptitle=1mm]

\# Příklad 4.7: Vizualizace hypergeometrického rozdělení v Observable JS
\{.unnumbered .unlisted\} Vizualizujte PMF a CDF hypergeometrického
rozdělení s parametry \(N=20\), \(K=5\) a \(n=6\). Umožněte interaktivně
měnit parametry \(N\), \(K\), \(n\) a pozorujte změny distribuce.
Zobrazte také kvantil pro zadanou kumulativní pravděpodobnost.

\begin{center}\rule{0.5\linewidth}{0.5pt}\end{center}

\#\#\#\# Řešení

\begin{Shaded}
\begin{Highlighting}[]
\NormalTok{// Parametry}
\NormalTok{N\_pop\_input = view(Inputs.range([10, 100], \{ step: 1, label: "Velikost populace (N)", value: 20 \}))}
\NormalTok{K\_success\_input = view(Inputs.range([1, N\_pop\_input.value], \{ step: 1, label: "Počet úspěchů v populaci (K)", value: 5 \}))}
\NormalTok{n\_sample\_input = view(Inputs.range([1, N\_pop\_input.value], \{ step: 1, label: "Velikost vzorku (n)", value: 6 \}))}

\NormalTok{// Kvantil k zobrazení}
\NormalTok{target\_prob\_input\_hyper = view(Inputs.range([0, 1], \{ step: 0.01, label: "Cílová kumulativní pravděpodobnost", value: 0.5 \}))}

\NormalTok{// Vypočet PMF a CDF}
\NormalTok{// Určení rozsahu pro k}
\NormalTok{min\_k\_hyper = Math.max(0, n\_sample\_input.value {-} (N\_pop\_input.value {-} K\_success\_input.value));}
\NormalTok{max\_k\_hyper = Math.min(n\_sample\_input.value, K\_success\_input.value);}

\NormalTok{// Zajištění, že max\_k\_hyper není menší než min\_k\_hyper}
\NormalTok{if (min\_k\_hyper \textgreater{} max\_k\_hyper) \{}
\NormalTok{    min\_k\_hyper = max\_k\_hyper;}
\NormalTok{\}}

\NormalTok{data\_hyper = Array.from(\{ length: max\_k\_hyper {-} min\_k\_hyper + 1 \}, (\_, i) =\textgreater{} \{}
\NormalTok{  const k = min\_k\_hyper + i;}
\NormalTok{  const pmf = dhyper\_js(k, N\_pop\_input, K\_success\_input, n\_sample\_input);}
\NormalTok{  return \{ k, pmf \};}
\NormalTok{\});}

\NormalTok{// Pomocná funkce pro dhyper (PMF) {-} implementace}
\NormalTok{function dhyper\_js(k, N, K, n) \{}
\NormalTok{  function combinations(n\_comb, k\_comb) \{}
\NormalTok{    if (k\_comb \textless{} 0 || k\_comb \textgreater{} n\_comb) return 0;}
\NormalTok{    if (k\_comb === 0 || k\_comb === n\_comb) return 1;}
\NormalTok{    if (k\_comb \textgreater{} n\_comb / 2) k\_comb = n\_comb {-} k\_comb;}
\NormalTok{    let res = 1;}
\NormalTok{    for (let i = 1; i \textless{}= k\_comb; ++i) \{}
\NormalTok{      res = res * (n\_comb {-} i + 1) / i;}
\NormalTok{    \}}
\NormalTok{    return res;}
\NormalTok{  \}}
\NormalTok{  const comb\_K\_k = combinations(K, k);}
\NormalTok{  const comb\_N\_K\_n\_k = combinations(N {-} K, n {-} k);}
\NormalTok{  const comb\_N\_n = combinations(N, n);}
  
\NormalTok{  if (comb\_N\_n === 0) return 0; // Ošetření dělení nulou}
\NormalTok{  return (comb\_K\_k * comb\_N\_K\_n\_k) / comb\_N\_n;}
\NormalTok{\}}

\NormalTok{// Výpočet CDF pro vizualizaci}
\NormalTok{data\_cdf\_hyper = data\_hyper.map((d, i, arr) =\textgreater{} \{}
\NormalTok{  const cdf = arr.slice(0, i + 1).reduce((sum, item) =\textgreater{} sum + item.pmf, 0);}
\NormalTok{  return \{ k: d.k, cdf: cdf \};}
\NormalTok{\});}

\NormalTok{// Najít kvantil (nejmenší k, kde CDF \textgreater{}= target\_prob\_input\_hyper)}
\NormalTok{// Upravit pro případ prázdného data\_cdf\_hyper, kdy min\_k\_hyper \textgreater{} max\_k\_hyper}
\NormalTok{quantile\_k\_hyper = data\_cdf\_hyper.find(d =\textgreater{} d.cdf \textgreater{}= target\_prob\_input\_hyper.value)?.k ?? min\_k\_hyper;}
\NormalTok{if (data\_cdf\_hyper.length === 0) quantile\_k\_hyper = NaN;}


\NormalTok{Plot.plot(\{}
\NormalTok{  width: width,}
\NormalTok{  height: 400,}
\NormalTok{  title: "Hypergeometrické rozdělení (PMF a CDF)",}
\NormalTok{  x: \{ domain: [min\_k\_hyper {-} 0.5, max\_k\_hyper + 0.5], label: "Počet úspěchů ve vzorku (k)" \},}
\NormalTok{  y: \{ domain: [0, 1], label: "Pravděpodobnost", tickFormat: "\%" \},}
\NormalTok{  grid: true,}
\NormalTok{  color: \{}
\NormalTok{    domain: ["PMF", "CDF", "Kvantil"],}
\NormalTok{    range: ["brown", "teal", "grey"],}
\NormalTok{    legend: true}
\NormalTok{  \},}
\NormalTok{  marks: [}
\NormalTok{    Plot.barY(data\_hyper, \{ x: "k", y: "pmf", fill: "brown", tip: true, title: d =\textgreater{} \textasciigrave{}P(X=$\{d.k\}) = $\{d.pmf.toFixed(4)\}\textasciigrave{} \}),}
\NormalTok{    Plot.step(data\_cdf\_hyper, \{ x: "k", y: "cdf", stroke: "teal", marker: "circle", tip: true, title: d =\textgreater{} \textasciigrave{}P(X\textless{}=$\{d.k\}) = $\{d.cdf.toFixed(4)\}\textasciigrave{} \}),}
\NormalTok{    // Zobrazení kvantilu}
\NormalTok{    Plot.ruleY([target\_prob\_input\_hyper.value], \{ stroke: "grey", strokeDasharray: "4 4", label: \textasciigrave{}Cílová pravděpodobnost ($\{(target\_prob\_input\_hyper.value * 100).toFixed(0)\}\%)\textasciigrave{} \}),}
\NormalTok{    Plot.ruleX([quantile\_k\_hyper], \{ stroke: "grey", strokeDasharray: "4 4", label: \textasciigrave{}Kvantil ($\{quantile\_k\_hyper\})\textasciigrave{} \}),}
\NormalTok{    Plot.text([}
\NormalTok{        \{ k: quantile\_k\_hyper, prob: target\_prob\_input\_hyper.value \}}
\NormalTok{      ],}
\NormalTok{      \{}
\NormalTok{        x: "k",}
\NormalTok{        y: "prob",}
\NormalTok{        text: \textasciigrave{}k=$\{quantile\_k\_hyper\}\textasciigrave{},}
\NormalTok{        dy: {-}10,}
\NormalTok{        dx: {-}10,}
\NormalTok{        fill: "grey",}
\NormalTok{        fontSize: 12}
\NormalTok{      \}}
\NormalTok{    )}
\NormalTok{  ]}
\NormalTok{\})}
\end{Highlighting}
\end{Shaded}

\end{tcolorbox}

\section{Klíčové poznatky}\label{kluxedux10dovuxe9-poznatky-2}

\begin{itemize}
\tightlist
\item
  \textbf{Diskrétní náhodné veličiny} nabývají spočetného počtu hodnot,
  typicky celých čísel (např. počet událostí, počet úspěchů).
\item
  \textbf{Pravděpodobnostní funkce (PMF)} \(P(X=k)\) udává
  pravděpodobnost, že náhodná veličina \(X\) nabude přesně hodnoty
  \(k\).
\item
  \textbf{Kumulativní distribuční funkce (CDF)} \(F(x)=P(X \le x)\)
  udává pravděpodobnost, že náhodná veličina \(X\) nabude hodnoty menší
  nebo rovné \(x\).
\item
  \textbf{Střední hodnota} \(E[X]\) představuje dlouhodobý průměr nebo
  očekávanou hodnotu náhodné veličiny.
\item
  \textbf{Rozptyl} \(Var[X]\) měří rozptýlenost hodnot náhodné veličiny
  kolem její střední hodnoty.
\item
  \textbf{Binomické rozdělení} modeluje počet úspěchů v pevném počtu
  \(n\) nezávislých pokusů s konstantní pravděpodobností úspěchu \(p\).
\item
  \textbf{Poissonovo rozdělení} modeluje počet událostí v pevném
  intervalu času nebo prostoru, když události nastávají s konstantní
  průměrnou četností \(\lambda\) a jsou nezávislé. Typické pro vzácné
  jevy.
\item
  \textbf{Geometrické rozdělení} modeluje počet pokusů potřebných k
  dosažení prvního úspěchu (včetně tohoto úspěšného pokusu) s konstantní
  pravděpodobností úspěchu \(p\).
\item
  \textbf{Hypergeometrické rozdělení} modeluje počet úspěchů ve vzorku
  vybraném \textbf{bez vracení} z konečné populace. Pravděpodobnost
  úspěchu se s každým taženým prvkem mění.
\end{itemize}

\begin{tcolorbox}[enhanced jigsaw, arc=.35mm, bottomrule=.15mm, bottomtitle=1mm, breakable, colback=white, colbacktitle=quarto-callout-tip-color!10!white, colframe=quarto-callout-tip-color-frame, coltitle=black, left=2mm, leftrule=.75mm, opacityback=0, opacitybacktitle=0.6, rightrule=.15mm, title=\textcolor{quarto-callout-tip-color}{\faLightbulb}\hspace{0.5em}{AI Pomocník pro tuto kapitolu}, titlerule=0mm, toprule=.15mm, toptitle=1mm]

Chtěli byste lépe porozumět této kapitole? Zkuste použít následující
prompt ve vašem oblíbeném AI nástroji (např. ChatGPT, Gemini):

\begin{Shaded}
\begin{Highlighting}[]
\NormalTok{"Jsem student/ka a právě studuji kapitolu o Diskrétních distribucích z učebnice Havi Notes.}
\NormalTok{Tato kapitola pojednává o Binomickém, Poissonově, Geometrickém a Hypergeometrickém rozdělení, jejich pravděpodobnostních a kumulativních funkcích, střední hodnotě a rozptylu.}

\NormalTok{Mohl/a bys mi prosím:}
\NormalTok{1. Vysvětlit klíčové pojmy \textquotesingle{}pravděpodobnostní funkce\textquotesingle{}, \textquotesingle{}kumulativní distribuční funkce\textquotesingle{}, \textquotesingle{}střední hodnota\textquotesingle{} a \textquotesingle{}rozptyl\textquotesingle{} pro diskrétní rozdělení jednodušším jazykem s reálnými příklady?}
\NormalTok{2. Porovnat hlavní rozdíly mezi Binomickým a Hypergeometrickým rozdělením a mezi Binomickým a Geometrickým rozdělením, a uvést, kdy které použít?}
\NormalTok{3. Navrhnout 3{-}5 praktických cvičení nebo úloh (s fiktivními daty), které mi pomohou ověřit si porozumění látce, a zaměřit se na výpočet pravděpodobností (PMF i CDF), kvantilů, střední hodnoty a rozptylu pro každé z probíraných rozdělení?}
\NormalTok{4. Identifikovat 2{-}3 nejčastější chyby nebo nedorozumění, ke kterým studenti v této oblasti docházejí, a jak se jim vyhnout?}
\NormalTok{5. Navrhnout téma pro malý výzkumný projekt, kde bych mohl/a aplikovat znalosti z této kapitoly (např. v oblasti řízení kvality, sportu nebo obchodu)?}

\NormalTok{Při odpovědích se snaž být stručný/á a jasný/á, a pokud je to možné, používej český jazyk. Děkuji!"}
\end{Highlighting}
\end{Shaded}

\textbf{Tipy pro studenty:} * Můžete upravit prompt dle vašich
specifických potřeb. * Nahraďte zástupné texty \texttt{{[}\ {]}}
konkrétními informacemi z kapitoly, kterou právě studujete. *
Experimentujte s různými AI nástroji a jejich nastaveními.

\end{tcolorbox}

\begin{verbatim}

::: {.quarto-book-part}


`<!-- quarto-file-metadata: eyJyZXNvdXJjZURpciI6Ii4ifQ== -->`{=html}

```{=html}
<!-- quarto-file-metadata: eyJyZXNvdXJjZURpciI6Ii4iLCJib29rSXRlbVR5cGUiOiJwYXJ0IiwiYm9va0l0ZW1OdW1iZXIiOm51bGwsImJvb2tJdGVtRGVwdGgiOjB9 -->
\end{verbatim}

\chapter{\texorpdfstring{\textbf{Pravděpodobnost}}{Pravděpodobnost}}\label{pravdux11bpodobnost}

:::

\protect\phantomsection\label{refs}
\begin{CSLReferences}{1}{0}
\bibitem[\citeproctext]{ref-azad_betterexplained}
Azad, K. (2024). \emph{Better Explained: Math Lessons That Click}.
\url{https://betterexplained.com/}

\bibitem[\citeproctext]{ref-grinstead2012introductory}
Grinstead, C. M., \& Snell, J. L. (2012). \emph{Introductory
Probability} (Second Edition). American Mathematical Society.
\url{https://stats.libretexts.org/Bookshelves/Probability_Theory/Introductory_Probability_(Grinstead_and_Snell)}

\bibitem[\citeproctext]{ref-R-stats}
R Core Team. (2023). \emph{R: A Language and Environment for Statistical
Computing}. R Foundation for Statistical Computing.
\url{https://www.R-project.org/}

\bibitem[\citeproctext]{ref-R-hist-bins}
R Core Team. (2024). \emph{R: A Language and Environment for Statistical
Computing} (hist). R Foundation for Statistical Computing.
\url{https://www.R-project.org/}

\end{CSLReferences}




\end{document}
